% Pais gap loop -- mathematics run sheet, part 2 (loaded by \input from pais_math_runsheet.tex).
% Run items live in parts 3, 4 and 5 (loaded below); the paths table and changelog are here; references are in pais_math_runsheet_refs.tex.
%=====================================================================================================
\section{Momentum and thrust}

\subsection{A closed static system carries no momentum}
\meta{[P] run 3 $\cdot$ gap 4 $\cdot$ from I.2 (von Laue)}
\begin{theorem}\label{thm:com}
For an isolated system with a symmetric stress--energy tensor (energy density $\varepsilon$, energy flux $\mathbf S$, momentum density
$\mathbf g=\mathbf S/c^2$), the total momentum is
\begin{equation}
\mathbf P=\frac{d}{dt}\,\frac1{c^2}\int\mathbf x\,\varepsilon\,dV .\label{eq:com}
\end{equation}
A static system has $\mathbf P=0$, and the centre of energy of any isolated system moves uniformly.
\end{theorem}
\begin{proof}
With $\partial_t\varepsilon+\nabla\!\cdot\mathbf S=0$:
$\int S_i\,dV=\int S_j\partial_jx_i\,dV=-\int x_i\,\partial_jS_j\,dV=\int x_i\,\partial_t\varepsilon\,dV$, the surface term vanishing for a localized
system. Divide by $c^2$.
\end{proof}
\noindent Consequences: the field momentum $\varepsilon_0\!\int\mathbf E\times\mathbf B\,dV$ of a static charged magnet is cancelled by ``hidden''
mechanical momentum \cite{coleman}; spinning, charging or vibrating a closed device cannot move its centre of energy. Thrust needs something
ejected (mass or radiation) or something outside to push on.

\subsection{The photon-rocket bound}
\meta{[D] run 3 $\cdot$ from III.1}
\begin{statement}If the only thing a device ejects is radiation of power $P$, then $F\le P/c=3.34\un{nN/W}$.\end{statement}
\nums{100 W $\to$ 0.33 $\mu$N. Laser ablation (45 $\mu$N/W) is $1.35\e4\times P/c$ because it throws target mass. Claims of order $\mu$N/W from
closed resonant cavities exceed $P/c$ several hundred times; the Dresden group's balances bound such thrust below photon pressure for equal
power \cite{tajmar2}. Rig: $P_{\rm rad}=1.1\e{-7}\un{W}\Rightarrow F\le3.7\e{-16}\un{N}$. Pais Eq.~2's $2.5\e{10}\un{W}$ would push with 84 N
(and vaporize the disk).}

\subsection{Thermal-noise floor of a torsion thrust balance, and what really limits it}
\meta{[E] run 3 $\cdot$ gap 4 protocol}
\begin{statement}
For $I\ddot\theta+\gamma\dot\theta+\kappa\theta=\tau(t)$ with $\gamma=I\omega_0/Q$, the fluctuation--dissipation theorem gives the one-sided torque
noise $S_\tau=4k_BT\gamma$ \cite{saulson}, so averaging for a time $\tau$ resolves
\begin{equation}
\delta F\approx\frac1L\sqrt{\frac{4k_BTI\omega_0}{Q\,\tau}} .\label{eq:balance}
\end{equation}
\end{statement}
\inputs{$I=0.5\un{kg\,m^2}$, period 30 s, $Q=100$, arm $L=0.25\un{m}$, $\tau=100\un{s}$, 300 K.}
\nums{$\delta F\approx1.7\e{-12}\un{N}$. Real balances resolve 10--100 nN \cite{soni,tajmar1}, $10^4$ times worse: systematics dominate (thermal
drift and tilt, electromagnetic and electrostatic coupling). For the rig the image pressure $\varepsilon_0E^2/2=0.71\un{Pa}$ at 0.40 MV/m pulls
0.050 N on each face; staying below 100 nN needs face-to-face symmetry better than $2\e{-6}$.}
\verdict{The rig protocol must cancel by symmetry and reversal, not by sensitivity: polarity reversal, the device rotated 180$^\circ$ in vacuum,
then discharged, spin-only and vibration-only runs. Known physics predicts a null to 8--10 orders below the balance floor.}

%=====================================================================================================
\section{Vacuum thresholds}

\subsection{Critical fields}
\meta{[D] run 0 $\cdot$ gap 3}
\begin{statement}
$E_S=m_e^2c^3/(e\hbar)=1.323\e{18}\un{V/m}$, the field that does work $m_ec^2$ on an electron across one reduced Compton wavelength
$\hbar/m_ec$; $B_Q=E_S/c=4.414\e9\un{T}$. The atomic unit of field is $E_{\rm at}=e/(4\pi\varepsilon_0a_0^2)=\alpha^3E_S$.
\end{statement}
\begin{proof}[Check]
$a_0=4\pi\varepsilon_0\hbar^2/(m_ee^2)$ gives $E_{\rm at}/E_S=e^6/[(4\pi\varepsilon_0)^3\hbar^3c^3]=\alpha^3$.
\end{proof}

\subsection{Schwinger pair production}
\meta{[D] run 0 $\cdot$ gap 3 $\cdot$ from IV.1}
\begin{statement}
For a constant electric field (no magnetic field) the pair rate per volume is \cite{schwinger}
\begin{equation}
w=\frac{(eE)^2}{4\pi^3\hbar^2c}\sum_{n\ge1}\frac1{n^2}\exp\!\Big(-\frac{n\pi E_S}{E}\Big).\label{eq:schwinger}
\end{equation}
\end{statement}
\nums{($n=1$) $w(E_S)=4.7\e{54}$, $w(0.1E_S)=2.5\e{40}$, $w(0.03E_S)=3\e7\un{m^{-3}s^{-1}}$. At the field-evaporation limit of any surface,
$5.7\e{10}\un{V/m}=4.3\e{-8}E_S$, the exponent is $7.3\e7$: zero for every purpose. Matter can never hold the field.}

\subsection{Invariants: why a craft cannot boost the vacuum, and how E-320 does}
\meta{[P]/[D] run 3 $\cdot$ gap 3 $\cdot$ from IV.2}
\begin{statement}
$\mathcal F=\tfrac12(E^2-c^2B^2)$ and $\mathcal G=c\,\mathbf E\cdot\mathbf B$ are Lorentz invariant, and the constant-field pair rate depends on the
fields only through them \cite{schwinger}. A plane wave has $\mathcal F=\mathcal G=0$ and makes no pairs at any intensity; a craft moving at any
speed carries the same invariants. A boost changes what a \emph{probe particle} sees: for an electron of Lorentz factor $\gamma$ meeting a wave
at angle $\alpha$ between their directions,
\begin{equation}
E^*=\gamma(1-\beta\cos\alpha)E,\qquad\chi=E^*/E_S .\label{eq:chi}
\end{equation}
\end{statement}
\nums{E-320 at FACET-II: 10 GeV, $\gamma=1.96\e4$; $\alpha=150^\circ$ gives $\gamma(1-\beta\cos\alpha)=3.65\e4$. $\chi=0.5$ needs a laboratory field of
$1.8\e{13}\un{V/m}$ (35 times atomic), $4.4\e{19}\un{W/cm^2}$, $a_0=4.5$ peak ($a_0^2\approx10$ rms). This matches the May 2025 run:
$a_0^2\gtrsim10$, $\chi\gtrsim0.5$, radiation reaction in the quantum regime seen, nonlinear Breit--Wheeler positrons not yet directly
observed \cite{e320,e320p}. Run-4 re-check (Nov 2026): a tracker now resolves single positrons (foil proxy, 0.120 per shot); no
nonlinear Breit--Wheeler signal reported yet \cite{e320t}.}
\verdict{Gap 3 settled for a craft by invariance (vacuum wall). The physics is reachable only by probe particles in boosted frames.}

\subsection{Gertsenshtein conversion}
\meta{[D] run 0 $\cdot$ gap 5}
\begin{statement}
Photon--graviton mixing in a transverse static field over length $L$: $P=\sin^2(\Delta L)\approx(\Delta L)^2$ with
$\Delta=\sqrt{4\pi G/(\mu_0c^4)}\,B$, so
\begin{equation}
P=\frac{4\pi G\,B^2L^2}{\mu_0c^4}=8.26\e{-38}\Big(\frac{B}{\rm T}\Big)^2\Big(\frac{L}{\rm m}\Big)^2 .\label{eq:gert}
\end{equation}
\end{statement}
\nums{45 T over 10 m: $1.7\e{-32}$ \cite{gertsenshtein}.}

\subsection{Zero-point energy below a cutoff}
\meta{[D] run 0, unit corrected in R0 source check}
\begin{statement}
Two polarizations give $\omega^2/(\pi^2c^3)$ modes per volume per unit $\omega$, each holding $\hbar\omega/2$, so
\begin{equation}
u=\int_0^{\omega_c}\frac{\hbar\omega^3}{2\pi^2c^3}\,d\omega=\frac{\hbar\omega_c^4}{8\pi^2c^3}\qquad(\omega_c\text{ in rad/s}).\label{eq:zpe}
\end{equation}
\end{statement}
\nums{$\omega_c=2\pi\times10\un{THz}$: $7.7\e{-7}\un{J/m^3}$.}

%=====================================================================================================
\section{Pairing and the 10 THz drive}

\subsection{The BCS gap ratio}
\meta{[D] $\cdot$ gap 8, 9 $\cdot$ leads to V.2}
\begin{statement}In weak-coupling BCS theory $\Delta_0/k_BT_c=\pi e^{-\gamma_E}=1.764$, so $2\Delta_0=3.528\,k_BT_c$.\end{statement}
\begin{proof}
The gap equation, with $\lambda=N(0)V$ and $E=\sqrt{\xi^2+\Delta^2}$, is
\begin{equation}
\frac1\lambda=\int_0^{\hbar\omega_D}\frac{d\xi}{E}\tanh\frac{E}{2k_BT}.\label{eq:gap}
\end{equation}
At $T=0$ the tanh is 1:
\[
\frac1\lambda=\operatorname{arsinh}\frac{\hbar\omega_D}{\Delta_0}\approx\ln\frac{2\hbar\omega_D}{\Delta_0}
\quad\Longrightarrow\quad \Delta_0=2\hbar\omega_D\,e^{-1/\lambda}.
\]
At $T_c$ ($\Delta\to0$), with $X=\hbar\omega_D/2k_BT_c$ and $\int_0^X\frac{\tanh x}{x}dx=\ln X+\ln\frac{4e^{\gamma_E}}{\pi}+o(1)$:
\[
\frac1\lambda=\ln\frac{2e^{\gamma_E}\hbar\omega_D}{\pi k_BT_c}
\quad\Longrightarrow\quad k_BT_c=\frac{2e^{\gamma_E}}{\pi}\,\hbar\omega_D\,e^{-1/\lambda}.
\]
Divide \cite{bcs,tinkham}.
\end{proof}

\subsection{The pairing photon}
\meta{[D] run 0 $\cdot$ from V.1}
\nums{$f=2\Delta_0/h=3.528k_BT_c/h$: 22.05 THz for $T_c=300\un{K}$; 1.47 THz for K$_3$C$_{60}$ ($T_c=20\un{K}$). The K$_3$C$_{60}$ drive resonance at
41 meV (9.9 THz, about 16 meV wide) is 6.8 times its own gap frequency: a mode resonance (a radial T$_{1u}$ phonon, an on-ball orbital
excitation or a dressed exciton; no accepted theory) \cite{rowe}, not a pairing energy.}

\subsection{Holding a driven state: pulse bookkeeping}
\meta{[D] runs 2--3 $\cdot$ gap 9, bridge B4 $\cdot$ canary in run 3}
\begin{statement}
With fluence $F$ needed per pulse, the spot is $d=2\sqrt{E/\pi F}$, the average power $P=Ef_{\rm rep}$ and the average intensity $Ff_{\rm rep}$.
\end{statement}
\inputs{$F=400\,\mu\mathrm{J/cm^2}=4\un{J/m^2}$ at the 10 THz resonance \cite{rowe}.}
\nums{FELBE (13 MHz, $\sim$20 W): $E=1.54\,\mu$J, $d=0.70\un{mm}$ (canary; logged 0.69), $5.2\un{kW/cm^2}$. The authors' plan (100 MHz, 100 $\mu$m):
31 nJ, 3.1 W; with the order-of-magnitude gain in coupling the authors say steady operation needs, 0.31 W. Shanghai ERL-FEL design (2--10 THz,
83.33 MHz, $\sim$1 kW \cite{shanghai}): $E=12\,\mu$J, enough for a 1.95 mm spot at full fluence, $33\un{kW/cm^2}$.}
\par\noindent\emph{Re-drive before it decays} ($f_{\rm rep}\ge1/\tau$, 100 $\mu$m spot, diamond substrate, V.4): at 100 K the state lasted
$\tau\approx10\un{ns}$, so 100 MHz, 3.1 W and a 10 K rise; at room temperature only $\tau\ge0.5\un{ns}$ is predicted \cite{rowe}, so 2 GHz,
63 W and a 200 K rise, or 6.3 W and 20 K with the tenfold coupling gain.
\verdict{Bridge B4 kept, sharpened: a facility-class source matching the 100 K plan exists on paper; room temperature needs gigahertz
repetition and the coupling gain. First test: does the state build up across a FELBE pulse train?}

\subsection{Heating of a driven spot}
\meta{[P] run 2 correction of run 1 $\cdot$ gap 9 $\cdot$ from V.3}
\begin{statement}
A point source $P$ on the insulated surface of a half-space of conductivity $k$ gives $T(r)=P/(2\pi kr)$; a uniform flux $q=P/\pi a^2$ over a
disk gives a centre rise
\begin{equation}
\Delta T=\int_0^a\frac{q}{2\pi k\rho}\,2\pi\rho\,d\rho=\frac{qa}{k}=\frac{P}{\pi ak},\label{eq:spot}
\end{equation}
and a film of thickness $t$ on a heat sink adds $qt/k_{\rm film}$.
\end{statement}
\begin{proof}Steady state: the flux through a hemisphere of radius $r$ is $2\pi r^2k\,|dT/dr|=P$. Superpose over the disk.\end{proof}
\nums{The authors' plan on diamond ($k=2000\un{W/m/K}$): 10 K. FELBE spot on diamond: 9.1 K, plus the film term.}

\subsection{No mechanical route to 10 THz (bridge B3)}
\meta{[E] run 3 $\cdot$ gap 9}
\nums{Acoustic wavelength $\lambda=v_s/f$: at 10 THz with $v_s\approx5\un{km/s}$, $\lambda\approx0.5\un{nm}$, the lattice scale, where only
zone-boundary or optical modes exist. The K$_3$C$_{60}$ 41 meV mode is intramolecular and infrared-active: it is driven by the electric field
of light \cite{rowe}. Electrically pumped coherent-phonon sources (sasers) work from sub-THz to about 1 THz \cite{saser}.}
\verdict{Bridge B3 (Pais's ``vibration'' at the pairing resonance) cut.}

%=====================================================================================================
\section*{S\quad Standalone items}
\addcontentsline{toc}{section}{S Standalone items}
\noindent\textbf{S.1 Earth's binding energy} [D, R0]. Assemble a uniform sphere shell by shell: $dU=Gm(r)\,dm/r$ with $m=Mr^3/R^3$, so
$U=3GM^2/5R=2.24\e{32}\un{J}$. At Pais's $1.04\e{10}\un{W}$: $6.8\e{14}$ years (ledger 8 wording corrected to $1.04\e{10}$ W).

\smallskip\noindent\textbf{S.2 Thermal occupation} [D, R0]. $\hbar\omega/k_BT=1.6\e{-5}$ at 100 MHz and 300 K: a 100 MHz vibration is classical
and far below any pairing energy (V.2).

\smallskip\noindent\textbf{S.3 Kinetic impactor} [D, R0]. Head-on, $(M_A+m)\Delta v=\beta mu$. Pais Eq.~3's inputs ($10^{10}\un{kg}$, $10^4\un{kg}$,
10 km/s, $\beta=1$) give 0.010 m/s, not the printed 0.001. DART measured $\beta=3.61$ and an along-track $\Delta v=2.70\un{mm/s}$ (a projection
set by the impact geometry) \cite{dart}.

%=====================================================================================================
% Pais gap loop -- mathematics run sheet, part 3: run-4 items (loaded by \input from part 2, before the paths table).
% Items keep their chain numbers (counters set below). Numbers: run4_numbers.json from run4_math.py.
%=====================================================================================================
\section*{Run 4 additions to chains I, III and V}
\addcontentsline{toc}{section}{Run 4 additions: I.11--I.14, III.4, V.6, V.7}
\noindent Each item keeps its chain number and names the item it builds on. Numbers come from \texttt{run4\_math.py}; the run opened with a
blind re-derivation of four logged run-3 numbers from their written inputs (\texttt{canary4.py}, 4 of 4 within 1\,\%).

\setcounter{section}{1}\setcounter{subsection}{10}
\subsection{The fibre ladder: specific strength at buildable length}
\meta{[D] run 4 $\cdot$ gap 2 (K1) $\cdot$ from I.4, I.7, I.8}
\begin{statement}
A hoop spinning at rim speed $v$ needs $\sigma/\rho\ge v^2$ (I.4). In textile units $1\un{N/tex}=1\un{GPa/(g\,cm^{-3})}=10^6\un{J/kg}$, so gap 2
($R=2\un{m}$, 30{,}000 rpm, $v=6.28\un{km/s}$) needs $39.5\un{N/tex}$. A unidirectional composite with fibre fraction $V_f$ in a matrix of density
$\rho_m$ needs fibre strength
\begin{equation}
\sigma_f\ge v^2\,\frac{V_f\rho_f+(1-V_f)\rho_m}{V_f}.\label{eq:dilution}
\end{equation}
\end{statement}
\begin{proof}
$1\un{tex}=1\un{g/km}=10^{-6}\un{kg/m}$, so $1\un{N/tex}=10^6\un{N\,m/kg}=10^6\un{J/kg}$. For \eqref{eq:dilution}: the composite's hoop stress
is $\rho_cv^2$ with $\rho_c=V_f\rho_f+(1-V_f)\rho_m$, and the fibres carry it, $V_f\sigma_f\ge\rho_cv^2$.
\end{proof}
\inputs{commercial CNT yarn at 100--120 mm gauge, $0.911\pm0.137$ N/tex \cite{dessureault}; continuous CNT yarn (150--500 m lengths made),
$1.77\pm0.07$ N/tex \cite{liang23}; Kevlar 49, Dyneema, Zylon and T1100G from strength and density data \cite{wikiss}; ultralong CNT bundles
80 GPa at 2 cm \cite{bai}, $\rho=1300$ (assumed); graphene 130 GPa, 2200 kg/m$^3$ \cite{lee}; $\rho_m=1200$.}
\nums{Factor short of 39.5 N/tex: CNT yarn at 0.1 m, $43\times$; continuous CNT yarn, $22\times$; Kevlar, $16\times$; Dyneema, Zylon and T1100G,
$10\times$ (T1100G 3.91 N/tex, rim 1.98 km/s, $3.2\times$ short in speed). Over the line: only the 2 cm bundles (61.5 N/tex) and ideal graphene
(59.1). As a $V_f=0.6$ composite the bundles need 82.9 GPa, $1.04\times$ their cm-scale strength; graphene would need 91\,\% of ideal.}
\verdict{At buildable length nanotube yarns are $2$--$4\times$ \emph{weaker} per kilogram than carbon fibre. Gap 2 at buildable size is $10\times$ in
$\sigma/\rho$, and only the dry cm-scale bundle is over the line. K1 kept with a lowered prior.}

\subsection{How many specimens decide the Weibull modulus}
\meta{[P] run 4 $\cdot$ gap 2 (K1 test design) $\cdot$ from I.8}
\begin{theorem}\label{thm:weibullN}
For $N$ strengths drawn from a Weibull law with unknown scale, the maximum-likelihood modulus $\hat m$ is asymptotically normal with
\begin{equation}
\operatorname{Var}\hat m=\frac{6}{\pi^2}\,\frac{m^2}{N};\qquad\text{telling $m$ from $m+\Delta m$ at $z$ standard errors needs}\quad
N\ge\Big(\frac{z\sqrt6}{\pi}\,\frac{m}{\Delta m}\Big)^2 .\label{eq:weibullN}
\end{equation}
\end{theorem}
\begin{proof}
$y=\ln\sigma$ follows the smallest-extreme-value (Gumbel) law with location $\mu=\ln\sigma_0$ and scale $b=1/m$. With $w=e^{(y-\mu)/b}\sim$
Exp(1), the Fisher information per observation is
$b^{-2}\bigl(\begin{smallmatrix}1&1-\gamma_E\\ 1-\gamma_E&\pi^2/6+(1-\gamma_E)^2\end{smallmatrix}\bigr)$ (moments of $\ln w$ and $w\ln w$ from
$\Gamma'(k)$, $\Gamma''(k)$). Its determinant is $\pi^2/(6b^4)$, so the $(b,b)$ element of the inverse is $6b^2/\pi^2$. With $m=1/b$ the delta
method gives $\operatorname{Var}\hat m=\operatorname{Var}\hat b/b^4=6m^2/(\pi^2N)$.
\end{proof}
\noindent Two more estimators follow from \eqref{eq:weibull}: across two gauges, $m=\ln(L_2/L_1)/\ln(\bar\sigma_1/\bar\sigma_2)$; from the scatter at
one gauge, $\mathrm{CoV}=\sqrt{\Gamma(1+2/m)/\Gamma(1+1/m)^2-1}\approx1.2/m$.
\nums{$z=2$, $\Delta m=4.5$: $N=12$ if $m=10$, $N=25$ if $m=14.5$. Over a decade of length the mean strength falls $1.26\times$ ($m=10$) or
$1.17\times$ ($m=14.5$). From scatter: commercial yarn at 110 mm, CoV 0.150, $m\approx7.9$ \cite{dessureault}; continuous yarn, CoV 0.040,
$m\approx32$ if the $\pm$ is one standard deviation (producer's own data, grade C) \cite{liang23}. A yarn is a bundle, so Daniels averaging
(I.9) raises its $m$ and lowers its mean: neither number transfers to the 80 GPa bundles.}
\verdict{K1 test: 12--25 bundles at each of two gauges a decade apart (20 and 200 mm); $m$ by maximum likelihood and by the two-gauge ratio. Cut
the bundle rim route if $m<10$.}

\subsection{Shear lag: what the matrix must hold}
\meta{[D] run 4 $\cdot$ gap 2, composite route $\cdot$ from I.10}
\begin{statement}
A fibre or bundle of diameter $d$, reloaded from a break through an interface that yields at shear stress $\tau$, regains stress linearly,
$d\sigma/dx=4\tau/d$, and carries $\sigma_f$ again $\sigma_fd/4\tau$ from the break on each side. The ineffective zone is the Kelly--Tyson length
\begin{equation}
\delta=\frac{\sigma_f\,d}{2\tau}\qquad\Longleftrightarrow\qquad\tau\ge\frac{\sigma_f\,d}{2\delta}.\label{eq:shearlag}
\end{equation}
\end{statement}
\begin{proof}Force balance on a slice $dx$: $\tfrac{\pi d^2}{4}\,d\sigma=\pi d\,\tau\,dx$ \cite{kellytyson}.\end{proof}
\nums{I.10 needs $\delta\lesssim0.1\un{mm}$. With $\sigma_f=80\un{GPa}$: $d=1\,\mu$m needs $\tau\ge400\un{MPa}$; $d=100\un{nm}$, 40 MPa;
$d=10\un{nm}$, 4 MPa. At $\tau=50\un{MPa}$ (an epoxy-class interface, assumed): $\delta=0.8\un{mm}$ for 1 $\mu$m bundles, 0.08 mm for 100 nm.}
\verdict{The composite route needs bundles about 100 nm across or thinner, each reloaded by the matrix; micron bundles need an interface no
polymer provides. This is the R\&D item between ``80 GPa at 2 cm'' and a hoop.}

\subsection{Cooldown: contraction mismatch and the pre-compression window (bridge B1)}
\meta{[D]+[E] run 4 $\cdot$ gap 6 (K2) $\cdot$ from I.6; supplies the pre-compression term I.6 left out}
\begin{statement}
A REBCO pack and a band shrink freely by $s_1$ and $s_2$ from room temperature to 4.2 K; their axial stiffnesses are $k_i=E_iA_i$.
\begin{enumerate}
\item Bonded or co-wound, both take the common shrinkage $\bar s=(k_1s_1+k_2s_2)/(k_1+k_2)$ and the tape is left with
\begin{equation}
\varepsilon_1=\frac{k_2\,(s_1-s_2)}{k_1+k_2}\qquad(\text{positive}=\text{tension}).\label{eq:cte}
\end{equation}
\item An unbonded overband loses contact by $s_1-s_2$ unless it is wound warm with pre-strain $\varepsilon_{\rm pre}\ge s_1-s_2$.
\item A tape pre-compressed by $\varepsilon_c$ at 4.2 K can take $\varepsilon_{\rm irr}+\varepsilon_c$ of tension, so the strain-matched field ceiling
\eqref{eq:strain} rises by $\sqrt{1+\varepsilon_c/\varepsilon_{\rm irr}}$.
\end{enumerate}
\end{statement}
\begin{proof}
(1) Equal final length and zero net force: $k_1(\bar s-s_1)+k_2(\bar s-s_2)=0$; the tape's strain is $s_1-\bar s$. (2) The free shrinkages differ by
$s_1-s_2$; contact needs that much stored band strain. (3) $B_{\max}^2\propto E\varepsilon$ in \eqref{eq:strain}.
\end{proof}
\inputs{$s_1=0.26\,\%$ for a Hastelloy-substrate tape pack (run-4 search, 0.26--0.27\,\%), stainless 304 $0.30\,\%$ and unidirectional carbon
along the fibre $\approx0$ \cite{ekin}; $E_1=150$ GPa (tape, assumed, 130--170), M55J hoop $0.6\times540$ GPa, stainless 193 GPa; equal areas;
$\varepsilon_{\rm irr}=0.63\,\%$; $\varepsilon_c=0.2\,\%$ (assumed, low end of reported compressive limits).}
\nums{Carbon co-wound: $\varepsilon_1=+0.178\,\%$, 28\,\% of the tensile window spent by cooling alone. Stainless co-wound: $-0.023\,\%$, a small
pre-compression. Carbon overband: $\varepsilon_{\rm pre}\ge0.26\,\%$, i.e.\ 842 MPa of warm band stress, pressing 42 MPa on the pack at
$t/R=0.05$. Pre-compression of 0.2\,\%: $1.32\times$ the strain window and $1.15\times$ the field ceiling (Hastelloy-class 59 $\to$ 68 T).}
\verdict{B1 revised. A carbon or nanotube band (the K1$\times$K2 bridge) either separates on cooldown or, bonded, pre-tensions the tape by
$\sim$0.18\,\%. Prefer a co-wind that shrinks at least as much as the tape (stainless, Inconel-class); use carbon only as a warm-pre-tensioned
overband sized by (2). Co-winding also lowers screening-current-induced stress, with diminishing returns as the co-wind thickens
\cite{scis22,yan19}.}

\setcounter{section}{3}\setcounter{subsection}{3}
\subsection{Thrust per watt: the ladder from light to ion wind (Biefeld--Brown)}
\meta{[D] run 4 $\cdot$ exploration, cut $\cdot$ from III.1, III.2}
\begin{statement}
A device that ejects mass at speed $v\ll c$ gets $F/P=2/v$ ($F=\dot mv$, $P=\tfrac12\dot mv^2$). Ions of mass $m$ and charge $q$ through voltage
$V$: $F/P=\sqrt{2m/(qV)}$. Electrons: $F/P=\gamma\beta/[(\gamma-1)c]$. Light: $1/c$. In a gas, ions drifting with mobility $\mu$ across a gap $d$
hand their momentum to neutrals:
\begin{equation}
F=\frac{I\,d}{\mu},\qquad \frac FP=\frac{d}{\mu V}.\label{eq:ehd}
\end{equation}
\end{statement}
\begin{proof}
The drift speed is $\mu E$, so the space charge per unit length is $I/(\mu E)$ and the electric force on it is $\int_0^d\frac{I}{\mu E}E\,dx=Id/\mu$;
in steady state collisions pass this to the gas. With no gas, Theorem~\ref{thm:com} leaves only what is ejected.
\end{proof}
\inputs{air $\mu=2\e{-4}\un{m^2/(V\,s)}$; $d=3$ cm, $V=30$ kV, $I=1$ mA (lifter-like, assumed); singly charged ions.}
\nums{Air: 0.15 N, 5 mN/W. Vacuum: protons 0.83 $\mu$N/W, Cu$^+$ 6.6, Xe$^+$ 9.5; 30 keV electrons 0.020 $\mu$N/W; light 3.3 nN/W. A vacuum claim
of 0.1 mN/W sits $10\times$ above xenon ejection: it needs a visible heavy-ion exhaust, or it is an artefact.}
\verdict{Biefeld--Brown cut: the force is ion wind and vanishes in vacuum \cite{talley,tajmarbb}. Any closed electrostatic device sits at zero
(III.1); one that ejects sits on the $2/v$ line.}

\setcounter{section}{5}\setcounter{subsection}{5}
\subsection{What field enhancement buys: the driven-volume theorem (B5)}
\meta{[P] run 4 $\cdot$ gap 9 (B5) $\cdot$ from V.3, V.4}
\begin{theorem}\label{thm:enh}
A metasurface that raises the local field by $F$ over a fraction $f$ of the illuminated area cuts the incident intensity needed for a given local
field by $F^2$, leaves the absorbed power per driven volume unchanged, and changes the driven volume per incident watt by
\begin{equation}
G=f\,F^2 .\label{eq:enh}
\end{equation}
\end{theorem}
\begin{proof}
The absorbed power density $\tfrac12\sigma_1E_{\rm loc}^2$ is fixed by the target field. Bare film: power $I_0A$ drives volume $At$. Metasurface:
power $I_0A/F^2$ drives volume $fAt$. Take the ratio of volume per watt.
\end{proof}
\inputs{hold fluence 4 J/m$^2$ \cite{rowe} delivered in 1 ps (assumed): local field 0.55 MV/cm. SRR gap 2.2 $\mu$m, enhancement $\propto1/g$
\cite{fan13}; SRR arm 4--10 $\mu$m in a 50 $\mu$m cell, dipole gap $1\times5\,\mu$m in a $170\times100\,\mu$m cell (geometry assumed).}
\nums{A tenfold intensity saving needs $F=\sqrt{10}=3.2$ (incident 0.17 MV/cm). Measured $F$: split-ring resonator $\sim$4 for a broadband pulse and
$\sim$15 on resonance \cite{fan13}; 1 $\mu$m-gap rings $\sim$20, reaching 3.5 MV/cm \cite{fan12}; dipole gap 170 at the metal edge, 60 mid-gap
\cite{tifee}. So $G=0.06$--0.14, 0.8--2.0, 0.6--1.6 and 1.1--1.7. Field ladder in a gap: needed 0.55 MV/cm; GaAs impact ionization 0.64
\cite{fan13}; electromigration damage seen above 2 \cite{electromig}; 3.5 reached \cite{fan12}; field emission 10 \cite{tifee}.}
\verdict{B5 kept only as a \emph{test enabler}: a smaller source can reach the hold field in micron gaps, $3.6\times$ below the damage seen. Cut as a
power route for gap 9: $G\approx0.1$--2, and the driven islands cover $\lesssim1\,\%$ of the film without connecting, so a test needs a
gap-resolved probe, not a transport measurement.}

\subsection{The London moment and the Tate anomaly}
\meta{[D] run 4 $\cdot$ exploration lead K5 $\cdot$ from V.1; meets the gravity path cut in run 3}
\begin{statement}
Deep inside a superconductor rotating at $\boldsymbol\omega$ no current flows relative to the lattice, so the pairs co-rotate, and fluxoid
quantization ($n=0$) gives a uniform field
\begin{equation}
\mathbf B=-\frac{2m^*}{e^*}\,\boldsymbol\omega\approx-\frac{2m_e}{e}\,\boldsymbol\omega=-1.137\e{-11}\un{T}\times\frac{\omega}{\rm rad\,s^{-1}} .\label{eq:london}
\end{equation}
\end{statement}
\begin{proof}
Canonical momentum $\mathbf p=m^*\mathbf v_s+e^*\mathbf A$ with $\mathbf v_s=\boldsymbol\omega\times\mathbf r$. On a loop of area $S$ inside,
$\oint\mathbf p\cdot d\mathbf l=0$ and $\oint(\boldsymbol\omega\times\mathbf r)\cdot d\mathbf l=2\omega S$, so $2m^*\omega S+e^*BS=0$. Pairs have
$m^*/e^*=m_e/e$ up to small corrections (theory for Nb: $m^*/2m_e=0.999992$).
\end{proof}
\nums{1000 rad/s gives 11.4 nT, one flux quantum through a 1 cm-radius ring per 0.58 rad/s. Tate \emph{et al.} (Nb rings, 1989--90):
$m^*/2m_e=1.000084(21)$, $+92$ ppm, $4.4\sigma$ \cite{tate}; an average over rotation-rate bins gives $1.000075(25)$. The gravitomagnetic
explanation needs $B_g=1.84\e{-4}\,\omega$ \cite{tajmargm}, about $5\e{22}$ times a classical estimate for a 1 cm Nb cylinder (E), and the same
authors' density-ratio model falls $23\times$ short of Tate's value \cite{tajmargm2}. For the craft: a REBCO band at 300 rad/s makes 3.4 nT.}
\verdict{K5 kept as a measurement lead only: an unrepeated $4\sigma$ precision result with no accepted explanation. Deciding experiment: repeat at
$\le10$ ppm; SQUID flux resolution is not the limit (10 ppm of 1728 flux quanta is 0.017 $\Phi_0$), the ring's cold dimensions and stray flux
are (E). Engineering consequence for the craft: none.}

%=====================================================================================================
\section*{Run 5 additions to chains I and III}
\addcontentsline{toc}{section}{Run 5 additions: I.15--I.17, III.5}
\noindent\textbf{Corrections found in run 5.} (i) I.11: the best metre-scale fibre is Toray T1200 (8.0 GPa at 1.82 g/cm$^3$, announced October
2023 \cite{toray23}), 4.40 N/tex, so gap 2 is $9.0\times$ short in $\sigma/\rho$ ($3.0\times$ in speed), not $10\times$; the run-4 value came from a
stale table. (ii) I.14: measured cooldown offsets imply a free-tape shrinkage of 0.31--0.32\,\% to 4.2 K, not 0.26\,\% (I.16). (iii) I.7 and
ledger 16: ``rims $\le7.7$ km/s'' holds for a rim that carries its own load; with a stator the virial bound applies to the whole craft (I.17).
Numbers below come from \texttt{run5\_math.py}; the run opened with a blind re-derivation of seven run-4 numbers (\texttt{canary5.py}, 7 of 7).

\setcounter{section}{1}\setcounter{subsection}{14}
\subsection{Schulz--Flory: the length budget of continuous nanotubes}
\meta{[D] run 5 $\cdot$ gap 2 (K1) $\cdot$ from I.8, I.13}
\begin{statement}
If a growing tube's catalyst survives each growth step $\Delta$ with probability $\alpha$, the number of tubes reaching length $L$ is
\begin{equation}
N(L)=N_0\,\alpha^{L/\Delta}=N_0\,2^{-L/L_{1/2}},\qquad L_{1/2}=\frac{\Delta\ln2}{\ln(1/\alpha)},\qquad \bar L=\frac{L_{1/2}}{\ln2}.\label{eq:sf}
\end{equation}
\end{statement}
\begin{proof}Deactivation is memoryless, so survival probabilities multiply step by step: a geometric distribution, the Schulz--Flory
law of polymer chain lengths \cite{zhang13}.\end{proof}
\inputs{$L_{1/2}=74.4$ mm for semiconducting ultralong tubes \cite{zhu19}; longest grown 550 mm \cite{zhang13}; tube diameter 2 nm, useful
length 2 cm, $V_f=0.6$, rim section 1 cm$^2$ (assumed).}
\nums{Fraction reaching 20 mm: 83\,\%; 200 mm: 16\,\%; 1 m: $9.0\e{-5}$; the 12.57 m hoop: $10^{-50.8}$. Deactivation probability
$9.3\e{-6}$ per $\mu$m of growth. A path round the hoop crosses about 117 tube ends ($\bar L=107$ mm), and a 1 cm$^2$ rim section needs
$1.2\e{16}$ tubes 2 cm long. With ordinary (defective) tubes, nanoscale Weibull statistics give $m\approx2.7$ and 34 GPa at 100 nm
\cite{pugno06}: a single 1 m tube would hold 0.08 GPa, and simulated macroscopic cables stay below about 10 GPa (7.7 N/tex) \cite{pugno08}.}
\verdict{No continuous tube spans a rim: it must be a spliced composite (I.13) of defect-free cm-long tubes, about $10^{16}$ per cm$^2$ of
section. The 20/200 mm test of I.12 stays feasible (16\,\% of tubes reach 200 mm). K1 is waiting on data, with its prior lowered again.}

\subsection{Bonded bands through the cooldown: the measured offset and the radial tension}
\meta{[D]+[E] run 5 $\cdot$ gap 6 (K2) $\cdot$ from I.14}
\begin{statement}
(a) A tape bonded to a much stiffer holder of shrinkage $s_h$ keeps the mechanical strain $\varepsilon\approx s_{\rm tape}-s_h$ (I.14 with
$k_2\gg k_1$). (b) A bonded outer band that shrinks less than the pack, by $\Delta s$, pulls it outward with radial tension
\begin{equation}
q=\frac{\Delta s}{R/(E_wt_w)+R/(E_bt_b)} .\label{eq:qrad}
\end{equation}
\end{statement}
\begin{proof}(a) is \eqref{eq:cte} as $k_2/k_1\to\infty$. (b) A thin ring of hoop stiffness $Et$ under a radial pressure $q$ moves by $qR^2/(Et)$;
the pack's outward and the band's inward displacements must add up to the free mismatch $R\,\Delta s$.\end{proof}
\inputs{Ti-6Al-4V Walters spring, 0.17\,\% to 4.2 K; zero-strain offsets for five makers' tapes 0.144--0.152\,\% at 4.2 K; irreversible limits
0.45--0.72\,\% at 4.2 K, 19 T \cite{barth}. Pack $R=0.1$ m, 20 mm at 150 GPa; M55J band 5 mm at $0.6\times540$ GPa (assumed geometry).}
\nums{(a) Free-tape shrinkage $0.17+0.144\ldots0.152=0.31$--$0.32\,\%$. Re-running I.14: a bonded carbon co-wind leaves $+0.18$ to
$+0.22\,\%$, i.e.\ 31--48\,\% of the five tapes' windows; stainless leaves $-0.02$ to $+0.01\,\%$, neutral. (b) $q=27$--$34$ MPa, against
c-axis delamination strengths of 18--53 MPa \cite{parkcaxis}; radial tension from contraction mismatch is the known cause of delamination in
impregnated pancakes \cite{delam}.}
\verdict{I.14's cooldown offset is a measured fact (grade B). Bonded carbon or nanotube bands are cut twice: strain window and delamination.
Metal co-winds are neutral; K2 reduces to one number, the free shrinkage of the actual conductor.}

\subsection{A rotor held by a stator: the virial theorem for a closed craft with magnetic support (bridge B6)}
\meta{[P] run 5 $\cdot$ gap 2 (new bridge B6) $\cdot$ from I.2, I.4, I.7; corrects the reading of I.7}
\begin{theorem}\label{thm:stator}
For a closed system of rotating matter, stationary matter and steady axisymmetric fields, $\int\tr\sigma\,dV$ over all matter equals $2T+U$.
A rotor can be free of stress where $\mathbf J\times\mathbf B=\rho_rv^2/R$ locally; the stator then carries $2T+U$, so
$M_{\rm stator}\ge\rho_s(2T+U)/\sigma_s$. A fully supported thin rotor needs
\begin{equation}
S_m\equiv\frac{J_e\,B\,R}{\rho_r}\;\ge\;v^2 :\label{eq:sm}
\end{equation}
its speed is bounded by this magnetic specific strength, not by its own $\sigma/\rho$.
\end{theorem}
\begin{proof}
$\partial_j(\sigma_{ij}+T_{ij})=\rho a_i$, with no $\partial_t\mathbf g$ term because the fields are steady. Multiply by $x_i$ and integrate over
all space: $-\int(\tr\sigma+\tr T)\,dV=\int\rho\,\mathbf x\cdot\mathbf a\,dV=-2T$, and $\tr T=-u$. For the rotor, the radial momentum balance
of a thin ring is $J_eB=\rho_rv^2/R$.
\end{proof}
\begin{corollary}[tilt]\label{cor:tilt}
If a fraction $\beta$ of a thin ring's centripetal load is magnetic, in a uniform axial field, then $4I_\perp K=\beta L^2$, where $L$ is the
spin angular momentum and $K$ the magnetic overturning stiffness. The ring is stable against tilting iff $\beta\le1$: full magnetic support
is exactly marginal, and the margin is $1-\beta$.
\end{corollary}
\begin{proof}
For tilt $\xi=\theta_x+i\theta_y$: $I_\perp\ddot\xi-iL\dot\xi-K\xi=0$, so $s=[iL\pm\sqrt{4I_\perp K-L^2}\,]/(2I_\perp)$, stable iff
$L^2\ge4I_\perp K$. An inward force needs the ring's moment anti-parallel to $\mathbf B$, so $K=|I|\pi R^2B$ with $|I|B=\beta M\omega^2/2\pi$, giving
$K=\tfrac12\beta M\omega^2R^2$; with $I_\perp=\tfrac12MR^2$ and $L=MR^2\omega$, $4I_\perp K=\beta M^2\omega^2R^4=\beta L^2$. A
flux-conserving ring changes $K$ only at second order in a uniform field.
\end{proof}
\inputs{$R=2$ m, 30\,000 rpm, $v^2=3.95\e7$ J/kg; REBCO $\rho_r=8900$; $J_e$ at 4.2 K, 20 T: 0.6 kA/mm$^2$ commercial \cite{kar19}, $\approx1$
kA/mm$^2$ record tape (2019) \cite{aries}, $\approx4.6$ kA/mm$^2$ thick-film R\&D samples \cite{selva24}; $J_e$ held to 40 T for
$\mathbf B\parallel ab$ (optimistic; $I_c$ falls $\sim$22\,\% to 31 T); rotor 2 mm $\times$ 100 mm (assumed).}
\nums{$S_m$ ranges from $2.7\e6$ (commercial, 20 T: $15\times$ short) through $9.0\e6$ (record tape, 40 T: $4.4\times$) to $3.1\e7$ and
$4.1\e7$ (R\&D, 30 and 40 T: $1.3\times$ short, then meets it); the best fibre is $9.0\times$ short. The rotor weighs 22 kg and holds
$T=442$ MJ; the stator needs $\ge289$ kg of T1200 hoop ($\ge6.1$ t at REBCO-coil $\sigma/\rho$); a 5 cm annular field gap stores only 10--40
MJ. A REBCO-only rotor can carry at most 0.4\,\% of the load itself ($\beta\ge0.9963$), so its tilt margin is 0.4\,\%.}
\verdict{B6 kept (D). It turns gap 2 from a materials wall into an engineering target: within $1.3\times$ at 30 T with lab-best conductor,
$4$--$15\times$ short with today's tape. Tilt stability decides it (exactly marginal at $\beta=1$): field shaping or active control must add
stiffness. Smallest deciding experiment: a REBCO ring of 10 cm radius in 10 T spun to 500 m/s (full magnetic support would need $\rho v^2/BR=2.2$ kA/mm$^2$), $1.3\times$ its own burst speed (381 m/s), at
$\beta\approx0.5$, which needs about 1.1 kA/mm$^2$. Prior art for levitated superconducting rings exists, with feedback and eddy-current
damping \cite{cooper70}; superconductors and spin both evade Earnshaw's static theorem \cite{post}.}

\setcounter{section}{3}\setcounter{subsection}{4}
\subsection{The flyby anomaly against frame dragging}
\meta{[E] run 5 $\cdot$ exploration, cut $\cdot$ from III.1}
\begin{statement}
Anderson's empirical formula is $\Delta V_\infty/V_\infty=K(\cos\delta_{\rm in}-\cos\delta_{\rm out})$ with $K=2\omega_ER_E/c=3.099\e{-6}$
\cite{anderson08}. The Lense--Thirring acceleration of a probe near the rotating Earth is of order $a\approx2GJ_Ev/(c^2r^3)$.
\end{statement}
\inputs{$J_E=5.86\e{33}$ kg m$^2$/s; NEAR perigee 539 km, 12.7 km/s (assumed speed); about 1000 s near perigee.}
\nums{NEAR: formula 13.28 mm/s against 13.46 observed. Frame dragging: $3.4\e{-10}$ m/s$^2$, about $3\e{-4}$ mm/s over 1000 s, $4\e4$ times
too small. The formula's predictions failed: no anomaly at Galileo's second flyby or Rosetta's second and third; about 7 mm/s predicted for
Juno (2013), which showed none or a threshold signal \cite{iorio15,acedo18}. No velocity-independent force field reproduces the formula.}
\verdict{Cut as a lead for the craft: no mechanism, and the anomaly does not repeat. Revive only if a dedicated flyby, tracked continuously
through perigee, shows it.}

\subsection{Charged-dielectric thrust claims: what electrostatics alone gives (X1)}
\meta{[D]+[E] run 5, side research (not scored) $\cdot$ exploration lead X1 $\cdot$ from I.1, III.1, III.2, III.4}
\begin{statement}
A conductor or electret face carrying surface charge $\sigma$ is pulled towards nearby conductors with the pressure $p=\sigma^2/2\varepsilon_0$, so a
face of area $A$ feels
\begin{equation}
F=\frac{\sigma^2A}{2\varepsilon_0},\qquad \sigma\le\varepsilon_0E_{\rm bd}\quad(\text{air, }E_{\rm bd}\approx3\un{MV/m}:\ \sigma\le26.6\,\mu\mathrm{C/m^2}).\label{eq:espress}
\end{equation}
A sealed device that ejects nothing has $F\le P/c$ (III.2), which is zero with the power off.
\end{statement}
\begin{proof}The field just outside a charged surface facing a conductor is $\sigma/\varepsilon_0$ and the normal Maxwell stress there is
$\varepsilon_0E^2/2$ (I.1); breakdown caps $E$. The second part is Theorem~\ref{thm:com}.\end{proof}
\inputs{claimed: $\sim$1 mN at $\sim$480 V, $\sim$2.5 mN in vacuum on a pendulum with grounded walls, 10 mN peak from stacked devices of 30--40 g,
thrust persisting on stored charge with the power off \cite{buhler}; face $5\times5$ cm (assumed).}
\nums{10 mN on 25 cm$^2$ needs $\sigma=8.4\,\mu$C/m$^2$, 32\,\% of the air-breakdown limit: ordinary attraction to walls and images is the right size,
and grounding the walls creates it rather than removing it. With nothing ejected and no power, $P/c=0$; photon thrust of 10 mN would need 3.0 MW.
Ejecting N$_2^+$ at 30 kV would give 4.4 $\mu$N/W (III.4). Weight: conductors of $\sim$1 kg charged from 10 V to 200 kV, both polarities, changed
weight by less than 2 g, i.e.\ $2\e{-3}$ \cite{hathaway21}.}
\verdict{X1 kept as a lead, grade E, contested; no independent test found. Deciding experiment: a battery-powered device sealed in a conducting
box on a torsion balance in high vacuum, with polarity reversal and an uncharged dummy; kill if the force vanishes inside the box, and a real
effect must exceed 3.34 nN/W with nothing ejected. The same items cut St.~Clair's rotating electrostatic thruster (US 2003/0209637, cited by the
project's triangular-spacecraft filing): a steadily spinning axisymmetric charge radiates nothing (II.2) and cannot move its centre of energy
(III.1).}


%=====================================================================================================
% Pais gap loop -- mathematics run sheet, part 4: run-6 items (loaded by \input from part 2, after part 3).
% Items keep their chain numbers (counters set below). Numbers: run6_numbers.json from run6_math.py.
%=====================================================================================================
\section*{Run 6 additions to chains I and V}
\addcontentsline{toc}{section}{Run 6 additions: I.18--I.21, V.8}
\noindent\textbf{Notes on earlier items found in run 6.}
\begin{itemize}
\item \emph{I.6 and I.14 (correction).} A band wound warm with pre-strain spends that pre-strain from its own failure strain. The ``72 T'' quoted for
an M55J band since run 4 is unreachable; I.21 gives 50 T. The error is logged.
\item \emph{I.11 (open input).} The 2 cm bundles' 80 GPa is computed on the tubes' own cross-section; on the engineering cross-section the
same bundles carry 43 GPa \cite{physt18}. Their specific strength is 61.5 N/tex only if 80 GPa goes with a density of at most 2{,}030
kg/m$^3$; if the assumed 1{,}300 kg/m$^3$ goes with the engineering area it is 33 N/tex, below the 39.5 N/tex the rim needs. To be checked
against the paper's methods.
\item \emph{I.17 (conductor).} The 4.6 kA/mm$^2$ R\&D value comes from thick REBCO films. Torque magnetometry to 45 T found that films
3.0--5.0 $\mu$m thick always flux-jump at 4.2 K during field sweeps, films 1.4--1.8 $\mu$m sporadically, a 1.12 $\mu$m film never;
production tapes carry 700--2000 A at 30 T in the ab plane (width not stated) \cite{jaroszynski25}.
\item \emph{III.6 (X1).} No independent replication had been published as of May 2026; the company's own patent cites a vacuum null for
asymmetric capacitors (Canning, 2004) \cite{exodus891}; a force that persists after power-off matches a balance-drift signature seen in
EMDrive tests \cite{tajmar15}. X1 stays at grade E.
\end{itemize}
Numbers below come from \texttt{run6\_math.py}; the run opened with a blind re-derivation of thirteen run-5 numbers (\texttt{canary6.py}, 13
of 13).

\setcounter{section}{1}\setcounter{subsection}{17}
\subsection{A current ring in an axisymmetric field: the translational stiffness has zero trace (B6)}
\meta{[P] run 6 $\cdot$ gap 2 (bridge B6) $\cdot$ from I.17; Earnshaw's theorem for a rigid ring}
\begin{theorem}\label{thm:ring}
A rigid thin ring of radius $R$ carries a fixed current $I$ in a static, current-free, axisymmetric field whose axial component at the ring,
$B_0$, pulls each element inward with force $f=|IB_0|$ per unit length. With field index $n=-(R/B_z)\,\partial B_z/\partial r$, the stiffnesses
against translation are
\begin{equation}
k_x=k_y=-\pi nf,\qquad k_z=2\pi nf,\qquad k_x+k_y+k_z=0 .\label{eq:ringk}
\end{equation}
For any rigid loop in any current-free field, the trace of the translational stiffness at fixed current is zero: one direction at least is
unstable unless all three are neutral.
\end{theorem}
\begin{proof}
At fixed current the energy is $U=-I\Phi_{\rm ext}(\boldsymbol\xi)$, with $\Phi_{\rm ext}(\boldsymbol\xi)=\int B_z(\mathbf x+\boldsymbol\xi)\,dA$ for a
translation $\boldsymbol\xi$. Hence $\sum_i\partial^2U/\partial\xi_i^2=-I\int\nabla^2B_z\,dA=0$, because $\nabla^2\mathbf B=0$ where there is no
current. For the ring: a shift $x$ puts the element at angle $\phi$ at radius $R+x\cos\phi$, where the inward force is $f(1-nx\cos\phi/R)$;
summing the $x$-components, $F_x=\int_0^{2\pi}f(1-nx\cos\phi/R)(-\cos\phi)R\,d\phi=\pi nfx$. Curl-free means $\partial B_r/\partial z=\partial
B_z/\partial r$, so $B_r(R,z)=-nB_0z/R$; the force on the azimuthal current is $-IB_r$ per unit length along $z$, and
$F_z=-2\pi R\,IB_r=-2\pi nfz$ (with $IB_0=-f$).
\end{proof}
\begin{corollary}[rates]\label{cor:rates}
If a share $\beta$ of the centripetal load of a ring of mass $M$ spinning at $\omega$ is magnetic, $f=\beta Mv^2/(2\pi R^2)$ and
\begin{equation}
k_\perp=-\tfrac12 n\beta M\omega^2,\qquad k_z=n\beta M\omega^2 .\label{eq:ringrates}
\end{equation}
For $n<0$ the ring is held in its plane at $\omega\sqrt{|n|\beta/2}$ and falls out of it vertically at the rate $\omega\sqrt{|n|\beta}$; for $n>0$
the roles swap. The instability is a fixed fraction of the spin rate.
\end{corollary}
\noindent\emph{Flux conservation.} A closed superconducting ring keeps $LI+\Phi_{\rm ext}$ fixed, so $U=(\Phi-\Phi_{\rm ext})^2/2L$ and each stiffness
gains $(\partial\Phi_{\rm ext}/\partial\xi)^2/L\ge0$. At the symmetric position $\partial\Phi_{\rm ext}/\partial x=\partial\Phi_{\rm ext}/\partial z=0$,
so it gains nothing. Off the midplane by $z_0$, $\partial\Phi_{\rm ext}/\partial z=-2\pi RB_r=2\pi nB_0z_0$, and passive vertical stability needs
$z_0\ge\sqrt{IL/(2\pi|n|B_0)}$.
\inputs{rotor of I.17: $R=2$ m, 30\,000 rpm, $2\times100$ mm REBCO section, $\rho=8900$; $\beta=1$; single-turn-equivalent inductance
$L=\mu_0R(\ln 8R/a-2)$ with $a=8$ mm. Numerical check: two coaxial loops of radius 2.5 m at $z=\pm0.6$ m, $10^7$ A-turns each (assumed),
and a test ring of $10^4$ A-turns.}
\nums{$M=22.4$ kg, $\omega=3142$ rad/s, centripetal acceleration $2.0\e6\,g$, $f=3.51\e7$ N/m (1.17 MA-turns at 30 T, i.e.\ 5.86 kA/mm$^2$).
At $|n|=0.1$ the unstable mode grows at 993 s$^{-1}$ ($n<0$, vertical) or 702 s$^{-1}$ ($n>0$, in plane); at $|n|=10^{-3}$, at 99 or 70
s$^{-1}$. Holding the rotor within 20 $\mu$m vertically takes 442 N at $|n|=0.1$ and 4.4 N at $10^{-3}$ (its weight is 219 N). A levitated
dipole held against gravity grows at roughly $\sqrt{g/\ell}=3$--10 s$^{-1}$ (estimate, $\ell=0.1$--1 m). Numerical check: the coil pair gives $B_z=5.62$ T and $n=0.666$ at the ring;
summing $I\,d\mathbf l\times\mathbf B$ around displaced rings gives $k_x=-1.17733\e5$ and $k_z=2.35463\e5$ N/m, equal to \eqref{eq:ringk} to
$10^{-5}$, with a trace of $-1\e{-5}k_z$. Flux conservation: $L=1.41\e{-5}$ H; passive vertical stability would need $z_0=0.94$ m at $|n|=0.1$
and 30 T, and still 0.30 m at $|n|=1$.}
\verdict{No field shape makes B6 passively stable in translation, and flux conservation does not rescue it at any practical offset. One
translation must be held actively, as levitated dipoles do: APEX-LD holds a REBCO floating coil to $\sigma_z<20\,\mu$m for over three hours
\cite{apexld}, and LDX flew a 560 kg coil on feedback for more than 80 hours \cite{ldx07,ldxfpa}. Because the rate scales with spin, the stator
field should be uniform to $|n|\lesssim10^{-3}$ over the rotor: about 100 s$^{-1}$ and a few newtons of control force, a kilohertz loop on a 22
kg ring. B6 kept (D), with this control requirement added.}

\subsection{Supercritical whirl: damping on the rotor destabilizes, damping on the stator stabilizes}
\meta{[D] run 6 $\cdot$ gap 2 (B6 kill search) $\cdot$ from I.18}
\begin{proposition}\label{prop:whirl}
A rotor of mass $M$ on an in-plane suspension of stiffness $k$ ($\omega_n=\sqrt{k/M}$), with non-rotating damping $c_n$ and damping $c_r$ in parts
spinning at $\Omega$, whirls forward unstably when
\begin{equation}
\Omega>\omega_n\Big(1+\frac{c_n}{c_r}\Big).\label{eq:whirl}
\end{equation}
\end{proposition}
\begin{proof}
With $\zeta=x+iy$, a damping force $-c_r(\dot\zeta-i\Omega\zeta)$ acts in the rotating frame, so $M\ddot\zeta+(c_n+c_r)\dot\zeta+(k-i\Omega c_r)\zeta=0$.
Put $\zeta=e^{st}$ with $s=i\omega_n+\delta$ and keep first order in the damping: $2iM\omega_n\delta=-i[(c_n+c_r)\omega_n-\Omega c_r]$. The forward
mode decays only while $(c_n+c_r)\omega_n>\Omega c_r$ \cite{genta,ismb10}.
\end{proof}
\nums{With the in-plane stiffness of \eqref{eq:ringrates} ($n<0$, $\beta=1$), $\Omega/\omega_\perp=\sqrt{2/|n|}$: 4.5 at $|n|=0.1$, 14 at 0.01 and
45 at $10^{-3}$, so the stator-side damping must exceed the rotor-side damping 3.5, 13 and 44 times. Eigenvalue check with $c_n=3c_r$: the real
part of the forward root goes from $-0.0089$ to $+0.0089$ s$^{-1}$ as $\Omega$ crosses $4\omega_n$ by $\pm1\,\%$.}
\verdict{A uniform field eases I.18's control loop but makes the rotor more supercritical. B6 needs its rotor nearly loss-free
(rotating damping comes from the rotor's own losses, I.20) and its damping on the stator side: eddy-current dampers or active damping in
the control loop, as rotordynamics prescribes for supercritical machines \cite{genta}.}

\subsection{Ripple loss in a spinning REBCO rotor}
\meta{[D]+[E] run 6 $\cdot$ gap 2 (B6) $\cdot$ from I.17, I.19}
\begin{statement}
A thin superconducting strip of width $2a$ and sheet critical current $K_c=J_cd$ in a perpendicular field of amplitude $H_m$ dissipates per
cycle and unit length \cite{brandt93}
\begin{equation}
Q=4\mu_0a^2K_cH_m\,g\!\Big(\frac{H_m}{H_c}\Big),\quad H_c=\frac{K_c}{\pi},\quad g(x)=\frac2x\ln\cosh x-\tanh x;\qquad
Q\simeq\frac{2\pi^3}{3}\,\frac{\mu_0a^2H_m^4}{K_c^2}\quad(H_m\ll H_c).\label{eq:brandt}
\end{equation}
\end{statement}
\begin{proof}[Small-amplitude limit]
$\ln\cosh x=x^2/2-x^4/12+O(x^6)$ and $\tanh x=x-x^3/3+O(x^5)$ give $g(x)=x^3/6+O(x^5)$; substitute $x=\pi H_m/K_c$.
\end{proof}
\inputs{rotor wound from 4 mm $\times$ 0.1 mm tape (assumed): 500 tapes of 12.57 m, 6{,}283 m in all; $K_c=J_e\times0.1$ mm; ripple at the
rotation frequency (500 Hz) or its 24th harmonic (12 kHz); zero transport current (a lower bound).}
\nums{For 1 W at 4.2 K the perpendicular ripple must stay below 6.3 mT at 500 Hz or 2.9 mT at 12 kHz with R\&D tape ($K_c=4.6\e5$ A/m), and
below 2.3 or 1.0 mT with commercial tape. The loss grows as the fourth power: 10 mT at 12 kHz dissipates 148 W at 4.2 K, at least 70 W of
room-temperature power per watt (Carnot) and several hundred in practice.}
\verdict{An axisymmetric stator has no azimuthal ripple in principle; joints, winding errors and rotor run-out must keep the perpendicular
component near or below $10^{-4}$ of 30 T. Transport current near $I_c$ and eddy currents in the stabilizer add to this bound, and every
watt the rotor dissipates is rotating damping (I.19).}

\subsection{A pre-strained band spends its pre-strain from its own failure strain (correction to I.6, I.14)}
\meta{[D] run 6 $\cdot$ gap 6 (K2) $\cdot$ from I.6, I.14, I.16}
\begin{statement}
A band wound warm with pre-strain $\varepsilon_{\rm pre}\ge s_{\rm tape}-s_{\rm band}$ (I.14) shares the tape's working strain once loaded, so
$\varepsilon_{\rm pre}+\varepsilon_{\rm work}\le\varepsilon_f$ and
\begin{equation}
\varepsilon_{\rm work}\le\min(\varepsilon_{\rm irr},\;\varepsilon_f-\varepsilon_{\rm pre}),\qquad
B_{\max}\approx\sqrt{2\mu_0E_{\rm band}\varepsilon_{\rm work}},\qquad \frac MU\ge\frac{\rho_{\rm band}}{E_{\rm band}\varepsilon_{\rm work}} .\label{eq:prestrain}
\end{equation}
\end{statement}
\begin{proof}Strain compatibility with the pre-strain added, then \eqref{eq:strain} and \eqref{eq:mmin}.\end{proof}
\inputs{Toray M55J fibre 4.02 GPa, 540 GPa, 0.8\,\%, 1.91 g/cm$^3$; unidirectional composite 338 GPa and 0.61\,\% strain, 1.90 GPa at 55\,\% $V_f$
\cite{toraym55j}; composite density 1{,}630 kg/m$^3$; $s_{\rm tape}=0.31$--0.32\,\% (I.16), carbon $\approx0$; Hastelloy-class co-wind 220 GPa
at $\varepsilon_{\rm irr}=0.63\,\%$.}
\nums{M55J band: $\varepsilon_{\rm work}=0.29$--0.30\,\%, $B_{\max}=49.6$--50.5 T, 1.6--1.7 t/GJ (54.5 T with run 4's 0.26\,\%).
Hastelloy-class co-wind: 59.0 T, 6.4 t/GJ. Run 4's 72 T used $0.6\times540$ GPa at the tape's 0.63\,\% with no pre-strain; the composite fails at
0.61\,\% before any pre-strain is applied. Cross-check: $338\un{GPa}\times0.61\,\%=2.06$ GPa, close to the measured 1.90 GPa laminate strength.}
\verdict{A carbon band gives no field gain over a metal co-wind; it is still 3.9 times lighter per stored joule. Gap 6 stays closed (it needs
about 20 T); K2 reduces to metal co-winds, plus carbon only where mass matters more than field.}

\setcounter{section}{5}\setcounter{subsection}{7}
\subsection{Strain as pressure: what a substrate can do for a superconductor (gap 8)}
\meta{[D]+[E] run 6 $\cdot$ gap 8, in-between pairing $\cdot$ from V.1}
\begin{statement}
A film clamped to a substrate with biaxial in-plane strain $\varepsilon$ and free out of plane has $\varepsilon_{zz}=-2\nu\varepsilon/(1-\nu)$,
in-plane stress $\sigma=E\varepsilon/(1-\nu)$ and volume strain $2\varepsilon(1-2\nu)/(1-\nu)$. The hydrostatic pressure that gives the same
volume change is
\begin{equation}
P_{\rm eq}=-K\frac{\Delta V}{V}=\frac23|\sigma|=\frac{2E|\varepsilon|}{3(1-\nu)} .\label{eq:peq}
\end{equation}
\end{statement}
\begin{proof}Hooke's law with $\sigma_{zz}=0$, and $K=E/3(1-2\nu)$.\end{proof}
\inputs{$E=200$ GPa, $\nu=0.3$ (assumed oxide film). Bilayer nickelate films superconduct at ambient pressure under $-1.2$ to $-2.0\,\%$ strain,
onset near 40--42 K \cite{nickfilm}; bulk La$_3$Ni$_2$O$_7$ needs about 14 GPa for about 80 K \cite{sun23}. Clamped Nb--H films stay elastic
to about $-10$ GPa below 6--8 nm \cite{nbh}. Superhydrides need 100--260 GPa (ledger, gap 8).}
\nums{$\varepsilon=1.2$, 2.0 and 3.0\,\%: $P_{\rm eq}=2.3$, 3.8 and 5.7 GPa; in-plane stress 3.4, 5.7 and 8.6 GPa. Strained nickelate films
supply about a quarter of the bulk pressure by volume (two-fifths by in-plane stress) and show about half the bulk $T_c$. Clamping supplies
at most about 10 GPa, 10--26 times less than superhydrides need.}
\verdict{Strain is a route only where the pressure needed is below about 10 GPa (nickelates, fine-tuning cuprates); it cannot bring a
superhydride to 1 atm. Gap 8's factor of 2 in $T_c$ is not closed by strain: cut for hydrides, kept as context for nickelates.}


%=====================================================================================================
% Pais gap loop -- mathematics run sheet, part 5: run-7 items (loaded by \input from part 2, after part 4).
% Items keep their chain numbers (counters set below). Numbers: run7_numbers.json from run7_math.py; run7_delay.json from run7_delay.py.
%=====================================================================================================
\section*{Run 7 additions to chains I and V, and S.4}
\addcontentsline{toc}{section}{Run 7 additions: I.22--I.24, V.9, S.4}
\noindent\textbf{Notes on earlier items found in run 7.}
\begin{itemize}
\item \emph{I.11 (correction; K1 closed).} The Supplementary Information of the 80 GPa bundle paper lists, for each bundle, the tube diameters
$d$, two areas and two stresses \cite{bai18si}. The areas are reproduced by $A_1=\sum\pi d\times0.34$ nm (the tube walls: 3.632 vs 3.630 nm$^2$
and 6.302 vs 6.299 nm$^2$ for two bundles) and $A_2=\sum\pi d^2/4$ (filled tubes: 4.555 and 9.621 nm$^2$, both exact). The headline 80 GPa is
therefore on the wall area, whose density is that of graphitic walls, about 2{,}260 kg/m$^3$; the ``engineering'' 43 GPa is on the filled-tube
area (implied mean $d\approx2.5$ nm, density about 1{,}215 kg/m$^3$). Either pairing gives the same specific strength, as it must:
\[
\frac{\sigma}{\rho}=\frac{F}{\lambda}=\frac{80\un{GPa}}{2260\un{kg\,m^{-3}}}=35.4\un{N/tex},
\]
where $\lambda$ is the mass per unit length. The rim needs $v^2=39.5$ N/tex, so even the best 2 cm bundles of near-perfect tubes are 0.90 times
short (5.95 against 6.28 km/s) before any length penalty. Runs 3--6 quoted 61.5 N/tex by pairing the wall-area stress with a bundle density; the
error is logged. K1 is cut as a rim route.
\item \emph{I.17 (conductor; open check closed).} Production-length REBCO now carries more than 1500 A at 4.2 K and 30 T with the field in the
ab plane, in tape 4.0 mm wide and 45 $\mu$m thick (30 $\mu$m substrate; the LBC6 coil, which reached 42.6 T) \cite{lbc6}. That is
$J_e=8.3$ kA/mm$^2$. B6 needs 5.86 kA/mm$^2$, i.e.\ it runs at $i=I/I_c=0.70$: 1{,}111 tapes of 12.57 m (14 km), 1{,}054 A each. In the rotor the
axial stator field lies in the tape plane, which is the ab orientation. Run 6's ``1.2--3.3 times short'' assumed 0.1 mm tape.
\item \emph{I.18 (cut revived).} Run 6 cut passive translational hold because the only flux-conserving element it considered was the rotor
itself, whose coupling vanishes at the symmetric position. A dedicated superconducting loop pair couples at that position; I.23 shows it
holds the unstable axis.
\item \emph{III.6 (X1).} The company's own slide shows 4.46 mN at 0.065 $\mu$A and 3{,}500 V (228 $\mu$W) at $1.4\e{-5}$ Torr \cite{exodusslides}.
Anything ejected can push at most $I\sqrt{2m_iV/e}$: 2.9 nN for N$_2^+$, 7.8 nN even for 200 u ions; photons give $P/c=0.76$ pN. The reading is
$5.7\e5$ times the ejecta bound, so it is a force between the article and its surroundings or an artefact. It equals the electrostatic pull
$\varepsilon_0E^2A/2$ on a 25 cm$^2$ face at $E=6.3\e5$ V/m, i.e.\ 3.5 kV across 5.5 mm; charge trapped in a dielectric (an electret) would
keep that pull after switch-off. No outside replication as of June 2026; X1 stays E.
\end{itemize}
Numbers below come from \texttt{run7\_math.py} and \texttt{run7\_delay.py}; the run opened with a blind re-derivation of fourteen run-6 numbers
(\texttt{canary7.py}, 14 of 14).

\setcounter{section}{1}\setcounter{subsection}{21}
\subsection{Holding the unstable axis actively: the delay and bandwidth it needs (B6)}
\meta{[P] run 7 $\cdot$ gap 2 (B6) $\cdot$ from I.18, I.19; prior art: levitated dipoles, tokamak vertical control, Beams' rotors, ESGs}
\begin{proposition}[critical delay]\label{prop:delay}
Let a translation obey $M\ddot x=k_ux+F$ with $\gamma=\sqrt{k_u/M}$, and let the controller apply proportional and derivative feedback with a
delay $\tau$: $F=-pk_u\,x(t-\tau)-dM\gamma\,\dot x(t-\tau)$. Some $(p,d)$ makes the motion asymptotically stable if and only if
\begin{equation}
\gamma\tau<\sqrt2 .\label{eq:delay}
\end{equation}
\end{proposition}
\begin{proof}[Derivation]
In units of $1/\gamma$ the characteristic function is $D(s)=s^2-1+(p+ds)e^{-s\tau}$. Roots cross the imaginary axis only on the D-curves
\cite{stepan89}: $s=0$ gives the line $p=1$, and $p<1$ is unstable because $D(0)<0$ while $D(s)\to+\infty$ on the positive real axis. $s=i\omega$
gives the curve
\[
p(\omega)=(1+\omega^2)\cos\omega\tau,\qquad d(\omega)=\frac{(1+\omega^2)\sin\omega\tau}{\omega},
\]
which leaves the line $p=1$ at $(1,\tau)$ as $\omega\to0$ with $p(\omega)-1=\omega^2(1-\tau^2/2)+O(\omega^4)$. The stable set is the region
between the line and the first arc of the curve; it has positive area only if the arc starts into $p>1$, i.e.\ $\tau^2<2$, and it shrinks to the
point $(1,\sqrt2)$ as $\tau\to\sqrt2$.
\end{proof}
\noindent\emph{Numerical check.} Chebyshev collocation of the delay equation \cite{breda05} ($N=32$), maximizing the decay rate over $(p,d)$: the
best decay rate is $\gamma$ at $\gamma\tau=0.5$, $0.268\gamma$ at 1.0, $0.061\gamma$ at 1.3 and $0.007\gamma$ at 1.40. Beyond the bound the
best rightmost root is unstable: $+0.008\gamma$ at 1.43 and $+0.027\gamma$ at 1.47. Interpolating, the decay rate vanishes at $\gamma\tau=1.414$.
\begin{proposition}[bandwidth]\label{prop:bode}
If the loop transfer function $L(s)$ has relative degree at least 2 and one unstable pole at $s=\gamma$, the sensitivity $S=1/(1+L)$ satisfies
Bode's integral $\int_0^\infty\ln|S(i\omega)|\,d\omega=\pi\gamma$ \cite{freudenberg85,stein03}. If the loop has negligible gain above $\omega_a$, then
$\max|S|\ge e^{\pi\gamma/\omega_a}$, and keeping $\max|S|\le2$ needs
\begin{equation}
\omega_a\ge\frac{\pi\gamma}{\ln2}=4.53\,\gamma .\label{eq:bodebw}
\end{equation}
\end{proposition}
\begin{proof}
Bode's integral is Cauchy's theorem for $\ln S$ on the right half-plane, with the unstable pole entering as a zero of $S$. With $\ln|S|\approx0$
above $\omega_a$, $\pi\gamma=\int_0^{\omega_a}\ln|S|\,d\omega\le\omega_a\ln\max|S|$.
\end{proof}
\inputs{rotor of I.17 ($M=22.4$ kg, $\Omega=3142$ rad/s, $I_r=1.171\e6$ A-turns at 30 T), $\beta=1$, $n<0$ so the vertical axis is the unstable
one; sensors with $10$ nm/$\sqrt{\rm Hz}$ (capacitive) or $1\,\mu$m/$\sqrt{\rm Hz}$ (laser) noise; control field from an opposed coil pair of
radius $R$ at $\pm0.3$ m (near-wire limit $B_r=\mu_0NI/\pi h$); a 5 mm stainless wall between coil and rotor, delay $\sim\mu_0\sigma t\ell$ with
$\ell=0.3$ m.}
\nums{\begin{center}\small
\begin{tabular}{@{}rrrrrrr@{}}
\toprule
$|n|$ & $\gamma$ (s$^{-1}$) & $\tau_c=\sqrt2/\gamma$ & $\omega_a/2\pi$ min & spin$/\omega_a$ & force at 20 $\mu$m & laser-noise rms\\
\midrule
$10^{-3}$ & 99 & 14.2 ms & 72 Hz & 7.0 & 4.4 N & 8.5 $\mu$m\\
$10^{-2}$ & 314 & 4.5 ms & 227 Hz & 2.2 & 44 N & 15 $\mu$m\\
$10^{-1}$ & 993 & 1.4 ms & 717 Hz & 0.70 & 442 N & 27 $\mu$m\\
\bottomrule
\end{tabular}\end{center}
The loop stays a factor 3 below the 500 Hz spin (where run-out appears) for $|n|\le5.4\e{-3}$. The radial control field is $6.8\e{-8}$ T per newton:
holding 20 $\mu$m at $|n|=10^{-3}$ takes 0.30 $\mu$T, resisting 0.1 $g$ takes 1.5 $\mu$T, i.e.\ about 1 ampere-turn in the opposed pair (APEX-LD's
lifting coil is 2.9 kA-turns against 60 kA-turns \cite{apexld}). The stainless wall delays the field by about 2.6 ms: inside the 14 ms budget at
$|n|=10^{-3}$, not inside 1.4 ms at $|n|=0.1$.}
\verdict{Active hold of B6's one unstable axis is easy if the stator field is uniform to $|n|\lesssim5\e{-3}$: millisecond delays, a loop below
250 Hz, microtesla control fields and micron sensors, the regime of levitated dipoles \cite{apexld,ldx07}, of Beams' magnetically suspended rotors
(one servo axis, radial centring passive, spinning in vacuum since 1937--51 \cite{beams51}), and of tokamak vertical control, where a vertically
unstable plasma is held by feedback on the radial field and the vessel's eddy currents set the delay \cite{humphreys09}. At $|n|=0.1$ the loop must
run above the spin frequency, inside the cryostat, and cancel the run-out synchronously, as electrostatically suspended gyroscopes do \cite{esg}.}

\subsection{Holding it passively: flux-conserving loops beat the zero trace (B6)}
\meta{[P] run 7 $\cdot$ gap 2 (B6) $\cdot$ from I.18 (revives the run-6 cut), I.19; in-between pairing: B6 $\times$ null-flux suspension}
\begin{theorem}[flux-conserving loops]\label{thm:fluxloops}
A rigid body carries fixed currents in the field of fixed external currents and of $N$ closed superconducting loops with inductance matrix
$\mathsf L$, cooled with no current at the equilibrium $\boldsymbol\xi=0$. Let $\mathsf G_{j\alpha}=\partial\Phi_j/\partial\xi_\alpha$ be the rate at which
the body's flux through loop $j$ changes. The translational stiffness is
\begin{equation}
\mathsf K=\mathsf K_0+\mathsf G^{\mathsf T}\mathsf L^{-1}\mathsf G,\qquad \tr\mathsf K_0=0 .\label{eq:fluxloops}
\end{equation}
The added term is positive semi-definite with rank at most $N$: one loop, or one series pair, can remove one unstable direction.
Earnshaw's bound (Theorem~\ref{thm:ring}) does not apply to induced currents.
\end{theorem}
\begin{proof}
Flux conservation, $\mathsf L\boldsymbol\iota+\boldsymbol\Phi(\boldsymbol\xi)=0$, gives the loop currents $\boldsymbol\iota=-\mathsf L^{-1}\boldsymbol\Phi$ and stores
$\tfrac12\boldsymbol\iota^{\mathsf T}\mathsf L\boldsymbol\iota=\tfrac12\boldsymbol\Phi^{\mathsf T}\mathsf L^{-1}\boldsymbol\Phi$. The body's potential is
$U=U_0(\boldsymbol\xi)+\tfrac12\boldsymbol\Phi^{\mathsf T}\mathsf L^{-1}\boldsymbol\Phi$, with $U_0$ the fixed-current part of Theorem~\ref{thm:ring}
($\tr\mathsf K_0=0$). The force $-\nabla U$ includes $\mathsf G^{\mathsf T}\boldsymbol\iota$, the push of the induced currents. With
$\boldsymbol\Phi=\mathsf G\boldsymbol\xi+O(\xi^2)$ the Hessian is \eqref{eq:fluxloops}; $\mathsf L$ is positive definite, so the added form is
$\boldsymbol\xi^{\mathsf T}\mathsf G^{\mathsf T}\mathsf L^{-1}\mathsf G\boldsymbol\xi\ge0$.
\end{proof}
\begin{corollary}[null-flux pair]\label{cor:nullflux}
Two coaxial loops of radius $a$ at heights $\pm d$ about the rotor's plane, joined in series opposition, link no flux from the rotor at any
in-plane offset or tilt when $z=0$, and couple only to vertical motion, with
\begin{equation}
g=\frac{\partial\Phi}{\partial z}=-2I_r\,\partial_zM(R,a;d),\qquad k_c=\frac{g^2}{L_{\rm pair}},\qquad L_{\rm pair}=2L_{\rm loop}-2M(a,a;2d),
\label{eq:nullflux}
\end{equation}
where $M(a,b;z)=\mu_0\sqrt{ab}\,[(2/k-k)K(k)-(2/k)E(k)]$, $k^2=4ab/[(a+b)^2+z^2]$. For $d\ll a=R$,
$k_c\approx2\mu_0aI_r^2/[d^2\ln(2d/r_w)]$. $k_c$ does not depend on the number of turns.
\end{corollary}
\begin{proof}
Reflection $z\to-z$ swaps the loops. At $z=0$ it maps a ring offset in its plane to itself, and a ring tilted about the $x$ axis to the opposite
tilt, which a half-turn about $z$ maps back. So each loop links the same flux and the pair links none, at any offset or tilt. $M$ is even in the
separation, so $\partial_z[M(z-d)-M(z+d)]_{z=0}=-2M'(d)$. With $N$ turns, $g\propto N$ and $L\propto N^2$. For $s\ll a$,
$M\approx\mu_0a[\ln(8a/s)-2]$ and $L_{\rm loop}=\mu_0a[\ln(8a/r_w)-2]$.
\end{proof}
\begin{proposition}[drift and damping]\label{prop:drift}
If the pair has resistance $R_j$ (its joints), the vertical motion and the loop current obey $M\ddot z=k_uz+g\iota$ and $L\dot\iota+R_j\iota=-g\dot z$,
with characteristic polynomial
\begin{equation}
MLs^3+MR_js^2+(g^2-k_uL)s-k_uR_j=0 .\label{eq:driftpoly}
\end{equation}
For $k_c>k_u$ one root is real, positive and small, $s\simeq(R_j/L)\,k_u/(k_c-k_u)$: the hold drifts on that time scale and needs only a slow
trim. A resistive pair adds damping, largest when its $R_n/L_n$ is near the mode frequency.
\end{proposition}
\begin{proof}
Eliminate $\iota=-gsz/(Ls+R_j)$; then $(Ms^2-k_u)(Ls+R_j)+g^2s=0$. For $R_j\to0$ the root at $s=0$ moves to first order to
$k_uR_j/(g^2-k_uL)$.
\end{proof}
\inputs{rotor of I.17 as one filament ($R=2$ m, $I_r=1.171\e6$ A-turns); loops of radius $a=R$ with effective conductor radius $r_w=5$ mm;
$L_{\rm loop}$ for surface current; joints $10^{-13}$ to $10^{-8}\,\Omega$ (superconducting joints $<10^{-12}\,\Omega$ at 4 K, soldered ones
nanohms \cite{nimsjoint}); a copper pair with half the coupling as the damper.}
\nums{Check: at $d=0.3$ m the energy route ($\partial_z^2$ of $\Phi^2/2L$ with elliptic integrals) gives $k_c=1.533\e7$ N/m, equal to
\eqref{eq:nullflux}. Stiffness of one pair: $3.1\e8$ ($d=0.08$ m), $1.9\e8$ (0.1), $3.9\e7$ (0.2), $1.5\e7$ (0.3), $4.7\e6$ (0.5) and $8.2\e5$ N/m
(1.0 m); the loop carries 16 A per 20 $\mu$m of displacement at $d=0.3$ m (0.8 kA per mm). Pairs sized for $k_c=3k_u$:
\begin{center}\small
\begin{tabular}{@{}rrrrrr@{}}
\toprule
$|n|$ & $k_u$ (N/m) & $d$ for $3k_u$ & vertical mode & in-plane mode & whirl ratio $c_n/c_r$ (I.19)\\
\midrule
$10^{-3}$ & $2.2\e5$ & 1.08 m & 140 rad/s (0.04 spin) & 70 rad/s & 44\\
$10^{-2}$ & $2.2\e6$ & 0.43 m & 444 rad/s (0.14 spin) & 222 rad/s & 13\\
$10^{-1}$ & $2.2\e7$ & 0.157 m & 1405 rad/s (0.45 spin) & 702 rad/s & 3.5\\
\bottomrule
\end{tabular}\end{center}
Drift time $1/s$ from \eqref{eq:driftpoly}: 18 years at $10^{-13}\,\Omega$, 1.8 years at $10^{-12}\,\Omega$, 16 hours at 1 n$\Omega$ and 1.6 hours at
10 n$\Omega$ ($|n|=10^{-3}$; 0.72 times these at $|n|=0.1$). A copper pair with half the superconducting pair's coupling gives a damping ratio up to
0.086, at $R_n/L_n\approx1.2$ times the mode frequency.}
\verdict{B6's unstable axis can be held passively. An axisymmetric superconducting null-flux pair supplies the missing stiffness with room to
spare; the feedback loop shrinks to a trim acting over hours to years. The pair couples only to the mean height, so rotor run-out and
eccentricity induce no first-order current and no spin drag. The idea is not new: superconducting loops were proposed to stabilize a levitated
coil passively in 1980 \cite{tran80}, null-flux coils have held maglev magnets since 1969 \cite{powell69}, and shorted superconducting turns appear
in bearing patents \cite{us5789837}. Passive hold also frees the field index: at $|n|=0.1$ the in-plane suspension is ten times stiffer and the
stator-side damping needed against whirl falls from 44 to 3.5 times the rotor-side damping, at the price of loops 0.16 m from the rotor. Run-6 cut
revived; B6 kept (D, prior art C). Deciding test: the tabletop rotor of I.17 held by a null-flux pair, with no fast loop.}

\subsection{Dynamic resistance: a ceiling on the operating current (B6 kill search)}
\meta{[E] run 7 $\cdot$ gap 2 (B6) $\cdot$ from I.20}
\begin{statement}
A tape of width $w=2a$ carrying a DC current $I_t=iI_c$ in an alternating perpendicular field of amplitude $B_a$ develops a DC resistance once
$B_a$ passes a threshold that falls linearly with $i$ \cite{jiang17,st194}:
\begin{equation}
B_{\rm th}=\kappa\,\mu_0\frac{I_c}{w}(1-i),\qquad \frac{R_{\rm dyn}}{L}=\frac{4af\,(B_a-B_{\rm th})}{I_c},\qquad P=I_t^2R_{\rm dyn},\label{eq:rdyn}
\end{equation}
with $\kappa=\tfrac12$ in the slab (parallel-field) theory and measured perpendicular thresholds two to three times higher.
\end{statement}
\inputs{LBC6 tape at 30 T: $I_c=1500$ A, $w=4$ mm, so $\mu_0I_c/w=0.47$ T; ripple bound of I.20: 3--6 mT at 500 Hz; one tape of 12.57 m.}
\nums{$B_{\rm th}=71$--212 mT at $i=0.7$, 24--71 mT at 0.9, 4.7--14 mT at 0.98, 2.4--7.1 mT at 0.99. A 3--6 mT ripple stays below threshold up to
$i=0.975$--0.996. Above it the loss is prohibitive: at $i=0.995$ with 6 mT, 360 W per tape at 4.2 K.}
\verdict{The kill search did not kill B6. Dynamic resistance is a hard ceiling near $I_c$ (within 1--3\,\%), not a constraint at the
$i\approx0.7$ that quench margin needs and production tape now allows (I.17 note); I.20's hysteresis bound still sets the ripple budget.}

\setcounter{section}{5}\setcounter{subsection}{8}
\subsection{Driving inside the cavity: what a light-held state costs at the source (gap 9)}
\meta{[P] run 7 $\cdot$ gap 9, in-between pairing (driven state $\times$ FEL optical cavity) $\cdot$ from V.3, V.4}
\begin{statement}
A sample of single-pass absorption $A$ inside an optical cavity whose other round-trip losses (mirrors, outcoupling) are $\mathcal L$, driven at
fluence $F$ per pass on area $S$ at repetition rate $f$, needs from the gain medium
\begin{equation}
P_{\rm src}=(A+\mathcal L)\,FSf\qquad\text{instead of}\qquad FSf\ \text{in a single pass},\label{eq:intracav}
\end{equation}
a gain of $G=1/(A+\mathcal L)$. The heat $AFSf$ and the condition $f\ge1/\tau$ are unchanged.
\end{statement}
\begin{proof}In steady state the gain medium replaces exactly what the round trip loses; only the absorbed part heats the sample.\end{proof}
\inputs{V.3: $F=4$ J/m$^2$ at 10 THz, 100 $\mu$m spot, $f=100$ MHz (100 K, $\tau\approx10$ ns) or 2 GHz (room temperature); $\mathcal L=0.04$
(four mirrors, no outcoupling). FELICE: 10 $\mu$J outcoupled $\approx$ 400 $\mu$J intra-cavity at 13 $\mu$m, 3--100 $\mu$m, 1 GHz micropulses in
$\sim$8 $\mu$s macropulses at 10 Hz, $\le1$ mJ \cite{felice}. Diamond: $k=2000$ W/m/K, $\rho c=1.79\e6$ J/m$^3$/K.}
\nums{Single pass: 31 nJ per pulse, 3.1 W at 100 MHz, 63 W at 2 GHz. Inside the cavity, for $A=0.02$, 0.05, 0.1, 0.2: $G=17$, 11, 7.1, 4.2,
i.e.\ 0.19--0.75 W at 100 MHz and 3.8--15 W at 2 GHz. FELICE's measured ratio is 40 (2.5\,\% outcoupling). One 400 $\mu$J pulse covers 4 J/m$^2$
over an 11 mm spot. In one macropulse a K$_3$C$_{60}$ film on diamond would be re-driven 8{,}000 times at 1 ns spacing; with $A=0.1$ the flux is
$4\e8$ W/m$^2$ and the surface rises $2q\sqrt{t/\pi k\rho c}=21$ K in 8 $\mu$s.}
\verdict{The run-2 verdict that the source binds gap 9 softens by 4--40 times when the sample sits in the cavity; the lifetime $\tau$ and the
repetition it demands still bind. New lead B7 (D; hardware exists, C): K$_3$C$_{60}$ at FELICE's intracavity focus at 30 $\mu$m. Does the
light-induced state persist through an 8 $\mu$s train of 1 GHz re-drives? At 100 K that is ten re-drives per lifetime; at room temperature,
where $\tau\ge0.5$ ns is predicted, 1 GHz is within a factor of 2 of what is needed. A condition: the sample's loss must stay below the FEL gain.}

%=====================================================================================================
\smallskip\noindent\textbf{S.4 The craft's magnet as a magnetic sail} [P+E, R7, exploration]. A dipole $m$ in a wind of density $n$ and speed $v$
holds it off at the radius where $B^2/2\mu_0=\rho v^2/2$ \cite{magsailrev}:
\begin{equation}
R_{\rm MHD}=\Big(\frac{\mu_0m}{4\pi B_{\rm mp}}\Big)^{1/3},\quad B_{\rm mp}=\sqrt{\mu_0\rho v^2},\qquad
\frac{m_pv}{eB_{\rm mp}}=\sqrt{\frac{m_p}{\mu_0ne^2}}=\frac{c}{\omega_{pi}}=d_i .\label{eq:magsail}
\end{equation}
The identity follows by substituting $B_{\rm mp}=v\sqrt{\mu_0nm_p}$: at the standoff field the proton gyroradius equals the ion inertial
length. A fluid magnetosphere needs $R_{\rm MHD}\gg d_i$; below it the deflecting region shrinks to about the St\"ormer radius
$C_{\rm St}=\sqrt{\mu_0me/4\pi m_pv}$ \cite{shaikh11,khazanov04}.
\emph{Inputs:} 1 AU, $n=5$ cm$^{-3}$, $v=400$ km/s; rotor moment $I_r\pi R^2=1.47\e7$ A\,m$^2$; an uncompensated stator (I.18's coil pair scaled to
30 T) $2.1\e9$ A\,m$^2$; craft $\ge312$ kg (rotor and minimum stator, I.17).
\emph{Numbers:} $\rho v^2=1.34$ nPa, $B_{\rm mp}=41$ nT, $d_i=102$ km. $R_{\rm MHD}=330$ m and 1.72 km, i.e.\ 0.003--0.017 $d_i$: deeply kinetic.
$C_{\rm St}=19$ and 224 m, so the thrust $\sim\rho v^2\pi C_{\rm St}^2$ is 1.5 $\mu$N to 0.21 mN (0.46--12.5 mN if the fluid picture held), 0.15--21 m/s
per year. A 100 m$^2$ solar sail gives 0.91 mN. A fluid-scale magnetosphere at 1 AU needs $m\gtrsim4\e{14}$ A\,m$^2$.
\emph{Verdict:} the craft's field is a real but tiny sail. It pushes on outside momentum, so it does not conflict with III.1--III.2. Cut as
propulsion.


%=====================================================================================================
% Pais gap loop -- mathematics run sheet, part 6: run-8 items (loaded by \input from part 2, after part 5).
% Items keep their chain numbers (counters set below). Numbers: run8_numbers.json from run8_math.py; canary8.py re-derived 15 run-7 numbers.
%=====================================================================================================
\section*{Run 8 additions to chains I, III and V}
\addcontentsline{toc}{section}{Run 8 additions: I.25--I.29, III.7, V.10--V.11}
\noindent\textbf{Notes on earlier items found in run 8.}
\begin{itemize}
\item \emph{I.17 and I.24 (operating fraction).} Run 7 took $i=I/I_c=0.70$ as safe because dynamic resistance binds only within 1--3\,\% of $I_c$
(I.24). For a rotor made of closed loops the binding limit is different: when one loop quenches its current moves into its neighbours, and
I.25 shows that $i\le0.51$ is needed for the neighbours to survive. At 30 T that needs $J_e\ge11.5$ kA/mm$^2$, 1.4 times the production tape of
I.17. This replaces an assumption; it is logged with the errors.
\item \emph{I.17, I.22 and I.23 (the deciding test).} The run-7 tabletop ring at 500 m/s has $\rho v^2/\sigma_a=2.2$: its own strength would
carry 45\,\% of any mismatch, so it would not test the stress-free regime at all. I.29 re-sizes it.
\item \emph{III.6 (X1).} III.7 adds the time scale: charge left on a dielectric keeps pulling for the dielectric's relaxation time.
\end{itemize}
Numbers below come from \texttt{run8\_math.py}; the run opened with a blind re-derivation of fifteen run-7 numbers (\texttt{canary8.py}, 15 of 15).

\setcounter{section}{1}\setcounter{subsection}{24}
\subsection{Losing the hold: quench, cascade and stored energy (B6)}
\meta{[P] run 8 $\cdot$ gap 2 (B6) $\cdot$ from I.17, I.23; kill search: flywheel containment; prior art: parallel-coil and NI-coil quench}
\begin{proposition}[mismatch tolerance]\label{prop:mismatch}
A rotor of density $\rho$ at rim speed $v$ held by $J_eBR=\rho v^2$ carries no hoop stress. If its hold falls by a fraction $\delta$ (in $J_e$,
in $B$, or locally in a strip), that part of the rim must carry the hoop stress $\delta\rho v^2$. With allowable stress $\sigma_a$ the rotor
survives only while
\begin{equation}
\delta\le\delta_{\max}=\frac{\sigma_a}{\rho v^2}=\Big(\frac{v_{\rm self}}{v}\Big)^2,\qquad v_{\rm self}=\sqrt{\sigma_a/\rho}.\label{eq:mismatch}
\end{equation}
\end{proposition}
\begin{proof}
Per unit volume the radial balance is $\rho v^2/R=J_eB+\sigma_\theta/R$ for a thin ring. With $J_eB$ reduced to $(1-\delta)\rho v^2/R$,
$\sigma_\theta=\delta\rho v^2$.
\end{proof}
\begin{theorem}[a quenched loop's current]\label{thm:redistribute}
$N$ closed superconducting loops with inductance matrix $\mathsf L$ carry $I_0$ each. If loop $q$ goes normal and its current decays to zero,
flux conservation in the others ($s\ne q$) leaves them with
\begin{equation}
\Delta\mathbf I_s=I_0\,\mathsf L_{ss}^{-1}\mathsf L_{sq},\qquad
E_{\rm diss}=\tfrac12I_0^2\big(\mathsf L_{qq}-\mathsf L_{qs}\mathsf L_{ss}^{-1}\mathsf L_{sq}\big),\label{eq:redistribute}
\end{equation}
where $E_{\rm diss}$ (the Schur complement) is the energy left to heat loop $q$. A loop survives only if $i\,(1+\Delta I_{\max}/I_0)<1$.
\end{theorem}
\begin{proof}
For $s\ne q$, $\sum_j\mathsf L_{sj}I_j$ is conserved: $\mathsf L_{ss}(\mathbf 1+\mathbf x)I_0=\mathsf L_{ss}\mathbf1I_0+\mathsf L_{sq}I_0$, so
$\mathbf x=\mathsf L_{ss}^{-1}\mathsf L_{sq}$. The stored energy falls from $\tfrac12I_0^2(\mathbf1^{\mathsf T}\mathsf L_{ss}\mathbf1+2\mathbf1^{\mathsf T}\mathsf L_{sq}
+\mathsf L_{qq})$ to $\tfrac12I_0^2(\mathbf1+\mathbf x)^{\mathsf T}\mathsf L_{ss}(\mathbf1+\mathbf x)$; the difference is \eqref{eq:redistribute}.
\end{proof}
\inputs{rotor of I.17: $R=2$ m, $\rho=8900$ kg/m$^3$, $v=6283$ m/s, $B=30$ T, section 2 mm $\times$ 100 mm as $44\times25$ tapes of
4 mm $\times$ 45 $\mu$m, $I_0=1065$ A, $I_c=1500$ A ($i=0.71$); thin-strip geometric mean distances, $L_{ij}=\mu_0R[\ln(8R/g_{ij})-2]$.
$\sigma_a=1.0$ GPa (Hastelloy-substrate tape yields at 0.9--1.1 GPa \cite{lee26scs}). Copper 10 $\mu$m per tape at $2\e{-10}$--$2\e{-9}\,\Omega$\,m
(20--30 K, with or without 30 T magnetoresistance; estimate), Hastelloy 30 $\mu$m at $1.24\e{-6}\,\Omega$\,m. Dyneema SK75 3.4 GPa, 3.5\,\%,
970 kg/m$^3$ \cite{dyneema}.}
\nums{$\rho v^2=351$ GPa, so $\delta_{\max}=0.28\,\%$ and $v_{\rm self}=335$ m/s: losing the hold means shedding 99.7\,\% of $T=442$ MJ.
\emph{One loop lost.} Coupling to a radial neighbour is $k=0.9955$, to an axial one 0.82, to the far side 0.40; $\mathsf L$ is positive definite.
An outer tape's current goes almost all into its one radial neighbour ($+98\,\%$, $i=1.40$); a middle tape's splits two ways ($+50\,\%$,
$i=1.06$). Both exceed $I_c$, so one lost tape cascades. The total current is kept to $10^{-6}$ and only 0.05--0.10 J is dissipated, so the
cascade is electrical, not thermal. Soldering the 44 radial tapes at each height into one conductor lets a local defect bypass within it
($+2.3\,\%$), but losing a whole end conductor puts $+76\,\%$ on its neighbour ($i=1.25$; middle $+50\,\%$) and drops the total by 0.24\,\%.
Survival needs $i\le0.51$ (separate tapes) or 0.57 (soldered conductors).
\emph{Everything lost.} The rotor's own inductance is $L_{\rm eq}=11.4\,\mu$H (7.9 MJ of self-field). With all tapes normal the current decays in
$\tau_q=0.05$--0.23 s, so the rim reaches $\delta_{\max}$ after $\delta_{\max}\tau_q=0.13$--0.65 ms; braking to $v_{\rm self}$ that fast takes
0.7--3.3 TW. A stator dumping in 1--30 s gives 3--85 ms and 5--150 GW. Only a slow, detected warm-up allows a safe spin-down (1 h at 123 kW).
\emph{Containment.} $T=442$ MJ is 106 kg of TNT. A catcher that stores the energy in fibre strain needs at least $T/e$ with
$e=\sigma\varepsilon_b/2\rho=61$ kJ/kg (52.7 kJ/kg elastic): 7.2--8.4 t, 23--27 times the craft's minimum mass of 312 kg. Fragments carry
19.7 MJ/kg, above what vaporizes nickel alloys (about 7.5 MJ/kg, estimate), so a bumper might beat the bound; by how much at this scale is
untested. Practice to date: a 1.9 MJ rotor ring at about 0.84 km/s inside a stainless liner with a 12.7 cm lip and a 7.6 cm lid \cite{flywheelcontain}.}
\verdict{B6 cannot be made safe at craft mass as designed: one lost loop cascades, the cascade bursts the rim within a millisecond, and the
fragments need tonnes of catcher. The mechanism of the cascade is known (parallel and no-insulation coils
\cite{cernparallel,us10984934,niquench22}); the numbers for B6 are new. B6 stays as a ground experiment (I.29). For the craft it is blocked
by stored-energy safety; revive only with $i\le0.5$ at 30 T and containment or a fail-safe below craft mass.}

\subsection{The rotor's own forces and the current uniformity it needs (B6)}
\meta{[P] run 8 $\cdot$ gap 2 (B6) $\cdot$ from I.17, I.25; in-between pairing: K2's screening-current kill search $\times$ B6}
\begin{proposition}[self-pinch]\label{prop:pinch}
A thin current sheet of height $h$ and thickness $t$ carrying $I$ (sheet current $K=I/h$) is squeezed axially; at mid-height the axial stress is
\begin{equation}
\sigma_z=\frac{\mu_0K^2h\ln2}{2\pi t}=\frac{\mu_0J_e^2A\ln2}{2\pi},\qquad A=ht,\label{eq:pinch}
\end{equation}
and its radial self-field at height $z$ is $B_r=(\mu_0K/2\pi)\ln[(h/2+z)/(h/2-z)]$.
\end{proposition}
\begin{proof}
For $h\ll R$ the sheet is locally straight. The 2D field of a strip of sheet current $K$ has the stated $B_r$ on the strip. The force on the
upper half per unit circumference is $\int_0^{h/2}KB_r\,dz=(\mu_0K^2/2\pi)\int_0^{a}\ln[(a+z)/(a-z)]\,dz$ with $a=h/2$, and the integral is
$2a\ln2$. Dividing by $t$ gives \eqref{eq:pinch}.
\end{proof}
\inputs{as I.25; $K=1.17\e7$ A/m; per tape $(I_c/w)B=11.25$ MPa and centrifugal pressure $\rho t\omega^2R=7.91$ MPa; penetration field
$\mu_0(I_c/w)/\pi=0.15$ T; delamination strength across the tape 2--100 MPa, commercial tapes 22--38 MPa at 77 K \cite{vdlaan07}; striated tapes
with 10--46 filaments exist at a 12--25\,\% $I_c$ cost \cite{striated}.}
\nums{(a) \emph{Uniformity.} By \eqref{eq:mismatch} the current density must match $\rho\omega^2r/B$ everywhere to 0.28\,\%.
(b) \emph{Faces.} The sheet's own field is $\pm\mu_0K/2=\pm7.4$ T on its faces, so the outer layers are pulled harder; the stack is in radial
compression, 22 MPa at mid-thickness ($J_e\mu_0Kt/8$): harmless.
(c) \emph{Ends.} $B_r$ is 9--11 T at the end tapes and 2.6 T a quarter of the way in, against a penetration field of 0.15 T: every tape but the
central ones carries screening currents at $\pm J_c$. In the critical state a strip $(1-i)w/2=0.58$ mm wide runs at $-J_c$ and is pushed out by
$11.25+7.91=19.2$ MPa. Held by one 45 $\mu$m tape as a cantilever that is $3qL^2/t^2=9.6$ GPa of bending (yield 1 GPa); with $n_f$ filaments it
falls as $1/n_f^2$: 96 MPa at 10, 24 MPa at 20. This is the screening-current damage seen in LBC3 and LBC6 \cite{lee26scs}, at a radius where the
hoop path cannot help.
(d) \emph{Pinch.} \eqref{eq:pinch} gives 0.95 GPa for the 2 cm$^2$ section, at the tape's yield. $J_e$ is fixed by $J_eBR=\rho v^2$, so only the
section can change: $A\le2\pi\sigma/(\mu_0J_e^2\ln2)=0.63$ cm$^2$ for 0.3 GPa (1.05 cm$^2$ for 0.5 GPa), i.e.\ 139 MJ per rotor.
(e) \emph{Through the tape.} The film carries all the current, so it must hold the rest of the tape: 5.3 MPa for the 30 $\mu$m substrate and
1.2 MPa for 7 $\mu$m of stabilizer. With the substrate on the inside that interface is in compression; the stabilizer side sits at 1.2 MPa of
tension, inside the measured range.}
\verdict{Three new B6 requirements: striated tape ($\ge10$ filaments), a section of at most about 0.6--1 cm$^2$ per rotor at 30 T, and the
substrate facing the axis. With them the self-forces fit; the 0.28\,\% current uniformity sets how finely each loop must be trimmed.}

\subsection{Heat: a spinning rotor in vacuum can only radiate (B6)}
\meta{[D] run 8 $\cdot$ gap 2 (B6) $\cdot$ from I.20, I.24; prior art: LDX, APEX-LD and Mini-RT floating coils}
\begin{equation}
P_{\rm rad}=\epsilon\sigma_{SB}A_s\,(T^4-T_w^4),\label{eq:rad}
\end{equation}
\inputs{rotor surface $A_s=2.56$ m$^2$; wall 4.2 K; emissivity 0.1--0.5; ripple loss 1 W at 6.3 mT and 500 Hz scaling as $B^3$ below
penetration (I.20); LDX floated 6--8 h while its Nb$_3$Sn coil warmed from 5 to 10 K on an on-board helium store, then was lowered and re-cooled
\cite{ldx}; APEX-LD floats 3.25 h \cite{apexld}.}
\nums{$P_{\rm rad}$ = 2--12 mW at 20 K, 12--59 mW at 30 K, 37--186 mW at 40 K, 91--454 mW at 50 K; at 4.2 K, nothing. Ripple heating is 4 mW at
1 mT, 32 mW at 2 mT, 108 mW at 3 mT. So a rotor that must stay below 30 K takes ripple $\lesssim2$ mT ($7\e{-5}$ of $B$) and tens of milliwatts
of everything else.}
\verdict{Every levitated superconductor so far floats on stored cold and is then docked; a craft rotor cannot dock, so it must balance its heat
by radiation, at 20--40 K. The 1500 A rating (I.17) is at 4.2 K. Open: the tape's $I_c$ at 30 T and 20--40 K, which sets $i$ in a radiatively
cooled rotor.}

\subsection{B6 as a flywheel and a gyroscope}
\meta{[P] run 8 $\cdot$ exploration (B6 $\times$ energy storage, B6 $\times$ craft attitude) $\cdot$ from I.2, I.17}
\begin{proposition}[moment and energy]\label{prop:mBT}
For the stress-free rotor the magnetic energy scale equals the kinetic energy: $mB=T$, with $m=I\pi R^2$.
\end{proposition}
\begin{proof}
$IBR=J_eABR=\rho v^2A=Mv^2/(2\pi R)$, so $mB=\pi R^2IB=\tfrac12Mv^2$.
\end{proof}
\inputs{as I.25; minimum stator 289 kg of T1200 (I.17); T1200 4.40 N/tex; flywheels: isotropic $\sim$12 Wh/kg, wound 22--80 Wh/kg, $\sim$550 Wh/kg
for a thin ring of the best fibre \cite{rabenhorst}; a composite system 5.6 Wh/kg \cite{tamuflywheel}.}
\nums{Rotor: 19.7 MJ/kg (5{,}480 Wh/kg), nine times a thin ring's $\sigma/2\rho$ (2.2 MJ/kg for T1200), because the stator carries the load.
With the minimum stator: 1.42 MJ/kg (394 Wh/kg), under the virial cap $\sigma/\rho=4.4$ MJ/kg of the stator's fibre. Angular momentum
$MR^2\omega=2.8\e5$ N\,m\,s; turning at 0.1 rad/s takes $2.8\e4$ N\,m, a field tilt of only $\tau/mB=\tau/T=6.4\e{-5}$ rad. A craft needs rotors
in counter-rotating pairs, as vehicle flywheels do \cite{us5921505}.}
\verdict{As a store B6 would beat every flywheel 35--70 times at system level, which is exactly why I.25's safety problem is so large.}

\subsection{The deciding tabletop test, re-sized}
\meta{[E] run 8 $\cdot$ gap 2 (B6) $\cdot$ from I.22, I.23, I.25, I.26; host: a 300 mm warm-bore magnet \cite{magnet300}}
\inputs{ring $R=0.1$ m, section 1 mm $\times$ 4 mm (22 tapes soldered as one conductor), $\rho=8900$ kg/m$^3$, host 8 T with the field falling a few
per cent across the bore ($|n|=0.03$); $J_e(8\un{T},4.2\un{K},B\parallel ab)\approx15$ kA/mm$^2$ assumed (twice the 30 T value), run at 0.7 of it;
null-flux loops of radius 0.1 m with 1 mm conductor; joints $10^{-12}\,\Omega$.}
\nums{At 500 m/s: $\rho v^2/\sigma_a=2.2$ (tolerance 45\,\%), 2.8 kJ: too slow to test anything new. At 970 m/s: $J_e=10.5$ kA/mm$^2$, 41.9 kA,
92{,}600 rpm, 22 g, $T=10.5$ kJ (2.5 g of TNT, a steel catcher), $\rho v^2/\sigma_a=8.4$, tolerance 12\,\%. Vertical growth $\gamma=1{,}680$ s$^{-1}$
asks for a 1.2 kHz loop against a 1.5 kHz spin, so hold it passively: a null-flux pair at $\pm23.5$ mm gives $\omega_z=2{,}376$ rad/s and drifts
in 22 days. Whirl: stator damping $\ge7.2\times$ rotor damping (I.19). Pinch 61 MPa.}
\verdict{A 10.5 kJ experiment that checks levitation, passive hold and point-wise force balance at 12\,\% tolerance (full scale needs
0.28\,\%). It does not test craft safety. An MRI magnet is the wrong host: $|n|\sim10^{-6}$ makes the radial support so soft that the whirl
ratio is about 1{,}400.}

%=====================================================================================================
\setcounter{section}{3}\setcounter{subsection}{6}
\subsection{X1's lingering force is a dielectric relaxation time}
\meta{[D] run 8 $\cdot$ X1 $\cdot$ from III.6; in-between pairing: X1 $\times$ trapped charge in dielectrics}
\begin{equation}
F=\frac{\sigma_t^2A}{2\varepsilon_0},\qquad \sigma_t(t)=\sigma_t(0)\,e^{-t/\tau},\qquad \tau=\varepsilon_0\varepsilon_r\rho_e .\label{eq:relax}
\end{equation}
\inputs{4.46 mN on 25 cm$^2$ (III.6); PTFE $\rho_e>10^{16}\,\Omega$\,m, $\varepsilon_r=2.1$; polyimide film $10^{13}$--$1.5\e{15}\,\Omega$\,m,
$\varepsilon_r=3.4$ \cite{dielectricsheets}.}
\nums{$\sigma_t=5.6\,\mu$C/m$^2$, i.e.\ the same $6.3\e5$ V/m as III.6. $\tau\ge2.2$ days for PTFE, 5 min to 12.5 h for polyimide.}
\verdict{A force that outlasts the power is what trapped charge does, for minutes to days. Add a control to the X1 deciding experiment:
neutralize the article (UV, an ionizer or a bake) and see whether the lingering force goes. X1 stays E.}

%=====================================================================================================
\setcounter{section}{5}\setcounter{subsection}{9}
\subsection{B7: the sample inside the cavity, and the film that settles the dispute}
\meta{[P] run 8 $\cdot$ gap 9 (B7) $\cdot$ from V.9; kill search: Dodge's critique; in-between pairing: intracavity film $\times$ depth-profile dispute}
\begin{proposition}[thin-film loss and the lasing condition]\label{prop:film}
A free-standing film of sheet conductance $\sigma_1t$ absorbs, at normal incidence,
\begin{equation}
A=\frac{4y}{(2+y)^2},\qquad y=Z_0\sigma_1t,\qquad Z_0=376.7\,\Omega,\label{eq:film}
\end{equation}
at most 0.5 (at $y=2$). An oscillator with small-signal round-trip gain $g$ and cavity loss $\mathcal L$ keeps lasing with the film inside only
if $A<g-\mathcal L$, and then the source pays $G^{-1}=A+\mathcal L$ of a single pass (V.9).
\end{proposition}
\begin{proof}
The film is a shunt admittance $\sigma_1t$ across free space: $r=-y/(2+y)$, $t_f=2/(2+y)$, and $A=1-r^2-t_f^2$. In steady state the gain
saturates to the total loss.
\end{proof}
\inputs{ELBE U27: $g\approx0.06$ against $\mathcal L\approx0.07$ at 10 $\mu$m \cite{elbe06}; FELIX: $g\approx0.5$ at first lasing, $\mathcal L\approx0.05$;
$\sigma_1(10\un{THz})$ of K$_3$C$_{60}$ $10^4$--$10^5$ S/m (estimate; only powder data found); pump depth $\Lambda\approx220$ nm \cite{dodge23}.}
\nums{ELBE-class gain leaves no room ($g-\mathcal L<0$); FELIX-class allows $A\le0.45$. For $\sigma_1=10^4$, $3\e4$, $10^5$ S/m, $A=0.1$ at 30, 10
and 3 nm and $A=0.2$ at 67, 22 and 7 nm; $G=6.7$ at $A=0.1$, 4 at 0.2.}
\verdict{Dodge \emph{et al.} argue that the K$_3$C$_{60}$ spectra depend on an assumed, non-uniform photoconductivity profile
\cite{dodge23}. A film thinner than $\Lambda/10\approx22$ nm is pumped uniformly, and its $\sigma(\omega)$ follows from $r$ and $t_f$ with no
profile model. The intracavity film is that film. B7 kept (D) and sharpened: FELIX-class gain, a 10--20 nm film, and the same sample settles
the dispute. Open: no light-driven work on K$_3$C$_{60}$ films turned up.}

\subsection{A dark cavity instead of a drive (exploration, cut)}
\meta{[D] run 8 $\cdot$ gaps 8 and 9 $\cdot$ in-between pairing: light-induced state $\times$ cavity vacuum; kill search: no-go theorem}
\begin{equation}
E_{\rm vac}=\sqrt{\frac{\hbar\omega}{2\varepsilon_0V}}\quad\text{per mode}.\label{eq:vac}
\end{equation}
\inputs{10 THz ($\lambda=30\,\mu$m); mode volume $(\lambda/2)^3$, $10^{-3}\lambda^3$, $10^{-6}\lambda^3$; drives that make the light-induced state:
$10^8$--$10^9$ V/m. Experiments: NbSe$_2$ $+10\,\%$ (3.02 to 3.41 K) and $+5.4\,\%$; NbN superfluid density $-13\,\%$; $\kappa$-ET $-50\,\%$;
Casimir cavity $\le12\,\mu$K \cite{cavitysc,nogo26}.}
\nums{$E_{\rm vac}=333$ V/m, 3.7 kV/m and 118 kV/m: $10^3$ to $10^6$ times short of the drive. Collective coupling raises the Rabi splitting as
$\sqrt N$, not the field each molecule feels.}
\verdict{A passive cavity is at most a 10\,\% knob on $T_c$, and a 2026 no-go theorem says minimal coupling lowers it. Gap 8 needs 2$\times$;
gap 9 needs energy in. Cut; revive if a passive cavity raises a $T_c$ by $\ge2\times$ or above 77 K.}


%=====================================================================================================
% Pais gap loop -- mathematics run sheet, part 7: run-9 items (loaded by \input from part 2, after part 6).
% Items keep their chain numbers (counters set below). Numbers: run9_numbers.json from run9_math.py; canary9.py re-derived 21 run-8 numbers.
%=====================================================================================================
\section*{Run 9 additions to chains I, III and V}
\addcontentsline{toc}{section}{Run 9 additions: I.30--I.32, III.8, V.12}
\noindent\textbf{Notes on earlier items found in run 9.}
\begin{itemize}
\item \emph{I.25 (the catcher).} Run 8 sized a catcher by the energy its fibres can store, $T/e_{\rm abs}$ with $e_{\rm abs}=\sigma\varepsilon/2\rho$
(7.2 t). That is the right model for fragments slower than about 4 km/s. B6's debris carries 19.7 MJ/kg, 2.6 times what it takes to vaporize
itself, so its energy ends as heat; the structure has to hold the debris' load, not absorb its energy. I.30 replaces the estimate.
\item \emph{I.17 and I.25 (thick-film tape).} The 30 T data for 4 $\mu$m films give $J_e$ on the authors' own 50 $\mu$m basis (29 kA/mm$^2$);
run 8 used a 0.1 mm basis (14.5 kA/mm$^2$), the basis of the long-length tapes. Both are kept in I.32, because which one holds decides how cold
the rotor must be.
\item \emph{I.27 (heat).} The rotor's radiative temperature now sets its operating fraction (I.32).
\item \emph{III.7 (X1).} III.8 turns the relaxation time into a control: modulate the charge with UV light and watch the force follow it.
\end{itemize}
Numbers below come from \texttt{run9\_math.py}; the run opened with a blind re-derivation of twenty-one run-8 numbers (\texttt{canary9.py}, 21 of 21
within 0.04\,\%).

\setcounter{section}{1}\setcounter{subsection}{29}
\subsection{A burst rotor's containment inherits the stator's virial (B6)}
\meta{[P] run 9 $\cdot$ gap 2 (B6) $\cdot$ from I.2, I.4, I.17, I.25; prior art: rotatable flywheel liners; no prior art found for the statement}
\begin{theorem}[containment inherits the virial]\label{thm:contain}
Let a rotor of mass $M$, radius $R$ and kinetic energy $T=\tfrac12Mv^2$ be held stress-free by a stator, so that the stator's structure carries
the rotor's reaction: a radial line load $q=Mv^2/(2\pi R^2)$ and a total hoop force $N=qR=Mv^2/(2\pi R)$, i.e.\ $\int\tr\sigma\,dV=2T$ (I.4, I.17).
After the rotor loses its hold:
\begin{enumerate}
\item[(a)] \emph{Sliding debris.} Debris pressed against a coaxial wall at radius $R+g$ with $g\ll R$ and moving at $v$ loads the wall with the
same $q$ and $N$; as friction slows it, both fall as $v^2$.
\item[(b)] \emph{Thermalized debris.} If an energy $E\le T$ ends as an ideal gas of adiabatic index $\gamma$ in a closed vessel, the vessel's
structure carries
\begin{equation}
\int\tr\sigma\,dV=3PV=3(\gamma-1)E\le3(\gamma-1)T\le2T\qquad(\gamma\le5/3),\label{eq:contain}
\end{equation}
of which a cylinder's hoop carries $2PV$ and its ends $PV$.
\end{enumerate}
So the tension a stator already supplies for the spinning rotor bounds the tension its debris can demand, in either form.
\end{theorem}
\begin{proof}
(a) A debris element $dm$ moving at $v$ along a wall of radius $R+g$ needs the inward force $dm\,v^2/(R+g)$, which only the wall can supply;
integrating round the circle gives $q$, and the hoop balance of a thin wall gives $N=qR$. Before the loss the same $q$ was the stator's reaction to
$J_eB$ on the rotor (I.17). (b) For the vessel plus gas at rest, apply I.2 with the gas as the source: the gas stress is $-P\delta_{ij}$, so the
structure carries $\int\tr\sigma=3PV$; for an ideal gas $PV=(\gamma-1)E$ (the virial theorem of a gas, $2E_{\rm kin}=3PV$ for $\gamma=5/3$
\cite{virialgas}). Energy taken up as latent heat of vaporization or ionization lowers the thermal energy, so $E\le T$. For a closed
cylinder, $\sigma_\theta=PR/t$ over the wall gives $2PV$ and $\sigma_z=PR/2t$ gives $PV$.
\end{proof}
\begin{corollary}[B6]\label{cor:contain}
B6's debris carries $v^2/2=19.7$ MJ/kg, 2.63 times the 7.5 MJ/kg that heats and vaporizes nickel-class metal (any debris faster than
$\sqrt{2e_{\rm vap}}=3.87$ km/s carries more than its own vaporization energy). A craft whose stator also serves as the containment needs, beyond
the stator, only end closures, a free-spinning liner and an ablator, plus a dynamic margin on the hoop.
\end{corollary}
\inputs{$T=441.5$ MJ, $M=22.37$ kg, $v=6283$ m/s, $R=2$ m, rotor height $h=0.10$ m (25 tapes of 4 mm); stator $\ge289$ kg of T1200 carrying $2T$,
so $\sigma/\rho=3.06\times10^6$ J/kg; wall composite T1000G-class, $\sigma=3.82$ GPa, $E=176$ GPa, $\rho=1560$ (I.3); $\gamma=1.2$--$5/3$;
liner 50--100 kg; ablation 59.8 MJ/kg (carbon) to 7.5 MJ/kg (iron); dynamic factor on the hoop 1.0--1.5 (estimate).}
\nums{Load: $q R=7.03\times10^7$ N of hoop force; on a 0.1 m wall the debris presses 351 MPa, next to the 358 MPa magnetic pressure of 30 T.
Gas: $\int\tr\sigma/2T=1.0$ ($\gamma=5/3$), 0.6 (1.4), 0.3 (1.2); end closures $(\gamma-1)T\rho/\sigma=96$, 58, 29 kg.
Impact while crossing the gap: $v_r=v\sqrt{1-R^2/(R+g)^2}=199$, 443, 882 m/s for $g=1$, 5, 20 mm; impulse $Mv_r/2\pi Rh=3.5$, 7.9,
15.7 kN\,s/m$^2$; the stator's 230 kg/m$^2$ moves at 15, 34, 69 m/s, using 0.4, 2.2, 9\,\% of the strain energy it can store
($v_{\rm allow}=\sigma/\sqrt{E\rho}=230$ m/s). Friction: the debris slows in $R/\mu v=3.2$, 1.1, 0.6 ms for $\mu=0.1$, 0.3, 0.5 at peak heat
0.28--1.4 TW; the hoop's breathing period is 1.18 ms. Liner of 50 or 100 kg: debris and liner end at 1.94 or 1.15 km/s, 69 or 82\,\% of $T$
dissipated. Ablator 7.4 kg (carbon) to 58.9 kg (iron). Containment-capable craft: 398--711 kg, 1.3--2.3 times the 312 kg minimum (run 8's catcher:
7,510 kg, 24 times).}
\noindent\emph{Caution.} Multi-shot blast vessels sit far above their static bound: LANL's M-9 vessel for 10 kg of TNT may weigh up to 75 t, while
$\tfrac32PV\rho/\sigma$ with $PV=0.3E$ and HY-100 steel is 214 kg, a factor of 350 \cite{blastvessel}. Their shock phase and elastic reuse dominate.
The estimate above assumes B6's load moves from the field to the wall smoothly, at grazing incidence, once.
\verdict{The 23$\times$ catcher of run 8 becomes about 1.3--2.3$\times$ the craft (D). The stator can be the containment because the debris needs
no more tension than the rotor did. A free-spinning liner takes the angular momentum (I.31), as in flywheel practice \cite{rotliner}. Deciding test:
a light-gas gun fires a tape-like fragment at about 6 km/s tangentially into a curved, composite-backed liner, and the hoop strain and torque are
compared with $Mv^2/2\pi R$ \cite{whipple}.}
\refs{\cite{virialgas,rotliner,whipple,blastvessel}}

\subsection{Angular momentum after a burst (B6)}
\meta{[D] run 9 $\cdot$ gap 2 (B6) $\cdot$ from I.28, I.30; prior art: NASA IPACS, rotatable liners}
\begin{statement}
The debris carries $L=2T/\omega$. Taken up by the craft's body of moment of inertia $I_c$ it spins the craft at $\omega_c=L/I_c$. Taken up by a
free-spinning liner on a fluid film, the housing receives about $L/t_{\rm sd}$ over the spin-down time $t_{\rm sd}$. A counter-rotating pair that
loses one rotor keeps a momentum bias $L$; holding attitude while slewing at $\Omega$ then needs a torque $L\Omega$.
\end{statement}
\nums{$L=2.81\times10^5$ N\,m\,s. Stator ring alone, 289 kg at 2.05 m: $I_c=1{,}215$ kg\,m$^2$, so $\omega_c=231$ rad/s (2,210 rpm, rim 474 m/s).
With a liner and UT-CEM's 120 s spin-down: 2.3 kN\,m on average (their test measured at most 63,000 ft-lb, about 85 kN\,m, against 400,000 ft-lb
predicted without the liner). Survivor bias at 0.1 rad/s: 28 kN\,m. Spinning the survivor down in one hour: 123 kW.}
\verdict{A rotatable liner is required. After one burst the craft must spin the survivor down (about an hour at 123 kW) or live with a
$2.8\times10^5$ N\,m\,s bias. NASA studied exactly this contingency for the ISS flywheel pair \cite{ipacs}.}
\refs{\cite{ipacs,rotliner}}

\subsection{The conductor at the rotor's radiative temperature (B6)}
\meta{[E] run 9 $\cdot$ gap 2 (B6) $\cdot$ from I.17, I.25, I.27}
\begin{proposition}\label{prop:radT}
A rotor in vacuum with heat input $P=P_0(\Delta B/B_0)^3$ (ripple loss, I.27), emissivity $\varepsilon$, area $A$, facing a stator at $T_s$, sits at
\begin{equation}
T_r=\Big(\frac{P}{\varepsilon\sigma_{\rm SB}A}+T_s^4\Big)^{1/4}.\label{eq:Trad}
\end{equation}
If $I_c(T)=I_c(4.2\,{\rm K})\,e^{-(T-4.2\,{\rm K})/T^*}$, the quench-safe fraction $i\le0.506$ of I.25 needs
\begin{equation}
J_e(30\,{\rm T},4.2\,{\rm K})\;\ge\;\frac{J_{\rm need}}{0.506}\,e^{(T_r-4.2\,{\rm K})/T^*}.\label{eq:Jreq}
\end{equation}
\end{proposition}
\begin{proof}Steady state: absorbed $P$ equals the net radiated $\varepsilon\sigma_{\rm SB}A(T_r^4-T_s^4)$ (gray body in a cold enclosure,
$\varepsilon_{\rm stator}\approx1$). Then substitute $I_c(T_r)$ into $i=J_{\rm need}/J_e(T_r)\le0.506$.\end{proof}
\inputs{$P_0=1$ W at $B_0=6.3$ mT (I.27); $A=2.56$ m$^2$; $T_s=4.2$ K; $T^*=21$--23 K below 40 K (commercial tape, $B\perp$, 4.2--40 K
\cite{senatore}); $J_{\rm need}=5.86$ kA/mm$^2$. Tapes at 30 T, 4.2 K, $B\parallel ab$: production 8.3 kA/mm$^2$ (LBC6); thick film, $J_c=36.3$
MA/cm$^2$ in a 4.0 $\mu$m film \cite{majkic20}: 14.5 kA/mm$^2$ on a 0.1 mm basis, 29.0 on the authors' 50 $\mu$m basis.}
\nums{$T_r$ for $\varepsilon=0.5$ to 0.1: 6.5--9.4 K at 0.3 mT ripple, 9.2--13.6 K at 0.5 mT, 15.3--22.9 K at 1 mT, 25.8--38.5 K at 2 mT. Needed
$J_e(4.2\,{\rm K})$: 12.8--14.6, 14.5--17.8, 19.2--27.1, 30.8--55 kA/mm$^2$. Production tape is never quench-safe (it needs $T<4.2$ K) and passes
$i=1$ above about 12 K (11.5--12.2 K). Thick film is safe below 9.0--9.4 K on the 0.1 mm basis (ripple $\lesssim0.3$ mT) or below 23.5--25.4 K on the 50 $\mu$m
basis (ripple $\lesssim1$ mT). Each rotor loop is 12.57 m, so 40 m pieces of A-MOCVD tape \cite{selva24,selvafnal} give three loops: 371 pieces for
1,111 loops.}
\verdict{Radiation-only cooling rules out production tape for B6 at any ripple. Thick-film tape works with ripple of 0.3--1 mT
($1$--$3\times10^{-5}$ of 30 T, MRI-class homogeneity), depending on its thickness. Grade C: the temperature scaling comes from commercial tape in
perpendicular field; $I_c(T)$ of thick film at 30 T, $B\parallel ab$, is the number to measure \cite{htsdb}.}
\par\smallskip\noindent\emph{Note on I.29 (the tabletop ring).} With $I_c\propto B^{-\alpha}$, production tape gives $J_e(8\,{\rm T})=8.3\,(30/8)^\alpha
=16.1$ ($\alpha=0.5$) to 22.4 kA/mm$^2$ ($\alpha=0.75$ above 10 T \cite{senatore}). The ring's 10.47 kA/mm$^2$ then runs at $i=0.47$--0.65 at 4.2 K;
run 8's 15 kA/mm$^2$ was conservative.
\refs{\cite{senatore,majkic20,selva24,selvafnal,htsdb}}

\setcounter{section}{3}\setcounter{subsection}{7}
\subsection{A UV charge-modulation control for X1}
\meta{[D] run 9 $\cdot$ X1 $\cdot$ from III.6, III.7; bridge: LISA Pathfinder and Gravity Probe B charge management}
\begin{statement}
A body with trapped charge $Q$ on a face of area $A$ next to a grounded electrode feels $F=a+bQ+cQ^2$, with $c=1/(2\varepsilon_0A)$ (III.7) and $b$
set by stray potentials. A propulsive force is the $a$ term, independent of $Q$. Changing $Q$ with UV photoemission (or an electron flood) and fitting
$F(Q)$ separates them; an air ionizer cannot, because it needs gas.
\end{statement}
\inputs{X1's slide: 4.46 mN on 25 cm$^2$, so $\sigma_t=5.62\,\mu$C/m$^2$ (III.7); claimed resolution $\pm20\,\mu$N; LISA Pathfinder's expected
maximum test-mass charge $2.35\times10^8\,e$; 254 nm mercury line \cite{lpf}.}
\nums{$Q=1.41\times10^{-8}$ C $=8.8\times10^{10}\,e$ (373 times LPF's maximum). $c=2.26\times10^{13}$ N/C$^2$. Discharge time $Q/I$: 14 s at 1 nA,
1.4 s at 10 nA, 0.14 s at 100 nA. Force at 30, 10 and 1\,\% of $Q$: 401, 44.6 and 0.45 $\mu$N; the $\pm20\,\mu$N resolution is reached at 6.7\,\%
of $Q$. A 254 nm photon carries 4.88 eV.}
\verdict{A UV or electron-flood control makes X1's test a short one: an electrostatic artefact follows $Q$ (as $Q^2$), a thrust stays. LISA
Pathfinder used the same method to measure charge-induced forces, including by modulating the charge on purpose \cite{lpf}.}
\refs{\cite{lpf}}

\setcounter{section}{5}\setcounter{subsection}{11}
\subsection{B7 at 30 \texorpdfstring{$\mu$m}{um}: lasing headroom, the film in monolayers, and the Meissner bound}
\meta{[D] run 9 $\cdot$ gap 9 (B7) $\cdot$ from V.9, V.10; kill search: magnetic expulsion in driven K$_3$C$_{60}$}
\begin{statement}
With single-pass gain $g$ and round-trip loss $\mathcal L$ the laser keeps lasing while the sample absorbs $A<g-\mathcal L$ (V.10). Inverting
$A=4y/(2+y)^2$ on its thin branch, $y(A)=2\big[(1-A)-\sqrt{1-2A}\big]/A$, the thickest film allowed is $t_{\max}=y(A_{\max})/(Z_0\sigma_1)$.
A K$_3$C$_{60}$(111) film of thickness $t$ holds $t/d_{111}$ monolayers, $d_{111}=a/\sqrt3$.
\end{statement}
\inputs{FELIX FEL-2 $g\approx0.5$, $\mathcal L\approx0.05$ (V.10); FHI mid-IR FEL, computed $g>2$ for 1 ps bunches at 7--30 $\mu$m,
$\mathcal L\approx0.08$ \cite{fhifel}; ELBE U100 at 30 $\mu$m, $g=0.25$--0.5 (design), $\mathcal L=0.10$ measured \cite{elbeu100}; $\sigma_1=3\times10^4$
S/m (V.10); $a=14.24$ \AA; MBE at about 2 ML/h \cite{k3c60films}.}
\nums{Headroom: FELIX 0.45 ($t_{\max}=92$ nm), FHI 1.9 (any film up to 177 nm), ELBE U100 0.15--0.40 (16--68 nm), ELBE U27 at 10 $\mu$m $-0.01$
(none). $d_{111}=0.822$ nm: 10 nm is 12 ML, 20 nm is 24 ML, 22.5 nm is 27 ML, i.e.\ 6--14 h of growth. K$_3$C$_{60}$ films on Si(111) superconduct
down to at least 25 ML (20.6 nm), and Rb$_3$C$_{60}$ films at 3--9 ML \cite{k3c60films}. Thin films of K$_3$C$_{60}$ have already been driven
optically in ultrafast transport experiments \cite{wang23}. No watt-class tabletop source exists near 30 $\mu$m \cite{dfgmir}.}
\verdict{The film B7 needs can be grown (MBE, about 20 nm), and it fits three FELs at 30 $\mu$m. A FEL stays the only source. The light-induced state
remains contested: a Faraday-rotation search found no field expulsion at about 50 nT, so $\chi_v>-0.1$ \cite{devecchi25}. The claimants expect
a state like the equilibrium superconductor near $0.8\,T_c$, which would be consistent with that bound.}
\refs{\cite{fhifel,elbeu100,k3c60films,wang23,devecchi25,dfgmir}}

% Pais gap loop -- mathematics run sheet, part 8: run-10 items (loaded by \input from part 2, after part 7).
% Items keep their chain numbers (counters set below). Numbers: run10_numbers.json from run10_math.py; canary10.py re-derived 15 run-9 numbers.
%=====================================================================================================
\section*{Run 10 additions to chains I, III and V}
\addcontentsline{toc}{section}{Run 10 additions: I.33--I.35, III.9--III.10, V.13}
\noindent\textbf{Notes on earlier items found in run 10.}
\begin{itemize}
\item \emph{I.32 (emissivity).} The best cryogenic black coatings keep $\varepsilon\ge0.95$ only down to about 30 K \cite{birb}; others fall
steeply below 50--100 K because they turn partly transparent at long wavelengths \cite{acktar}. At 5--25 K the rotor radiates at
120--580 $\mu$m, so $\varepsilon$ there is unmeasured and probably small. I.33 shows this matters little: the allowed ripple scales as $\varepsilon^{1/3}$.
\item \emph{I.30 (sliding debris).} Ricochet data exist only for flat plates: at oblique incidence the largest ricochet fragments keep 24--36\,\% of
the impact speed \cite{schonberg}. Nobody has tested debris sliding along a concave wall at $\ge3$ km/s. I.34 designs that test.
\item \emph{I.32 (thick film).} Long thick-film tapes vary in $I_c$ along their length at 4.2 K, 30 T, and thick films on thin stabilizers show
flux jumps \cite{jaroszynski25}. I.35 sizes the second effect against temperature.
\end{itemize}
Numbers below come from \texttt{run10\_math.py}; the run opened with a blind re-derivation of fifteen run-9 numbers (\texttt{canary10.py}, 15 of 15
within $10^{-4}$).

\setcounter{section}{1}\setcounter{subsection}{32}
\subsection{The field quality B6's stator needs (B6)}
\meta{[P] run 10 $\cdot$ gap 2 (B6) $\cdot$ from I.27, I.32; evidence: Iseult \cite{iseult}, a 9.4 T all-REBCO NMR magnet \cite{rebconmr}}
\begin{proposition}\label{prop:ripple}
Let the rotor see an azimuthal field ripple $\Delta B$ at harmonic $n$ of its rotation, with hysteresis loss $P=nP_0(\Delta B/B_0)^3$
($P_0=1$ W at $B_0=6.3$ mT, I.27), and shed heat only by radiation (emissivity $\varepsilon$, area $A$) to a wall at $T_s$. It stays at or below
the temperature $T_m$ its tape tolerates (I.32) if and only if
\begin{equation}
\Delta B\le B_0\left[\frac{\varepsilon\sigma A\,(T_m^4-T_s^4)}{nP_0}\right]^{1/3}.\label{eq:ripple}
\end{equation}
The bound grows as $\varepsilon^{1/3}$ and falls as $n^{-1/3}$.
\end{proposition}
\begin{proof}
In steady state the absorbed ripple loss equals the radiated power, $nP_0(\Delta B/B_0)^3=\varepsilon\sigma A(T_r^4-T_s^4)$; $T_r$ increases
with $\Delta B$, so $T_r\le T_m$ is the stated inequality.
\end{proof}
\inputs{$A=2.5635$ m$^2$, $B=30$ T, $T^*=22$ K (21--23); $T_m$ from $i\le0.506$: 9.2 K (9.0--9.4) for thick film on a 0.1 mm basis,
24.5 K (23.5--25.4) on a 50 $\mu$m basis; $\varepsilon$ 0.02--0.95; $T_s$ = 4.2, 15, 20 K.}
\nums{At $T_s=4.2$ K and $n=1$: 0.17--0.62 mT (5.7--21 ppm of 30 T) for the 0.1 mm basis and 0.64--2.31 mT (21--77 ppm) for the 50 $\mu$m basis,
across $\varepsilon=0.02$--0.95. A 48-fold range in $\varepsilon$ moves the bound only 3.6$\times$. With $T_s=20$ K (50 $\mu$m basis): 0.52--1.90 mT
(17--63 ppm). Evidence: Iseult (NbTi, 11.72 T) measured 0.88 ppm peak to peak on a 22 cm sphere with passive shims only \cite{iseult};
at 30 T that is 26 $\mu$T, a loss of $7\e{-8}$ W, and the rotor stays at 4.20--4.23 K. A 9.4 T all-REBCO NMR magnet differed from its design
by more than 100 ppm before shimming (screening currents, tape thickness) and reached sub-ppm only after passive and active shimming
\cite{rebconmr}; 100 ppm of 30 T (3 mT) would heat the rotor to 35--62 K and no tape survives. Simulated screening-current fields reach 7.8\,\%
of the centre field in high-field REBCO coils \cite{scfsim}.}
\verdict{The emissivity question folds into a field-quality condition: B6's stator must hold its azimuthal field to about 6--80 ppm along the
rotor's 12.6 m path, depending on the tape. That lies between unshimmed REBCO coils ($\ge100$ ppm) and shimmed NMR/MRI magnets ($\le1$ ppm)
(C: met in small volumes, never along a 4 m ring). Shimming and current overshoot are known remedies; drift of screening currents is the open
item.}
\refs{\cite{iseult,rebconmr,scfsim,birb,acktar}}

\subsection{A tangential fragment test with the craft's own wall pressure (B6)}
\meta{[P] run 10 $\cdot$ gap 2 (B6) $\cdot$ tests I.30; facilities \cite{swrilgg}; no prior test of sliding on a concave wall found}
\begin{proposition}[similarity and the capstan law]\label{prop:capstan}
(a) A coupon cut from the rotor's cross-section, with height $h_c\le h$, the rotor's radial thickness and any length $\ell$, moving at $v$ along
a concave track of radius $R$, presses on the track with the craft's own wall pressure $p=\lambda v^2/(Rh)$, where $\lambda=M/2\pi R$.
(b) If it slides with friction coefficient $\mu$ and loses no mass, its speed after turning through angle $\theta$ is $v_0e^{-\mu\theta}$, so
\begin{equation}
\mu=-\frac{1}{\theta}\ln\frac{v_{\rm out}}{v_{\rm in}},\qquad \Delta E/E=1-e^{-2\mu\theta}.\label{eq:capstan}
\end{equation}
\end{proposition}
\begin{proof}
(a) The coupon's mass is $\lambda(h_c/h)\ell$ and it needs the inward force $\lambda(h_c/h)\ell v^2/R$, spread over the contact area $h_c\ell$:
the pressure is $\lambda v^2/(Rh)$ for any $h_c$ and $\ell$. (b) The track supplies the normal force $N=mv^2/R$, friction $\mu N$ slows the
coupon, $m\,dv/dt=-\mu mv^2/R$; with $ds=v\,dt=R\,d\theta$ this is $dv/d\theta=-\mu v$.
\end{proof}
\inputs{$\lambda=1.78$ kg/m, $h=0.1$ m, $R=2$ m, $v=6.28$ km/s; mean rotor density 8,000 kg/m$^3$ (assumed), so radial thickness 2.2 mm;
coupon $20\times25\times2.2$ mm; sabot 20 g; SwRI's 38 mm two-stage gun: 7 km/s top, 0.5 kg at 3 km/s \cite{swrilgg}.}
\nums{Coupon 8.9 g, normal force 176 kN, wall pressure 351 MPa (the craft's), kinetic energy 176 kJ; with sabot 570 kJ at 6.28 km/s, a quarter
of the gun's 2.25 MJ at 3 km/s. Over a 30$^\circ$ arc (1.05 m, 0.17 ms) $v_{\rm out}/v_{\rm in}$ = 0.949, 0.855, 0.770 for $\mu$ = 0.1, 0.3, 0.5;
1\,\% velocimetry resolves $\mu$ to $\pm0.027$. Over a full turn the same $\mu$ leave 0.53, 0.15, 0.04 of the speed. Entering from a gap of
1, 5, 20 mm reproduces the radial impacts of I.30: 199, 443, 882 m/s. At 3 km/s the coupon carries 4.5 MJ/kg (0.6 of vaporization) and presses
80 MPa; a 0.25 m coupon then weighs 89 g (400 kJ).}
\verdict{A 9 g coupon on a 2 m-radius track reproduces the craft's wall pressure exactly, and its exit speed measures the friction that
sets I.30's slowing time. Two speeds bracket the vaporization threshold (3 and $\sim$6.3 km/s). I.30 fails if the load exceeds $\lambda v^2/R$
by much beyond the entry impact, or if debris leaves the track inward and strikes elsewhere with more than the predicted energy. Facility: C
(energy within reach; launching a flat coupon edge-on at 6 km/s is unproven).}
\refs{\cite{swrilgg,schonberg}}

\subsection{Flux-jump margin against temperature, and a warm design point (B6)}
\meta{[P] run 10 $\cdot$ gap 2 (B6) $\cdot$ criterion from \cite{wilsonmag}; observation \cite{jaroszynski25}; heat capacity model \cite{kittel}}
\inputs{Adiabatic stability of a slab of half-thickness $a$ \cite{wilsonmag}: $\beta=\mu_0J_c^2a^2/(C\,T^*)<3$, with $T^*$ replacing
$T_c-T_0$ for $J_c\propto e^{-T/T^*}$. Thick film: $J_c=36.3$ MA/cm$^2$ (30 T, $B\parallel ab$, 4.2 K), $a=2$ $\mu$m. Heat capacity of the
Hastelloy-dominated tape $C=\gamma T+aT^3$ with $\gamma=0.119$ J kg$^{-1}$K$^{-2}$ (Ni-like) and $\theta_D=400$ K (assumed, $\pm50$\,\%);
REBCO alone 630 J m$^{-3}$K$^{-1}$ at 4.2 K (assumed).}
\nums{$\mu_0J_c^2a^2=6.6\e5$ J/m$^3$. At 4.2 K: $\beta=48$ for the film alone, 0.50 and 0.25 if the heat spreads through a 50 $\mu$m or 0.1 mm tape.
Relative to 4.2 K, $\beta(T)/\beta(4.2)=e^{-2(T-4.2)/T^*}C(4.2)/C(T)$ = 0.19, 0.057, 0.020, 0.007 at 10, 15, 20, 25 K. A warm design point,
wall at 15--20 K and rotor below 24.5 K on the 50 $\mu$m basis, tolerates ripple of 0.71--0.82 mT at $\varepsilon=0.05$ and 1.5--1.8 mT at
$\varepsilon=0.5$ (I.33).}
\verdict{Whether the film flux-jumps at 4.2 K depends on how fast heat leaves the 4 $\mu$m film, which only a measurement settles; flux jumps have
been seen in thick films on thin stabilizers \cite{jaroszynski25}. Running warm cuts the margin parameter 18--50 fold at 15--20 K while the
50 $\mu$m-basis tape still carries the load, and a 20 K wall matches the cryocooler practice of the 20 T SPARC model coil \cite{tfmc}. The
conductor test should therefore ramp the field at spin-up rates at 15--25 K as well as at 4.2 K (D).}
\refs{\cite{wilsonmag,kittel,jaroszynski25,tfmc}}

\setcounter{section}{3}\setcounter{subsection}{8}
\subsection{High-frequency gravitational-wave cavities as a test of EM-to-GW claims}
\meta{[P] run 10 $\cdot$ gap 5 and side thread $\cdot$ from ledger 1 and 7; detectors \cite{hfgwcav}}
\inputs{Plane-wave flux $S=c^3\omega^2h^2/(32\pi G)$; detector reach: $h\approx10^{-21}$ at 300 MHz, $10^{-20}$ at 3 GHz, $3.9\e{-21}$ at 5.311 GHz
\cite{hfgwcav}; Pais's flux formula 1.18$\e{24}$ W/m$^2$ (ledger 1); conversion $4\pi G/(\mu_0c^4)=8.26\e{-38}$ per (T m)$^2$ \cite{gertsenshtein}.}
\nums{The detectors need $S=1.4\e{10}$, $1.4\e{14}$, $6.8\e{13}$ W/m$^2$ respectively. Pais's flux, read as gravitational waves, gives
$h=9.1\e{-15}$, $2.7\e{-15}$, $5.1\e{-16}$ at 0.3, 1, 5.3 GHz: $1.3\e5$ above the 5.3 GHz limit. Known physics for a generous generator
(10 T m, 1 MW) gives $8.3\e{-30}$ W of gravitational waves, $6.6\e{-31}$ W/m$^2$ at 1 m and $h=2\e{-42}$ at 1 GHz, $1.5\e{20}$ below reach.}
\verdict{An existing cavity detector, tuned near a claimed generator's frequency, would see Pais-scale emission at once and known-physics emission
never: a clean null test of any HFGW-generator claim. Stress-energy quadratic in the fields radiates at $2f$, so tuning to $2f$ rejects direct
electromagnetic pickup at $f$; shielding must still exceed the claimed signal ratio.}
\refs{\cite{hfgwcav,gertsenshtein}}

\subsection{Patent sweep 2023--2025: every claim falls under a ledger bound}
\meta{[P] run 10 $\cdot$ side thread $\cdot$ sources \cite{patents25,zhang25pat,oldpatents}}
\inputs{Filings found: US 12,417,870 (gravity field from a varying EM field); US 2025/0055389 (pulsed field to ``decouple the Higgs field''); US
2024/0063731 (inertia attenuator); US 2024/0375795, 2025/0286480, 2025/0296708 (closed electromagnetic drives); US 2025/0132082 (accelerated
electrons radiate); US 12,442,361 (photon drive); older US 3,095,167, 3,610,971, 3,120,363, 3,071,705, 2003/0209637.}
\nums{Each falls in one class. Closed EM systems: thrust $\le P/c=3.34$ nN/W (ledger 17). Radiation or photon drives: exactly $P/c$. Mass or
inertia reduction: field mass is positive and $\le2\e{-9}$ (ledger 16); charge leaves weight unchanged to $2\e{-3}$ (ledger 25). Gravity from
EM fields: $8.26\e{-38}$ per (T m)$^2$ (ledger 7); nothing seen at nN (Tajmar 2024 \cite{tajmar24}). Air-breathing electrostatics: ion wind
(ledger 21). The one checkable number, US 3,095,167 (charge a vehicle against Earth's fair-weather field, $\sim$130 V/m): floating 1 kg needs
0.075 C, but a 1 m sphere at air breakdown (3 MV/m) holds $3.3\e{-4}$ C, 226$\times$ short; it would float 4.4 g.}
\verdict{No new physics thread: the 2023--2025 filings repeat bounded classes, with no data. The side thread's only open edge stays the
high-frequency coupling (K4), testable with balances or a resonant gravity detector (ETH, 42 Hz, 3\,\% agreement with theory \cite{brack22}).}
\refs{\cite{patents25,zhang25pat,oldpatents,tajmar24,brack22}}

\setcounter{section}{5}\setcounter{subsection}{12}
\subsection{FELICE: intracavity fluence against the hold fluence (B7)}
\meta{[P] run 10 $\cdot$ gap 9 (B7) $\cdot$ from V.9, V.12; source \cite{felice}}
\inputs{FELICE first lasing: 10 $\mu$J outcoupled $\approx$ 400 $\mu$J intracavity per micropulse; Rayleigh range 50--82 mm at the user
stations \cite{felice}; $\lambda=30$ $\mu$m (10 THz); hold fluence 0.4 mJ/cm$^2$ \cite{rowe}.}
\nums{Intracavity factor 40. Waist $w_0=\sqrt{\lambda z_R/\pi}$ = 0.69--0.88 mm; peak fluence $2E/\pi w_0^2$ = 33--53 mJ/cm$^2$, 81--133$\times$ the
hold fluence. A film placed 0.58--0.73 m from focus sees the hold fluence over an 8 mm spot.}
\verdict{FELICE's cavity holds far more fluence than B7 needs, so a film could sit well off focus with a large spot and small loss. Its stations
are built for gas-phase samples, with limited access during operation; a cryogenic film station would be new (hardware C).}
\refs{\cite{felice,rowe}}

%=====================================================================================================
\section*{Run 11 notes and the stop rule}
\addcontentsline{toc}{section}{Run 11 notes and the stop rule}
\noindent Run 11 was the last literature pass: kill searches on run 10's conclusions, since-date checks and one bridge (ten searches; numbers
in \texttt{run11\_math.py}; canary: fifteen run-10 numbers re-derived exactly).
\begin{itemize}
\item \emph{I.33, precedent and caution.} Superconducting cyclotrons hold the first azimuthal harmonic at the few-gauss level: 2--3 G already
disturbs the beam at Kolkata; RIKEN's trim coils leave errors below 15 G; a 1 mm trim-rod error at CYCIAE-230 gives 0.24 G \cite{cyclotrons}.
That is the absolute level B6 needs at $n=1$ (0.17--2.3 mT), reached with NbTi and iron at a few tesla. The only REBCO circular-orbit magnet
built for field accuracy, the 0.7 m Ultra-Baby Skeleton Cyclotron (1.6 T design, 540 A, 30 K), quenched at 250 A while charging (46\,\% of its
design current, 36\,\% of $I_c$) and damaged its main coils; screening-current stress and cooldown contraction were blamed \cite{ubsc}. The
technique exists; its transfer to REBCO is unproven and has a failure on record.
\item \emph{I.34, gouging and attitude.} Sled rails gouge above 1.5 km/s and gun barrels above 4 km/s \cite{gouging}; B6's debris at 6.28 km/s
is 1.6--4.2 times faster, so the coupon will gouge its track. The shot must record gouge depth and shape, and the track's material pair should
be chosen from the gouging literature. Entry from a 1, 5 or 20 mm gap meets the wall at 1.8$^\circ$, 4.0$^\circ$ or 8.1$^\circ$; fragment
simulants in a two-stage gun fly with mean pitch and yaw of 4.7--5.5$^\circ$ end-on but 32--43$^\circ$ side-on \cite{fsp}. So the coupon should
fly along its length, with its attitude recorded by flash X-ray.
\item \emph{I.35, buying rules.} Edge-slit thick-film tapes fluctuate in $I_c$ and centre-slit tapes do not; metal layers suppress flux
avalanches in films \cite{jaroszynski25,avalanche}. Buy centre-slit tape with a thicker copper stabilizer.
\item \emph{K4's instrument.} A 42 Hz resonant beam with $Q\approx10^4$ and $\sim$1 pm resolution resolves $F=m\omega^2x/Q\approx$ 7--28 pN for an
effective mass of 1--4 kg (assumed), 36--140 times below a nN balance, at one frequency \cite{brack22,brack23}. A shielded electromagnetic
source whose field energy oscillates at 42 Hz (driven at 21 Hz) is a sharper K4 test than a balance.
\item \emph{X1 and B7.} No outside test and no paper for X1, and no independent driven-K$_3$C$_{60}$ experiment, as of November 2026.
\end{itemize}
\verdict{Run 11 moved no gap: no status, grade or deciding measurement changed, and the tests were refined. Every gap is settled in the loop's
sense: closed (4, 6, 10), against a wall (1, 3, 5, 7), nothing to do (11), or reduced to deciding experiments (2: stator field map,
thick-film ramp, coupon shot; 9: the FEL film test), with gap 8 on a dated re-check (February 2027). The loop stops here. It restarts if a
re-check (gap 8, E-320) finds a new record, an outside group tests Exodus or driven K$_3$C$_{60}$, a REBCO ring is field-mapped, or any
measurement from the test plans comes in.}
\refs{\cite{cyclotrons,ubsc,gouging,fsp,avalanche,brack23}}


%=====================================================================================================
\section*{K\quad Paths kept and cut in runs 3--11}
\addcontentsline{toc}{section}{K Paths kept and cut}
{\small
\begin{longtable}{@{}p{3.6cm}p{1.5cm}p{11.0cm}@{}}
\toprule
Path & Verdict & Why (item)\\
\midrule
\endhead
H1 spinning charged triangle & cut & octupole $1.6\e{-35}$ W; 1 W needs a $0.024c$ rim (II.5, I.7)\\
H2 rotating ferroelectric & cut & $\le3\e{-13}$ W at burst; depolarizing field $8.9\e8$ V/m (II.6)\\
Mass reduction by field energy & cut & $\Delta m>0$ and $\Delta m/M\le2\e{-9}$ (I.7)\\
Pais Eq.~2 flux & cut & CMUTs would emit $10^{24}\times$ what they do (II.8)\\
B3 mechanical 10 THz & cut & no acoustic mode there; resonance is optical; sasers $\lesssim1$ THz (V.5)\\
B2 drive PQP Hg-1223 & cut & light-induced effect vanishes at optimal doping (YBCO) \cite{mankowsky}; PQP sample near optimal (assumed)\\
Spinning-superconductor gravity & cut & never seen independently in another laboratory \cite{mcculloch}\\
B1 strain-matched overband & kept & high-modulus band on REBCO, test at $\varepsilon_{\rm irr}$ (I.6)\\
Gap 2 nanotube composite & kept, edge & needs $m\ge14.5$ or $\delta\lesssim0.1$ mm at $V_f\ge0.62$ (I.8--I.10)\\
B4 facility 10 THz source & kept & ERL-FEL design matches the plan; FELBE build-up test first (V.3)\\
Gap 4 rig protocol & kept & symmetry $<2\e{-6}$ and reversals; null predicted (III.3)\\
\midrule
\multicolumn{3}{@{}l}{\emph{Run 4}}\\
CNT yarn as the rim (gap 2) & cut & 0.9--1.8 N/tex at $\ge0.1$ m against 39.5 needed; below carbon fibre (I.11)\\
K1 bundle Weibull test & kept, lower prior & 12--25 bundles per gauge, two gauges a decade apart (I.12); composite needs $d\lesssim100$ nm (I.13)\\
Unpre-tensioned carbon or CNT band on REBCO & cut & separates on cooldown, or bonded pre-tensions the tape 0.18\,\% (I.14)\\
B1 metal co-wind or pre-tensioned carbon band & kept, revised & co-wind shrinking $\ge$ the tape's; warm pre-strain $\ge0.26\,\%$ (I.14)\\
Biefeld--Brown lifter & cut & ion wind, 5 mN/W in air, zero in vacuum (III.4)\\
B5 as a power route & cut & driven volume per watt $G=fF^2\approx0.1$--2 (V.6)\\
B5 as a test enabler & kept & $F\ge3.2$ in micron gaps, $3.6\times$ below damage seen (V.6)\\
K5 Tate anomaly & kept, lead & unrepeated $4.4\sigma$; repeat at $\le10$ ppm (V.7)\\
\midrule
\multicolumn{3}{@{}l}{\emph{Run 5}}\\
Continuous nanotubes round a hoop & cut & $2^{-L/74\,{\rm mm}}$: $10^{-50.8}$ survive 12.57 m (I.15)\\
K1 spliced bundle composite & kept, waiting & $\sim10^{16}$ defect-free 2 cm tubes per cm$^2$; 20/200 mm test still feasible (I.15)\\
Bonded carbon band (second reason) & cut & radial tension 27--34 MPa vs 18--53 MPa delamination (I.16)\\
K2 metal co-wind & kept, narrowed & offset measured (B); one number left: conductor shrinkage (I.16)\\
B6 stator-borne rotor & kept, new (D) & $J_eBR/\rho$: $1.3\times$ short at 30 T (R\&D tape); tilt marginal (I.17)\\
New fibre classes (M5, graphene) & cut & 3.4--3.6 and 0.73 N/tex; T1200 4.40 is the best (I.11 corr.)\\
Flyby anomaly & cut & no mechanism; did not repeat (III.5)\\
X1 charged-dielectric thrust (side research) & kept, lead (E) & image forces of the claimed size; $P/c=0$ unpowered; sealed-box test (III.6)\\
St.~Clair rotating electrostatic thruster & cut & axisymmetric spin: no radiation, no self-push (II.2, III.1, III.6)\\
\midrule
\multicolumn{3}{@{}l}{\emph{Run 6}}\\
B6 passive translation & cut & trace of stiffness 0 at fixed current; flux conservation needs offsets of 0.3--0.9 m (I.18)\\
B6 with active control & kept (D) & uniform field $|n|\lesssim10^{-3}$: 100 s$^{-1}$, newtons; LDX and APEX-LD prior art (I.18)\\
B6 whirl & condition & stator damping $>(\sqrt{2/|n|}-1)\times$ rotor damping (I.19); ripple $\lesssim10^{-4}$ of $B$ (I.20)\\
Pre-strained carbon band & cut (field) & 50 T, below the 59 T of a metal co-wind; 3.9$\times$ lighter (I.21)\\
Strain for hydrides & cut & clamping gives $\le10$ GPa, superhydrides need 100--260 (V.8)\\
Strain for nickelates & context & 2\,\% gives 3.8 GPa equivalent, about half the bulk $T_c$ (V.8)\\
Passive EDS (Inductrack) for B6 & kept (D) & needs a patterned rotor; drag about 16 W per newton at 6.3 km/s (estimate)\\
\midrule
\multicolumn{3}{@{}l}{\emph{Run 7}}\\
B6 conductor & met (C) & production tape 4 mm $\times$ 45 $\mu$m, $>1500$ A at 30 T (ab): $J_e$ 8.3 kA/mm$^2$, runs at $i=0.70$ (I.17 note)\\
B6 active hold & kept & $\gamma\tau<\sqrt2$, loop $\ge4.5\gamma$; easy for $|n|\le5\e{-3}$ (I.22)\\
B6 passive hold & revived (D) & superconducting null-flux pair adds $g^2/L$; drift hours to years; no spin drag (I.23)\\
Dynamic resistance & ceiling & binds only within 1--3\,\% of $I_c$ for 3--6 mT ripple (I.24)\\
K1 nanotube bundles as rim & cut & 35.4 N/tex at 2 cm (80 GPa on wall area), 0.90$\times$ the 39.5 needed (I.11 note)\\
B7 intracavity drive (gap 9) & kept, new (D) & source power $\times(A+\mathcal L)$, 4--40$\times$ less; FELICE test at 1 GHz (V.9)\\
Craft magnet as magnetic sail & cut & kinetic regime, 1.5 $\mu$N--0.21 mN at 1 AU (S.4)\\
X1 (Exodus) & kept, lead (E) & 4.46 mN is $5.7\e5\times$ the ejecta bound; equals a 3.5 kV pull across 5.5 mm (III.6 note)\\
\midrule
\multicolumn{3}{@{}l}{\emph{Run 8}}\\
B6 as the craft's rim & blocked (safety) & one lost loop cascades at $i=0.71$ (needs $i\le0.51$); the rim bursts within $\sim$1 ms of losing its hold; 442 MJ needs $\ge7$ t of catcher (I.25)\\
B6 rotor build & conditions & striated tape ($\ge10$ filaments), section $\le0.6$--1 cm$^2$ (pinch), substrate inward, 0.28\,\% current uniformity (I.26)\\
B6 heat & condition & radiation only: $\le$12--59 mW at 30 K; ripple $\lesssim2$ mT; $I_c$ at 20--40 K open (I.27)\\
B6 tabletop test & kept, re-sized & 10 cm ring at 970 m/s in 8 T, null-flux hold, 10.5 kJ; tolerance 12\,\% (I.29)\\
B7 intracavity film & kept, sharpened (D) & needs $A<g-\mathcal L$ (FELIX-class, not ELBE); a 10--20 nm film also settles the depth-profile dispute (V.10)\\
Dark cavity for gaps 8--9 & cut & $\le10\,\%$ on $T_c$; vacuum field $10^3$--$10^6\times$ short of the drive; no-go theorem (V.11)\\
X1 lingering force & explained (E) & trapped charge, 5.6 $\mu$C/m$^2$, relaxes in minutes to days (III.7)\\
\midrule
\multicolumn{3}{@{}l}{\emph{Run 9}}\\
B6 containment & revised (D) & debris $>3.9$ km/s ends as heat; the stator's virial bounds its load: craft 1.3--2.3$\times$, not 24$\times$ (I.30)\\
B6 after a burst & condition & free-spinning liner; survivor spun down in about an hour at 123 kW (I.31)\\
B6 conductor & condition & radiative 9--23 K: production tape never safe; thick film with ripple $\le0.3$--1 mT (I.32)\\
B6 tabletop ring & eased & $J_e(8\,{\rm T})=16$--22 kA/mm$^2$: $i=0.47$--0.65 with production tape (I.29 note)\\
B7 film & makeable (B/C) & 12--24 ML by MBE; FELIX, FHI and ELBE U100 have headroom at 30 $\mu$m (V.12)\\
K$_3$C$_{60}$ light-induced state & contested & no field expulsion at 50 nT, $\chi_v>-0.1$ (V.12)\\
Tabletop 10 THz source for B7 & cut (exploration) & no watt-class source near 30 $\mu$m (V.12)\\
X1 (Exodus) & kept, lead (E) & no outside test; UV charge-modulation control decides it (III.8)\\
\multicolumn{3}{@{}l}{\emph{Run 10}}\\
B6 emissivity & folded in & allowed ripple scales as $\varepsilon^{1/3}$; a field-quality condition replaces it (I.33)\\
B6 stator field quality & condition (C) & azimuthal field within 6--80 ppm along 12.6 m; unshimmed REBCO $\ge100$ ppm, shimmed NMR $\le1$ ppm (I.33)\\
B6 tangential fragment test & designed (C) & 9 g coupon on a 2 m track gives the craft's 351 MPa; exit speed measures $\mu$ (I.34)\\
B6 warm operation & design point (D) & wall 15--20 K, rotor $\le24.5$ K: flux-jump margin 18--50$\times$ better, ripple 0.7--1.8 mT (I.35)\\
HFGW-generator claims & testable null & cavities would see Pais-scale emission ($10^5$ margin); known physics $10^{20}$ below (III.9)\\
2023--2025 gravity/propulsion patents & no new thread & every claim falls under ledger 7, 16, 17, 21 or 25 (III.10)\\
B7 at FELICE & fluence ample (C) & 81--133$\times$ the hold fluence; no cryogenic film station (V.13)\\
\multicolumn{3}{@{}l}{\emph{Run 11 (last literature pass; the loop stops)}}\\
B6 field quality & precedent / caution & cyclotrons hold the 1st harmonic at a few gauss (NbTi, iron); the REBCO skeleton cyclotron quenched at 46\,\% of design (part 8 notes)\\
B6 coupon shot & refined & debris 1.6--4.2$\times$ above gouging thresholds: measure gouges; fly the coupon along its length (part 8 notes)\\
K4 & sharper instrument & 42 Hz resonant detector, $\sim$7--28 pN (estimate) (part 8 notes)\\
\bottomrule
\end{longtable}}

\noindent\textbf{Run-sheet changelog.} Run 3: sheet created. Re-derived from earlier runs: I.5, II.1--II.4, II.7, IV.1, IV.2, IV.4, IV.5, V.1, V.2,
V.4, S.1--S.3. New in run 3: I.1--I.4, I.6--I.10, II.5, II.6, II.8, III.1--III.3, IV.3, V.3, V.5. New in run 4 (file part 3): I.11
fibre ladder, I.12 Weibull sample size (theorem), I.13 shear lag, I.14 cooldown mismatch and pre-compression, III.4 thrust per watt, V.6
driven-volume theorem, V.7 London moment and the Tate anomaly. Run 4 lesson: sources are logged at search time (a context reset lost them).
New in run 5 (file part 3): I.15 Schulz--Flory length budget, I.16 measured cooldown offset and radial tension, I.17 stator-borne rotor
(theorem and tilt corollary), III.5 flyby anomaly; corrections to I.7 (scope), I.11 (T1200) and I.14 (tape shrinkage). From the run-5 side
research (not scored): III.6 charged-dielectric thrust claims. New in run 6 (file part 4): I.18 ring stability (theorem, rates corollary,
flux conservation), I.19 supercritical whirl, I.20 ripple loss, I.21 pre-strained band (correction to I.6 and I.14), V.8 strain as pressure;
notes on I.11 (area definition, open), I.17 (flux jumps) and III.6 (no replication). New in run 7 (file part 5): I.22 active hold (critical
delay, Bode bandwidth), I.23 passive hold by flux-conserving loops (theorem, null-flux corollary, drift), I.24 dynamic resistance, V.9
intracavity drive, S.4 magnetic sail; corrections to I.3, I.8, I.10 and I.11 (K1 is 35.4 N/tex) and I.18 (passive hold revived); notes on I.17 (production tape
meets B6) and III.6 (ejecta bound). New in run 8 (file part 6): I.25 loss of hold (mismatch tolerance, redistribution theorem for a quenched loop, cascade, global loss, containment bound), I.26 the rotor's own forces (uniformity, screening edge loads, self-pinch, through-tape stress), I.27 heat by radiation only, I.28 flywheel and gyroscope (the identity $mB=T$), I.29 the tabletop test re-sized, III.7 X1's lingering force, V.10 B7's thin film and lasing condition, V.11 a dark cavity (cut); notes on I.17 and I.24 (safe $i$ is set by quench coupling, not dynamic resistance) and on the run-7 test speed. New in run 9 (file part 7): I.30 containment inherits the stator's virial (theorem; B6 corollary), I.31 angular momentum after a burst, I.32 the conductor at the rotor's radiative temperature, III.8 a UV charge-modulation control for X1, V.12 lasing headroom at 30 $\mu$m, the film in monolayers and the Meissner bound; notes on I.25 (catcher replaced), I.17 (thick-film thickness basis), I.27 and I.29 ($J_e$ at 8 T). New in run 10 (file part 8): I.33 the field quality the stator needs, I.34 a tangential fragment test with the craft's own wall pressure, I.35 flux-jump margin against temperature and a warm design point, III.9 high-frequency GW cavities as a test of EM-to-GW claims, III.10 the 2023--2025 patent sweep, V.13 FELICE intracavity fluence; notes on I.32 (emissivity, thick-film flux jumps) and I.30 (no concave-wall data). Run 11 (file part 8, notes): cyclotron first-harmonic precedent and the REBCO skeleton-cyclotron quench (I.33), gouging and launch attitude (I.34), tape buying rules (I.35), a resonant detector for K4; the loop stopped by its rule.

%=====================================================================================================
% Pais gap loop -- mathematics run sheet, references (loaded by \input from part 2). Each run appends its new bibitems here.
\begin{thebibliography}{99}\small
\bibitem{jackson} J.\,D. Jackson, \emph{Classical Electrodynamics}, 3rd ed.\ (Wiley, 1999), \S\S6.7, 9.2, 9.7--9.8.
\bibitem{laue} M. von Laue, ``Zur Dynamik der Relativit\"atstheorie'', Ann.\ Phys.\ \textbf{340}, 524 (1911).
\bibitem{moon} F.\,C. Moon, ``The virial theorem and scaling laws for superconducting magnet systems'' (Cornell, 1981); J.\ Appl.\ Phys.\ \textbf{53}, 9112 (1982).
\bibitem{hilal} M.\,A. Hilal, O. Arici, M. Cuban, ``Minimum structure magnetic energy storage solenoids with a toroidal insert'', IEEE Trans.\ Magn.\ \textbf{21}(2), 1044 (1985).
\bibitem{tsutsui} H. Tsutsui, S. Nomura, R. Shimada, ``Optimization of SMES coil by using virial theorem'', IEEE Trans.\ Appl.\ Supercond.
\bibitem{lbl} Lawrence Berkeley Laboratory report on force-compensated coils and the virial theorem, escholarship.org/content/qt0b08n8vh.
\bibitem{cr5} IEEE Council on Superconductivity, SMES review note CR5 (2008), snf.ieeecsc.org.
\bibitem{barth} C. Barth \emph{et al.}, ``Electro-mechanical properties of REBCO coated conductors from various industrial manufacturers at 77 K, self-field and 4.2 K, 19 T'', Supercond.\ Sci.\ Technol.\ \textbf{28}, 045011 (2015); arXiv:1502.06713.
\bibitem{jcsj} Teion Kogaku (J.\ Cryo.\ Soc.\ Jpn.) \textbf{46}(5), 2011: REBCO coil reversible to 1288 MPa / 0.63 \% hoop strain.
\bibitem{sandia} T. Lambert, ``Next generation composite materials for flywheel development'', DOE ESS peer review, Sandia (2012).
\bibitem{sae} ``Design, fabrication, and testing of 10 MJ composite flywheel energy storage rotors'', SAE 981282 (1998).
\bibitem{bai} Y. Bai \emph{et al.}, ``Carbon nanotube bundles with tensile strength over 80 GPa'', Nat.\ Nanotechnol.\ \textbf{13}, 589 (2018).
\bibitem{lee} C. Lee \emph{et al.}, Science \textbf{321}, 385 (2008).
\bibitem{ma} W. Ma \emph{et al.}, ``Monitoring a micromechanical process in macroscale carbon nanotube films and fibers'', Adv.\ Mater.\ (2009).
\bibitem{beyerlein} I.\,J. Beyerlein \emph{et al.}, ``Probabilistic strength of carbon nanotube yarns'' (Springer, 2010); Nanotechnology (2011).
\bibitem{gao} E. Gao, W. Lu, Z. Xu, ``Strength loss of carbon nanotube fibers explained in a three-level hierarchical model'' (2018).
\bibitem{daniels} H.\,E. Daniels, Proc.\ R.\ Soc.\ A \textbf{183}, 405 (1945).
\bibitem{rosen} B.\,W. Rosen, AIAA J.\ \textbf{2}, 1985 (1964).
\bibitem{pourrahimi} S. and N. Pourrahimi, US 6,534,718 and US 7,476,810, ``Reinforcement of superconducting coils by high-strength materials''.
\bibitem{cmut} ``Advances in capacitive micromachined ultrasonic transducers'', Micromachines \textbf{10}, 152 (2019).
\bibitem{ergun} A.\,S. Ergun \emph{et al.}, ``Broadband capacitive micromachined ultrasonic transducers ranging from 10 kHz to 60 MHz''.
\bibitem{coleman} S. Coleman, J.\,H. Van Vleck, Phys.\ Rev.\ \textbf{171}, 1370 (1968).
\bibitem{tajmar2} M. Tajmar \emph{et al.}, ``Thrust measurements of microwave-, superconducting- and laser-type EMDrives'' (TU Dresden).
\bibitem{tajmar1} M. Tajmar \emph{et al.}, ``SpaceDrive: thrust balance development'', IAC-18-C4.7-C3.5.5 (2018).
\bibitem{soni} J. Soni, S. Roy, ``Design and characterization of a nano-Newton resolution thrust stand'', Rev.\ Sci.\ Instrum.\ \textbf{84} (2013).
\bibitem{saulson} P.\,R. Saulson, ``Thermal noise in mechanical experiments'', Phys.\ Rev.\ D \textbf{42}, 2437 (1990).
\bibitem{schwinger} J. Schwinger, Phys.\ Rev.\ \textbf{82}, 664 (1951).
\bibitem{e320} E-320 collaboration talks: ELI (Sept 2025), ``first quantitative measurements of quantum radiation reaction''; LPHYS (2026); V. Yakimenko \emph{et al.}, Phys.\ Rev.\ Accel.\ Beams \textbf{22}, 101301 (2019).
\bibitem{e320p} ``Prospects for the production and detection of Breit--Wheeler tunneling positrons in Experiment 320'', arXiv:2506.04992 (2025).
\bibitem{gertsenshtein} M.\,E. Gertsenshtein, Sov.\ Phys.\ JETP \textbf{14}, 84 (1962).
\bibitem{bcs} J. Bardeen, L.\,N. Cooper, J.\,R. Schrieffer, Phys.\ Rev.\ \textbf{108}, 1175 (1957).
\bibitem{tinkham} M. Tinkham, \emph{Introduction to Superconductivity}, 2nd ed.\ (McGraw-Hill, 1996).
\bibitem{rowe} E. Rowe \emph{et al.}, ``Resonant enhancement of photo-induced superconductivity in K$_3$C$_{60}$'', Nat.\ Phys.\ \textbf{19}, 1821 (2023); arXiv:2301.08633.
\bibitem{shanghai} ``Numerical modeling of a high power terahertz source at Shanghai'', arXiv:1107.5619 (2011).
\bibitem{saser} C.\,L. Poyser \emph{et al.}, ``Terahertz sound lasers (sasers): recent developments and applications'', Proc.\ IOA; W. Maryam, PhD thesis, Nottingham (2013).
\bibitem{mankowsky} R. Mankowsky \emph{et al.}, Struct.\ Dyn.\ \textbf{4}, 044007 (2017); C.\,R. Hunt \emph{et al.}, Phys.\ Rev.\ B \textbf{94}, 224303 (2016).
\bibitem{mcculloch} M.\,E. McCulloch, ``Can the Podkletnov effect be explained by quantised inertia?'', arXiv:1108.3488 (2011).
\bibitem{dart} A.\,F. Cheng \emph{et al.}, Nature \textbf{616}, 457 (2023).
\bibitem{dessureault} Y.-S. Dessureault \emph{et al.}, ``Tensile strength and failure modes of continuous carbon nanotube yarns for structural composite application'', SAMPE 2020; Mater.\ Sci.\ Eng.\ A (2020).
\bibitem{liang23} ``Scalable high tensile modulus composite laminates using continuous carbon nanotube yarns for aerospace applications'', ACS Appl.\ Nano Mater.\ \textbf{6}(13), 11260 (2023), doi:10.1021/acsanm.3c01266.
\bibitem{wikiss} Manufacturer strength and density data as tabulated in ``Specific strength'', Wikipedia (accessed Nov 2026).
\bibitem{kellytyson} A. Kelly, W.\,R. Tyson, J.\ Mech.\ Phys.\ Solids \textbf{13}, 329 (1965).
\bibitem{ekin} J.\,W. Ekin, \emph{Experimental Techniques for Low-Temperature Measurements} (Oxford, 2006), thermal contraction tables.
\bibitem{scis22} ``Numerical research of potential methods to reduce screening-current-induced stress for no-insulation REBCO coils'', IEEE Trans.\ Appl.\ Supercond.\ (2022).
\bibitem{yan19} Y. Yan, C. Xin, Y. Tan, T. Qu, ``Screening current effect on the stress and strain distribution in REBCO high-field magnets'', arXiv:1909.07553.
\bibitem{talley} R.\,L. Talley, ``Twenty first century propulsion concept'', Phillips Laboratory PL-TR-91-3009 (1991).
\bibitem{tajmarbb} M. Tajmar, ``Biefeld--Brown effect: misinterpretation of corona wind phenomena'', AIAA J.\ \textbf{42}, 315 (2004).
\bibitem{fan13} K. Fan \emph{et al.}, ``Nonlinear terahertz metamaterials via field-enhanced carrier dynamics in GaAs'', Phys.\ Rev.\ Lett.\ \textbf{110}, 217404 (2013); arXiv:1208.1538.
\bibitem{fan12} K. Fan \emph{et al.}, ``Nonlinear metamaterials through enhanced impact ionization in GaAs at terahertz frequencies'', APS March Meeting 2012, Q17.11.
\bibitem{tifee} ``Terahertz radiation-induced sub-cycle field electron emission across a split-gap dipole antenna'', arXiv:1508.04737 (2015).
\bibitem{electromig} Report of THz-field-induced electromigration in sub-wavelength metallic gaps (single-cycle pulses to 400 kV/cm), core.ac.uk/download/43246148.
\bibitem{tate} J. Tate, B. Cabrera, S.\,B. Felch, J.\,T. Anderson, Phys.\ Rev.\ Lett.\ \textbf{62}, 845 (1989); Phys.\ Rev.\ B \textbf{42}, 7885 (1990).
\bibitem{tajmargm} M. Tajmar \emph{et al.}, ``Experimental detection of the gravitomagnetic London moment'', arXiv:gr-qc/0603033; M. Tajmar, C.\,J. de Matos, arXiv:gr-qc/0607086 (2006).
\bibitem{tajmargm2} ``Local photon and graviton mass and its consequences'', arXiv:gr-qc/0603032 (2006).
\bibitem{e320t} E-320 positron tracker preprint, arXiv:2604.23805 (April 2026).
\bibitem{toray23} Toray Industries, ``Toray develops TORAYCA T1200, the ultra high-strength carbon fiber'', press release, 29 Oct 2023 (1,160 ksi, 1.82 g/cm$^3$).
\bibitem{zhang13} R. Zhang, Y. Zhang, Q. Zhang, H. Xie, W. Qian, F. Wei, ``Growth of half-meter long carbon nanotubes based on Schulz--Flory distribution'', ACS Nano \textbf{7}, 6156 (2013).
\bibitem{zhu19} Z. Zhu \emph{et al.}, ``Rate-selected growth of ultrapure semiconducting carbon nanotube arrays'', Nat.\ Commun.\ \textbf{10}, 4467 (2019).
\bibitem{pugno06} N.\,M. Pugno, R.\,S. Ruoff, ``Nanoscale Weibull statistics'', J.\ Appl.\ Phys.\ \textbf{99}, 024301 (2006); arXiv:cond-mat/0504518.
\bibitem{pugno08} N.\,M. Pugno, F. Bosia, A. Carpinteri, ``Multiscale stochastic simulations for tensile testing of nanotube-based macroscopic cables'', Small \textbf{4}, 1044 (2008).
\bibitem{parkcaxis} M. Park \emph{et al.}, ``De-lamination characteristics of coated conductor for conduction cooled HTS coil'' (c-axis tensile strength 18--53 MPa).
\bibitem{delam} ``Delamination strengths of different types of REBCO-coated conductors and method for reducing radial thermal stresses of impregnated REBCO pancake coils'', IEEE Trans.\ Appl.\ Supercond.
\bibitem{kar19} S. Kar \emph{et al.}, ``Next-generation highly flexible round REBCO STAR wires with over 580 A\,mm$^{-2}$ at 4.2 K, 20 T'', Supercond.\ Sci.\ Technol.\ \textbf{32}, 10LT01 (2019).
\bibitem{aries} CERN Accelerating News, Bruker BHTS tape for H2020 ARIES: $J_E$ 600 A/mm$^2$, nearly 1000 A/mm$^2$ at 4.2 K, 20 T (2019).
\bibitem{selva24} ``40-meter-long REBCO tapes'' (University of Houston), 2024: short samples $J_e\sim4600$ A/mm$^2$ at 4.2 K, 20 T.
\bibitem{cooper70} R.\,K. Cooper, C.\,H. Woods, V.\,K. Neil, C.\,E. Taylor, ``Stability analysis of a levitated superconducting current ring stabilized by feedback and eddy currents'' (1970).
\bibitem{post} R.\,F. Post (LLNL), ambient-temperature passive magnetic bearings and Earnshaw's theorem, OSTI 632846; ``Earnshaw-stability of levitated rotors'', core.ac.uk 71164777.
\bibitem{anderson08} J.\,D. Anderson, J.\,K. Campbell, J.\,E. Ekelund, J. Ellis, J.\,F. Jordan, ``Anomalous orbital-energy changes observed during spacecraft flybys of Earth'', Phys.\ Rev.\ Lett.\ \textbf{100}, 091102 (2008).
\bibitem{iorio15} L. Iorio, ``Gravitational anomalies in the solar system?'', arXiv:1412.7673.
\bibitem{acedo18} L. Acedo, P. Piqueras, J.\,A. Mora\~no, ``A possible flyby anomaly for Juno at Jupiter'', Adv.\ Space Res.\ (2018); arXiv:1711.08893.
\bibitem{buhler} Exodus Propulsion Technologies (C. Buhler), charged-dielectric thrust claims: WO 2020/159603 A2 as cited in press reports; The Debrief and IFLScience coverage (claims, no independent test).
\bibitem{hathaway21} G. Hathaway, L.\,L. Williams, ``Null-result test for effect on weight from large electrostatic charge'', Physics \textbf{3}(2), 160 (2021); arXiv:2103.00700.
\bibitem{physt18} ``Carbon nanotubes share the load'', Physics Today \textbf{71}(7), 21 (2018), on Bai \emph{et al.}: about 43 GPa on the engineering cross-section.
\bibitem{jaroszynski25} J. Jaroszynski \emph{et al.}, ``Critical current, lengthwise fluctuations, and flux jumps in REBCO CC: a torque magnetometry study up to 45 T'', arXiv:2502.02706 (2025).
\bibitem{exodus891} US 11,511,891, ``System and method for generating forces using asymmetrical electrostatic pressure'' (Exodus Propulsion Technologies), citing Canning (2004).
\bibitem{tajmar15} M. Tajmar, G. Fiedler, ``Direct thrust measurements of an EM Drive and evaluation of possible side-effects'', AIAA 2015-4083 (TU Dresden).
\bibitem{apexld} ``Design and operation of APEX-LD: a compact levitated dipole for the confinement of electron--positron pair plasmas'', arXiv:2511.07191 (2025).
\bibitem{ldx07} ``First flight of the Levitated Dipole Experiment'', APS DPP 2007 (560 kg coil, 40 min first flight, digital feedback).
\bibitem{ldxfpa} D. Garnier for the LDX team, ``Report on the Levitated Dipole eXperiment'', Fusion Power Associates (2009): over 80 hours of flight.
\bibitem{genta} G. Genta \emph{et al.} (Politecnico di Torino), rotating and non-rotating damping in rotordynamics; eddy-current dampers and electrodynamic bearings (ISMB proceedings).
\bibitem{ismb10} ISMB-10 submission 78: forward whirl unstable above $\omega_{\lim}=\sqrt{k/m}\,(1+c_{nr}/c_r)$.
\bibitem{brandt93} E.\,H. Brandt, M. Indenbom, ``Type-II-superconductor strip with current in a perpendicular magnetic field'', Phys.\ Rev.\ B \textbf{48}, 12893 (1993).
\bibitem{toraym55j} Toray TORAYCA M55J fibre and unidirectional-composite data sheets (Toray CMA); M55J/EX-1515 laminate data (Toray TAC).
\bibitem{nickfilm} Strained bilayer nickelate films: arXiv:2608.01295; ``Structural modifications in strain-engineered bilayer nickelate thin films'', Nature \textbf{653}, 76 (2026); arXiv:2510.27613.
\bibitem{sun23} H. Sun \emph{et al.}, ``Signatures of superconductivity near 80 K in a nickelate under high pressure'', Nature \textbf{621}, 493 (2023).
\bibitem{nbh} Hydrogen-loaded epitaxial Nb films: about $-10$ GPa elastic stress below 6--8 nm (Pundt group, Göttingen and KIT).
\bibitem{post96} R.\,F. Post, D.\,D. Ryutov, ``The Inductrack: a simpler approach to magnetic levitation'' (LLNL); $L/D=2\pi vL/(\lambda R)$.
% ---- run 7
\bibitem{bai18si} Y. Bai \emph{et al.}, ``Carbon nanotube bundles with tensile strength over 80 GPa'', Nat. Nanotechnol. \textbf{13}, 589 (2018), Supplementary Information, Tables S2--S3 (areas $A_1$, $A_2$ and stresses $\sigma_1$, $\sigma_2$ per bundle).
\bibitem{lbc6} LBNL ``Little Big Coil'' LBC6, EUCAS 2025 submission (2025): Shanghai Superconductor EuBCO+BaHfO$_3$ tape, 4.0 mm $\times$ 45 $\mu$m (30 $\mu$m Hastelloy), $>1500$ A at 4.2 K, 30 T, $B\parallel ab$ for production-length tape; 42.6 T at first quench.
\bibitem{majkic20} G. Majkic \emph{et al.}, ``In-field critical current performance of 4.0 $\mu$m thick film REBCO conductor with Hf addition at 4.2 K and fields up to 31.2 T'', Supercond. Sci. Technol. (2020): 5812 A per 4 mm at 30 T, $B\parallel ab$.
\bibitem{exodusslides} Exodus Propulsion Technologies, company presentation (load-cell reading 4,460 $\mu$N at 0.065 $\mu$A, 3,500 V, $1.38\e{-5}$ Torr); status: okdiario, 22 June 2026, citing Popular Mechanics (May 2026).
\bibitem{stepan89} G. St\'ep\'an, \emph{Retarded Dynamical Systems: Stability and Characteristic Functions} (Longman, 1989): the D-subdivision method.
\bibitem{breda05} D. Breda, S. Maset, R. Vermiglio, ``Pseudospectral differencing methods for characteristic roots of delay differential equations'', SIAM J. Sci. Comput. \textbf{27}, 482 (2005).
\bibitem{freudenberg85} J.\,S. Freudenberg, D.\,P. Looze, ``Right half plane poles and zeros and design tradeoffs in feedback systems'', IEEE Trans. Autom. Control \textbf{30}, 555 (1985); H.\,W. Bode, \emph{Network Analysis and Feedback Amplifier Design} (1945).
\bibitem{stein03} G. Stein, ``Respect the unstable'', IEEE Control Syst. Mag. \textbf{23}(4), 12 (2003).
\bibitem{beams51} J.\,W. Beams, J.\,D. Ross, J.\,F. Dillon, ``Magnetically suspended vacuum-type ultracentrifuge'', Rev. Sci. Instrum. \textbf{22}, 77 (1951); J.\,W. Beams, double magnetic suspension with servo-stabilized vertical axis, Rev. Sci. Instrum. \textbf{34}, 1071 (1963).
\bibitem{humphreys09} D.\,A. Humphreys \emph{et al.}, ``Experimental vertical stability studies for ITER performance and design guidance'', Nucl. Fusion \textbf{49}, 115003 (2009); DIII-D vertical position control with vessel eddy currents (Nucl. Fusion, core.ac.uk 147937471).
\bibitem{esg} D.\,A. Hill, T. Letendre, H.\,A. Mills, ``Embedded, real-time DSP control of an electrostatically suspended gyroscope'', Proc. ACC 2004; US Patent 3,955,426 (mass-unbalance modulation removed from the suspension servo).
\bibitem{tran80} M.\,Q. Tran, A. Lee, R. Bollens, ``Passive feedback stabilization of one levitated coil by means of superconducting loops'', CRPP Lausanne report LRP 172 (1980) (record seen; full text not read).
\bibitem{powell69} J.\,R. Powell, G.\,T. Danby, null-flux suspension, US Patent 3,470,828 (1969); M. Pelovitz, W.\,J. Heim, J.\,G. Murray, ``Floating ring stabilization system'', IEEE Trans. Nucl. Sci. \textbf{18}(4), 91 (1971).
\bibitem{us5789837} US Patent 5,789,837, ``High-temperature superconducting magnetic bearing'' (superconducting shorted turns; flux creep in bulk bearings); US 5,481,146 (passive null-flux coil bearing).
\bibitem{nimsjoint} Y. Takeda \emph{et al.} (NIMS), slurry-processed REBCO joints below $10^{-12}\,\Omega$ at 4 K (EUCAS 2025); REBCO single-turn loop joint below $10^{-13}\,\Omega$ (IEEE Trans. Appl. Supercond., 2019); soldered joints: Mini-RT REBCO coil, 320 h decay at 25 K (Plasma Fusion Res. 9, 1405014, 2014).
\bibitem{jiang17} Z. Jiang \emph{et al.}, ``Dynamic resistance of a high-$T_c$ coated conductor wire in a perpendicular magnetic field at 77 K'', Supercond. Sci. Technol. \textbf{30}, 03LT01 (2017); M.\,D. Ainslie \emph{et al.}, Supercond. Sci. Technol. \textbf{31}, 074003 (2018).
\bibitem{st194} Dynamic resistance of YBCO-coated conductors in applied AC fields with DC transport currents (IEEE CSC slides ST194): $R_{\rm dyn}/L=4af(B_{ac}-B_{\rm th})/I_c$; perpendicular thresholds 2--3$\times$ the slab value.
\bibitem{felice} B.\,L. Militsyn \emph{et al.}, ``FELICE --- the free electron laser for intra-cavity experiments'', Nucl. Instrum. Methods A (2003); B. Redlich, A.\,F.\,G. van der Meer, FEL 2007 (10 $\mu$J outcoupled $\approx$ 400 $\mu$J intra-cavity); arXiv:2203.05348 (2--5\,\% hole outcoupling). Stations: molecular beam and FTICR cell (gas phase), Rayleigh range 55 and 82 mm (FELIX Laboratory wiki); 4-mirror 9 m resonator, 5 cm Rayleigh range at the experiment (FEL 2007, MOAAU04).
\bibitem{magsailrev} ``Propellantless space exploration'', arXiv:2510.21743 (2025), Sec.~V; R. Zubrin, D. Andrews, J. Spacecraft Rockets \textbf{28}, 197 (1991).
\bibitem{shaikh11} I.\,F. Shaikhislamov \emph{et al.}, mini-magnetosphere in the kinetic regime, arXiv:1110.4461 (2011): size equals the MHD standoff when $d_i<R_M$ and the St\"ormer radius otherwise.
\bibitem{khazanov04} G.\,V. Khazanov \emph{et al.}, ``Plasma sail concept fundamentals'', NASA/TP-2004-213143 (2004): fluid models invalid at M2P2 scales.
\bibitem{niquench22} Current distribution and induced edge currents in no-insulation REBCO coils when an adjacent coil quenches, J. Supercond. Nov. Magn. (2022), doi:10.1007/s10948-022-06344-z (one quenched NI pancake loses almost all its current; neighbours carry induced edge currents).
\bibitem{cernparallel} Protection of parallel-connected coils: the current of a quenched coil transfers to the coils in parallel through their inductive coupling; the quench stays local only if they stay below $I_c$ (CERN indico event 763185, contribution 3416442).
\bibitem{us10984934} US Patent 10,984,934, ``Fast inductive heaters for active quench protection of superconducting coil'' (in an NI coil a quenching region's current falls and nearby turns rise toward $I_c$ by inductive coupling).
\bibitem{lee26scs} J. Lee \emph{et al.}, ``Possible self-arresting screening current stress damage in a REBCO test coil at 42.6 T in 31 T background field'', IEEE Trans. Appl. Supercond. (2026) (EUCAS 2025); J. Bang \emph{et al.}, ``Elastic-plastic conductor damage evaluation at over 0.4\,\% strain using a high-stress REBCO coil'', Supercond. Sci. Technol. \textbf{37}, 095011 (2024) (Hastelloy-substrate tape yields at 0.9--1.1 GPa).
\bibitem{dyneema} DSM Dyneema SK75 technical brochure: tenacity 35.1 cN/dtex (3.4 GPa), modulus 113 GPa, elongation 3.5\,\%; SK staple 3.3--3.9 GPa.
\bibitem{flywheelcontain} A. J. Colozza, ``High energy flywheel containment evaluation'', NASA/CR-2000-210508 (2000); A. Buchroithner \emph{et al.}, Appl. Sci. \textbf{8}, 2622 (2018) (burst containment; housing dominates system mass); US Patent 6,889,577 (energy-absorbing housing; debris energy equals rotor energy).
\bibitem{vdlaan07} D. C. van der Laan, J. W. Ekin, C. C. Clickner, T. C. Stauffer, ``Delamination strength of YBCO coated conductors under transverse tensile stress'', Supercond. Sci. Technol. \textbf{20}, 765 (2007); H.-S. Shin, M. A. Diaz, MT28 (commercial tapes 22--38 MPa at 77 K); Fermilab/LBNL slides (2--100 MPa).
\bibitem{striated} R. Gyur\'aki \emph{et al.}, ``Interfilament resistance at 77 K in striated HTS coated conductors'', IEEE Trans. Appl. Supercond. \textbf{26} (2016) (10 filaments, 20 $\mu$m grooves); J. \v{S}ouc \emph{et al.}, ``AC losses of copper stabilized multifilament YBCO coated conductors'', arXiv:1212.0452 (stripes 0.04--0.07 mm); fully stabilized tapes with 12, 24, 46 filaments, $I_c$ 410 $\to$ 360, 310 A (Ohio State).
\bibitem{ldx} ``Status of the floating coil of the Levitated Dipole Experiment'', IEEE Trans. Appl. Supercond. \textbf{12}, 666 (2002) (6--8 h levitated as the coil warms from 5 to 10 K); A. Zhukovsky, D. T. Garnier, ``Operation of the LDX floating coil'', IEEE Trans. Appl. Supercond. \textbf{16}, 898 (2006); OSTI 781377.
\bibitem{rabenhorst} D. W. Rabenhorst, Johns Hopkins APL Technical Digest \textbf{1}(2) (flywheel energy densities: isotropic $\sim$12 Wh/kg in practice, wound 22--80 Wh/kg, $\sim$550 Wh/kg theoretical for a thin ring of the best materials).
\bibitem{tamuflywheel} Texas A\&M shaftless flywheel study (composite flywheel 42 kJ/kg with shaft; 5.6 Wh/kg for the whole system).
\bibitem{us5921505} US Patent 5,921,505, ``System and method for reducing mechanical disturbances from energy storage flywheels'' (counter-rotating pairs; residual torque sensed and nulled); US 6,232,671.
\bibitem{magnet300} S. Chen \emph{et al.}, ``Development of an 8-T conduction-cooled superconducting magnet with 300-mm warm bore for material processing application'', IEEE Trans. Appl. Supercond.; Magnex 9.4 T/310 mm specification TS1330H (2005).
\bibitem{dielectricsheets} PTFE: volume resistivity $>10^{18}\,\Omega$\,cm, $\varepsilon_r=2.1$ (Professional Plastics; Proto Labs data sheet 2026); polyimide film: $10^{15}\,\Omega$\,cm (CS Hyde) to $1.5\times10^{17}\,\Omega$\,cm, $\varepsilon_r=3.4$ (ERIKS).
\bibitem{elbe06} U. Lehnert \emph{et al.}, ``Laser gain and intra-cavity losses of the ELBE mid-IR FEL'', Proc. FEL 2006, TUPPH014 (round-trip loss 5.25--10\,\%, 6\,\% gain over 7\,\% loss); FELIX first lasing, Proc. EPAC 1992, p.~46 (single-pass gain of order 50\,\%).
\bibitem{dodge23} J. S. Dodge, L. Lopez, D. G. Sahota, ``Optical saturation produces spurious evidence for photoinduced superconductivity in K$_3$C$_{60}$'', Phys. Rev. Lett. \textbf{130}, 146002 (2023); reply arXiv:2303.10169; ``More proof that the evidence for photoinduced superconductivity in K$_3$C$_{60}$ is distorted'', APS March Meeting 2025.
\bibitem{cavitysc} H. Zhang \emph{et al.}, ``Cavity enhanced superconductivity'', arXiv:2606.19171 (2026) (NbSe$_2$, 3.02 $\to$ 3.41 K); vacuum-dressed NbN in a high-$Q$ THz cavity, arXiv:2601.08191; cavity-altered superconductivity ($\kappa$-ET next to hBN), arXiv:2505.17378; A. Thomas \emph{et al.}, ``Exploring superconductivity under strong coupling with the vacuum electromagnetic field''; QEDFT for MgB$_2$, arXiv:2404.08122.
\bibitem{nogo26} ``No-go theorem and routes towards cavity-enhanced superconductivity'', arXiv:2608.14784 (2026); Casimir-cavity $T_c$ test, no shift above 12 $\mu$K.
% ---- run 9
\bibitem{virialgas} The virial theorem for a gas in a vessel, $2\langle E_{\rm kin}\rangle=3pV$: e.g.\ ``What is osmosis?'', arXiv:physics/0305011, App.~A; with $P=(\gamma-1)u$, $3(\gamma-1)E_i+\Omega=0$ (coefficient 2 for $\gamma=5/3$): ``Axion astrophysics'', arXiv:2411.02492, \S2.6.1.
\bibitem{rotliner} Rotatable flywheel containment: US 6,889,577 (liner released above a torque threshold, then riding on a fluid film); US 11,143,277 (free-spinning liner absorbs the burst rotor's angular momentum); UT-CEM full-scale composite containment for the ATTB (SAE 2004-01-0005; burst at 39,500 rpm, housing torque $\le$ 63,000 ft-lb against 400,000 predicted, spin-down 120 s; core.ac.uk 322358635).
\bibitem{whipple} E.\,L. Christiansen, modified Cour-Palais Whipple-shield equation for $V_n\ge7$ km/s (NASA JSC memo, NTRS 19920010826); ballistic, shatter and melt regimes below 3, 3--7 and above 7 km/s (Acta Astronaut., 2024, Padova); stuffed Whipple shields (zenodo 1258555).
\bibitem{blastvessel} M.\,A. Polcyn, E.\,D. Esparza, M.\,G. Whitney, ``Design of the M-9 firing facility containment vessel for Los Alamos'' (10 kg TNT, weight limit 75 t; DTIC ADA525803) and its proof test (DTIC ADA499952); LLNL, ``Composite vessels for containment of extreme blast loadings'' (2 m vessel for 35 kg TNT; OSTI 15014483); Sandia EDS vessels (11 kg TNT, 1092 shots).
\bibitem{ipacs} C.\,M. Roithmayr \emph{et al.}, ``Integrated power and attitude control for a spacecraft with flywheels and control moment gyroscopes'', NASA/TP-2003-212178; C.\,M. Roithmayr, ``International Space Station attitude motion associated with flywheel energy storage'' (abrupt stop of one rotor of a counter-rotating pair), NTRS 20040086750; US 6,113,033.
\bibitem{senatore} C. Senatore (Univ.\ Geneva), HiTAT workshop slides, CERN indico 1220254: $I_c$/width 1094--1218 A/cm at 4.2 K, 19 T and 341--460 A/cm at 20 K, 19 T ($B\perp$); $\alpha=0.75$ above 10 T; $I_{c\parallel}/I_{c\perp}=3.5$--3.8; $T^*=21$--23 K below 40 K.
\bibitem{selvafnal} V. Selvamanickam, ``Conductor and cable challenges and opportunities'', FNAL (2023): 50 m A-MOCVD tape at 3.5$\times$ the $I_c$ of typical commercial tape at 4.2 K, 20 T; STAR wires $>2.5$ kA at 4.2 K, 30 T.
\bibitem{htsdb} S.\,C. Wimbush, N.\,M. Strickland, ``A public database of high-temperature superconductor critical current data'', IEEE Trans.\ Appl.\ Supercond.\ \textbf{27}, 8000105 (2017); figshare collection 2861821 ($I_c(T,B,\theta)$, 15--90 K, up to 8 T).
\bibitem{lpf} M. Armano \emph{et al.} (LISA Pathfinder), ``Charge-induced force noise on free-falling test masses'', Phys.\ Rev.\ Lett.\ \textbf{118}, 171101 (2017); ``Precision charge control for isolated free-falling test masses'', Phys.\ Rev.\ D \textbf{98}, 062001 (2018); charge modulation by phased UV pulses, Phys.\ Rev.\ D \textbf{106}, L101101 (2022); D. Dal Bosco, PhD thesis (Trento): torsion-pendulum tests of the LISA UV-LED charge management.
\bibitem{fhifel} W. Sch\"ollkopf \emph{et al.}, ``First lasing of the IR FEL at the Fritz Haber Institute'', FEL 2012, MOOB01 (computed small-signal gain $>200$\,\% for 1 ps bunches at 7--30 $\mu$m; 1\,\% loss per mirror).
\bibitem{elbeu100} ELBE far-IR FEL (U100): 10\,\% measured loss per pass at 30 $\mu$m with a 2 mm hole, FEL 2007 MOPPH036; design gain $>25$\,\% across the range, losses $\le6$\,\% without outcoupling, FEL 2005 TUPP031.
\bibitem{k3c60films} K$_3$C$_{60}$ on Si(111) by MBE, 1--140 ML at about 2 ML/h, superconducting down to $\ge25$ ML (TU Delft, Phys.\ Rev.\ B, 2003; arXiv:cond-mat/0206221); MBE growth of A$_x$C$_{60}$ films, AAPPS Bull.\ (2021), arXiv:2202.05968; Rb$_3$C$_{60}$ films 3--9 ML, $T_c$ up to 23 K, arXiv:2102.07914.
\bibitem{wang23} E. Wang \emph{et al.}, ``Superconducting nonlinear transport in optically driven high-temperature K$_3$C$_{60}$'' (2023): thin films on photoconductive switches, granular superconducting-like response.
\bibitem{devecchi25} G. de Vecchi \emph{et al.}, ``Upper limit on magnetic field expulsion in optically driven K$_3$C$_{60}$'', Phys.\ Rev.\ Research \textbf{7}, 043270 (2025); arXiv:2502.20276 (about 50 nT resolution, $\chi_v>-0.1$).
\bibitem{dfgmir} Mid-IR difference-frequency sources: C. Gaida \emph{et al.}, Light Sci.\ Appl.\ \textbf{7}, 94 (2018) (watt-scale intrapulse DFG); 6.7 W at 2.9 $\mu$m, Opt.\ Lett.\ \textbf{45}, 1248 (2020); none at watt level near 30 $\mu$m.
% ---- run 10
\bibitem{iseult} N. Boulant, L. Quettier \emph{et al.}, ``Commissioning of the Iseult CEA 11.7 T whole-body MRI'', MAGMA \textbf{36}, 175 (2023) (NbTi, 1.8 K, 338 MJ, 132 t; spec $<0.5$ ppm peak to peak on a 22 cm sphere, 0.05 ppm/h); ISMRM 2022 abstract 0382 (measured 0.88 ppm peak to peak, RMS 0.15 ppm, passive and cryo shims).
\bibitem{birb} T. Copp, J.\,G. Tuttle \emph{et al.}, ``The area-density-dependence of Ball IR Black's low-temperature emissivity'', Cryogenics \textbf{64}, 240 (2014) ($\varepsilon\ge0.95$ from 30 to 300 K); J. Tuttle \emph{et al.}, NASA NTRS 20180000042 (most black paints fall steeply below about 100 K).
\bibitem{acktar} Acktar Fractal Black, total hemispherical emissivity 5--300 K, TFAWS 2023 (TFAWS23-CR-1): 0.8--0.9 at 40--300 K, falling significantly below about 50 K; 30 $\mu$m effective thickness.
\bibitem{rebconmr} J. Bang \emph{et al.}, ``Harmonic errors of a 9.4 T all-REBCO NMR magnet affected by screening current and geometric inconsistency of coated conductors'', Sci.\ Rep.\ (2024), PMC11333634 (over 100 ppm between design and measurement; sub-ppm after passive and active shimming, 2018).
\bibitem{scfsim} S. Shi, R. Liang, ``Numerical analysis of REBCO high-temperature superconducting coils based on screening effect'', J.\ Supercond.\ Nov.\ Magn.\ \textbf{35}(12) (2022) (screening-current field up to 0.8 T, 7.8\,\%); current overshoot and field shaking remedies, J. Korea Inst.\ Appl.\ Supercond.\ Cryog.\ (2018).
\bibitem{schonberg} W.\,P. Schonberg, R.\,A. Taylor, J.\,R. Horn, ``An analysis of penetration and ricochet phenomena in oblique hypervelocity impact'', NASA TM (MSFC, 1988); ``Oblique hypervelocity impact response of dual-sheet structures'' (NTRS 19890011875; largest ricochet fragment up to 40\,\% of the projectile diameter at 24--36\,\% of its speed); W.\,P. Schonberg, A.\,R. Ebrahim, Int.\ J.\ Impact Eng.\ \textbf{23}, 823 (1999).
\bibitem{swrilgg} Southwest Research Institute, ``Large two-stage light-gas gun facility'' (38 mm launch tube, 22 m, 7 km/s top speed, 0.5 kg at 3 km/s, 1.4 kg at 2 km/s; fragment and flyer-plate launches); JHU HyFIRE (1--7 km/s); Illinois HIT lab ($>7$ km/s).
\bibitem{wilsonmag} M.\,N. Wilson, \emph{Superconducting Magnets} (Oxford, 1983), ch.~7 (adiabatic stability of a slab, $\beta<3$).
\bibitem{kittel} C. Kittel, \emph{Introduction to Solid State Physics} (Wiley), ch.~5--6 (Debye $T^3$ and electronic $\gamma T$ heat capacity; Ni $\gamma\approx7.0$ mJ mol$^{-1}$K$^{-2}$).
\bibitem{tfmc} Z.\,S. Hartwig \emph{et al.}, ``The SPARC Toroidal Field Model Coil Program'', IEEE Trans.\ Appl.\ Supercond.\ \textbf{34}, 0600316 (2024), arXiv:2308.12301 (20.1 T at 40.5 kA, 815 kN/m, 10,058 kg, 270 km of REBCO, 600 W at 20 K, local damage after an open-circuit quench at 31.5 kA); B. LaBombard, IEEE CSC SNF 57 (2024) (110 MJ, not self-protecting).
\bibitem{hfgwcav} L. Fischer \emph{et al.}, ``First characterisation of the MAGO cavity'', Class.\ Quantum Grav.\ \textbf{42}, 115015 (2025); Y. Kim \emph{et al.}, ``Search for high-frequency gravitational waves via re-analysis of cavity axion data'', arXiv:2511.17817 ($h_0\approx3.9\e{-21}$ near 5.3 GHz); cavity reach $h\sim10^{-21}$ at 250--330 MHz and $10^{-20}$ at 2.5--3.4 GHz (Valero \emph{et al.}, 2024).
\bibitem{brack22} T. Brack \emph{et al.}, ``Measurement and theory of gravitational coupling between resonating beams'', arXiv:2201.03420 (2022) (42 Hz; detector amplitude up to $10^{-11}$ m; $2.5\e{-20}$ J/s; $d^{-4}$; 3\,\% agreement).
\bibitem{tajmar24} M. Tajmar, M. K\"o{\ss}ling, O. Neunzig, ``In-depth experimental search for a coupling between gravity and electromagnetism with steady fields'', Sci.\ Rep.\ \textbf{14} (2024), doi:10.1038/s41598-024-70286-w; arXiv:2402.15640.
\bibitem{zhang25pat} X. Zhang, Z. Xiao, X. Li, ``Method and device for driving movement of target object by gravitational field generated by varying electromagnetic field'', US Patent 12,417,870 B1 (16 Sep 2025).
\bibitem{patents25} CPC B64G1/417 filings: US 2024/0063731, US 2024/0375795, US 2025/0132082, US 2025/0286480, US 2025/0296708, US 12,442,361 (ipiq.io listing); US 2025/0055389 (``Voorhees Anti-Gravity Engine II'').
\bibitem{oldpatents} US 3,095,167 (electrostatic charge against Earth's charge); US 3,610,971 (motional electric field generator); US 3,120,363 (Electronatom lifter); US 3,071,705 (electrostatic fluid propulsion); US 2003/0209637 A1 (rotating electrostatic field; abandoned).
% ---- run 11
\bibitem{cyclotrons} J. Pradhan \emph{et al.}, ``Study of magnetic field imperfections of Kolkata superconducting cyclotron'', Cyclotrons 2010, MOPCP078 (a 2--3 G first harmonic disturbs the beam); LNS Catania first-harmonic correction (EPAC 1994, 2301); NSCL K500 (Cyclotrons 1998, G-18); J.\,W. Kim \emph{et al.}, RIKEN SRC trim coils (PAC 1997, 6P013: errors below $\pm$15 G); D. Zhang \emph{et al.}, CYCIAE-230 trim rods (Cyclotrons 2016, MOP06: 0.24 G per mm).
\bibitem{ubsc} R. Kumagai \emph{et al.}, ``Fabrication and experiments on a 1/2-scale demonstration NI-REBCO coil system of a skeleton cyclotron'', IEEE Trans.\ Appl.\ Supercond.\ \textbf{34}(5) (2024); S. Noguchi; H. Ueda \emph{et al.}, ASC 2024 (CERN indico 1431972): quench at 250 A of 540 A, coil damage, screening-current and contraction stress; SCIF and overshoot, IEEE Trans.\ Appl.\ Supercond.\ \textbf{32}, doi:10.1109/TASC.2022.3171195 (2022).
\bibitem{gouging} R.\,D.\,M. Tachau, ``An investigation of gouge initiation in high-velocity sliding contact'' (Sandia, 1991): rails above 1.5 km/s, gun barrels above 4 km/s; C.\,H. Yew, T.\,G. Trucano, R.\,D.\,M. Tachau, Int.\ J.\ Impact Eng.\ \textbf{17}, 825 (1995); J.\,D. Cinnamon, \emph{Hypervelocity Gouging Impacts} (AIAA, 2009); P.\,M. Bassett \emph{et al.}, Wear \textbf{310}, 41 (2014).
\bibitem{fsp} Fragment-simulating projectiles in a two-stage light-gas gun with separating sabots (OSTI 2004149): mean apparent pitch and yaw 4.7--5.5$^\circ$ end-on, 32--43$^\circ$ side-on.
\bibitem{avalanche} Suppression of flux avalanches by metal layers: NbN films half-coated with Cu, arXiv:1211.3787; thermomagnetic instability with an adjacent metal layer, arXiv:1304.5405; J. Jaroszynski, CCA 2025 (centre-slit tapes stable, edge-slit tapes fluctuate).
\bibitem{brack23} T. Brack \emph{et al.}, ``Dynamic measurement of gravitational coupling between resonating beams in the hertz regime'', Nat.\ Phys.\ \textbf{18}, 952 (2022); ``Dynamic gravitational excitation of structural resonances in the hertz regime using two rotating bars'', Commun.\ Phys.\ \textbf{6}, 270 (2023) (detector $Q\approx10^4$, 42 Hz, $\sim$0.1\,\% uncertainty).
\end{thebibliography}

